% SEQUÊNCIA DIDÁTICA — Volume 4 (gerado por expandir_volume.py)

%% ============================================
%% UNIDADE 0 -- Bem-vindo ao 4º Ano
%% ============================================

\chapter{Unidade 0 --- Bem-vindo ao 4º Ano}

\metadata{I-00}{Silábico-Alfabético Avançado}{EF04LP01, EF04LP02, EF04LP03, EF04LP05, EF04LP06, EF04LP07, EF04LP09}{D1}{EF04LP01, EF04LP02, EF04LP03, EF04LP05, EF04LP06, EF04LP07, EF04LP09}{Resolucao CNE/CEB 7/2010, art. 12}

\textbf{Regra:} Revisar as vogais, as consoantes simples e a leitura de sílabas diretas.

\textbf{Palavra geradora do bloco:} \textbf{MÃE} --- a palavra mais afetiva da alfabetização, revisita a nasalidade com carinho.

\section{Lição 0.1 --- Acolhida e rotina do 4º ano}

\textbf{Foco da lição:} Acolhida e rotina do 4º ano


\begin{boquinha}[Boquinha da Técnica]
Na revisão, a boca produz cada vogal com a boca aberta e o ar saindo sem obstáculo. Use o espelho para comparar A, E, I, O, U.
\end{boquinha}

\noindent\textbf{Formação guiada:}

\begin{center}
\resizebox{\textwidth}{!}{%
\begin{tabular}{ccccc}
  \sil{M}{A} & \sil{S}{A} & \sil{T}{A} & \sil{L}{E} & \sil{P}{I} \\
  \sil{B}{O} & \sil{F}{U} & \sil{N}{A} \\
  \end{tabular}}
\end{center}

\noindent\textbf{Palavras da unidade:} mãe, sapo, tatu, leite, pipa, bola, foca, navio.

\noindent\textbf{Frases para leitura oral:}
\begin{itemize}[leftmargin=1.6em]
  \item A mãe lê para mim.
  \item O sapo pula na lagoa.
  \item A pipa voa no céu.
  \item A foca nada no mar.
\end{itemize}

\section{Lição 0.2 --- Leitura do texto curto}

\noindent\textbf{Leia o texto em voz alta, com atenção às palavras da unidade:}


\begin{quote}
A mãe lê para mim. O sapo pula na lagoa.

A pipa voa no céu. A foca nada no mar.
\end{quote}

\textbf{Depois de ler, responda por escrito:}
\begin{enumerate}

  \item O que acontece no texto que você leu?
  \begin{center}
  \underline{\hspace{12cm}}
  \end{center}

  \item Quantas frases o texto tem?
  \begin{center}
  \underline{\hspace{12cm}}
  \end{center}

  \item Circule no texto as palavras com a grafia ÃE.

  \item Leia o texto de novo em voz alta, agora sem ajuda.

\end{enumerate}

\section{Lição 0.3 --- Escrita com a grafia ÃE}

\noindent\textbf{Regra da unidade:} Revisar as vogais, as consoantes simples e a leitura de sílabas diretas.


\noindent\textbf{Ditado guiado:} o professor dita as palavras abaixo; o aluno escreve e depois corrige com o professor.

\begin{center}
  \underline{\hspace{2.3cm}} \quad \underline{\hspace{2.3cm}} \quad \underline{\hspace{2.3cm}} \quad \underline{\hspace{2.3cm}}

  \underline{\hspace{2.3cm}} \quad \underline{\hspace{2.3cm}} \quad \underline{\hspace{2.3cm}} \quad \underline{\hspace{2.3cm}}
\end{center}

\noindent\textbf{Complete as palavras com a grafia ÃE:}
\begin{center}
  mã\underline{~} \quad sap\underline{~} \quad tat\underline{~} \quad leit\underline{~} \quad pip\underline{~}
\end{center}

\noindent\textbf{Escreva a palavra que o professor mostrar (banco de palavras) e separe em sílabas:}

\begin{center}
  \underline{\hspace{2cm}} $=$ \underline{\hspace{2cm}} \quad \underline{\hspace{2cm}} $=$ \underline{\hspace{2cm}} \quad \underline{\hspace{2cm}} $=$ \underline{\hspace{2cm}}
  \end{center}

%% PLANCHETA 0.1

\plancheta{Folha de Prática 0.1 --- Bem-vindo ao 4º Ano (Parte 1)}

\begin{exercicio}[Folha de Prática 0.1 --- Bem-vindo ao 4º Ano (Parte 1)]

\noindent\textbf{Nome:} \underline{\hspace{3.2cm}} \quad \textbf{Data:} \underline{\hspace{1.6cm}} \quad \textbf{Nota:} \underline{\hspace{0.9cm}}/10


\subsection*{PARTE 1: RECONHECIMENTO}

\begin{enumerate}

  \item Forme as sílabas com os sons da unidade:

\begin{center}
\resizebox{\textwidth}{!}{%
\begin{tabular}{ccccc}
  \sil{M}{A} & \sil{S}{A} & \sil{T}{A} & \sil{L}{E} & \sil{P}{I} \\
  \sil{B}{O} & \sil{F}{U} & \sil{N}{A} \\
  \end{tabular}}
\end{center}

  \item Sublinhe as palavras da unidade:

\begin{center}
  \uline{mãe} \quad \uline{sapo} \quad \uline{tatu} \quad \uline{leite} \quad \uline{pipa} \quad \uline{bola} \quad foca \quad navio \quad gato \quad \uline{sapo} \quad rato \quad \uline{bola} \quad menino
\end{center}

  \item Identifique o som-alvo em cada palavra:
  \begin{center}
  MÃE $\rightarrow$ som-alvo: \underline{\hspace{1cm}} \quad SAPO $\rightarrow$ som-alvo: \underline{\hspace{1cm}} \quad TATU $\rightarrow$ som-alvo: \underline{\hspace{1cm}}
  \end{center}

  \item Separe em sílabas:
  \begin{center}
  MÃE $\rightarrow$ \underline{\hspace{1cm}} \quad SAPO $\rightarrow$ \underline{\hspace{1cm}} \quad TATU $\rightarrow$ \underline{\hspace{1cm}}
  \end{center}

  \item Conte as sílabas:
  \begin{center}
  MÃE $\rightarrow$ \underline{\hspace{0.9cm}} \quad SAPO $\rightarrow$ \underline{\hspace{0.9cm}} \quad TATU $\rightarrow$ \underline{\hspace{0.9cm}}
  \end{center}

\end{enumerate}

\subsection*{PARTE 2: PRODUÇÃO GUIADA}

\begin{enumerate}[resume]

  \item Complete com a grafia da unidade:
  \begin{center}
  ÃE\underline{~~~} \quad ÃE\underline{~~~} \quad ÃE\underline{~~~} \quad ÃE\underline{~~~} \quad ÃE\underline{~~~}
  \end{center}

  \item Escreva 3 palavras com o som-alvo:
  \begin{center}
  1. \underline{\hspace{3cm}} \quad 2. \underline{\hspace{3cm}} \quad 3. \underline{\hspace{3cm}}
  \end{center}

  \item Copie as palavras:
  \begin{center}
  mãe \quad sapo \quad tatu \quad leite \quad pipa
  \end{center}

  \item Separe em sílabas:
  \begin{center}
  MÃE = \underline{\hspace{1cm}} \underline{\hspace{1cm}} \quad SAPO = \underline{\hspace{1cm}} \underline{\hspace{1cm}} \quad TATU = \underline{\hspace{1cm}} \underline{\hspace{1cm}}
  \end{center}

  \item Escreva 2 frases usando palavras da unidade:
  \begin{enumerate}[label=\alph*)]
    \item \underline{\hspace{8cm}}
    \item \underline{\hspace{8cm}}
  \end{enumerate}

\end{enumerate}

\end{exercicio}

%% PLANCHETA 0.2

\plancheta{Folha de Prática 0.2 --- Bem-vindo ao 4º Ano (Parte 2)}

\begin{exercicio}[Folha de Prática 0.2 --- Bem-vindo ao 4º Ano (Parte 2)]

\noindent\textbf{Nome:} \underline{\hspace{3.2cm}} \quad \textbf{Data:} \underline{\hspace{1.6cm}} \quad \textbf{Nota:} \underline{\hspace{0.9cm}}/10


\subsection*{PARTE 3: INTEGRAÇÃO}

\begin{enumerate}[resume]

  \item Leia as palavras em voz alta:
  \begin{center}
  $\rightarrow$ mãe ~ sapo ~ tatu ~ leite ~ pipa ~ bola
  \end{center}

  \item Leia as frases em voz alta:
  \begin{center}
  $\rightarrow$ A mãe lê para mim.\\
  $\rightarrow$ O sapo pula na lagoa.\\
  $\rightarrow$ A pipa voa no céu.\\
  $\rightarrow$ A foca nada no mar.\\
  \end{center}

  \item Circule a opção correta:
  \begin{enumerate}[label=(\alph*)]

    \item O som da unidade é: ( ) oral ( ) nasal 

    \item A primeira sílaba de MÃE é: ( ) MÃ ( ) MA

    \item As palavras da unidade aparecem: ( ) no início ( ) no meio ( ) no fim das palavras

  \end{enumerate}

  \item Leia e complete:
  \begin{center}
  mãe\underline{~~~} \quad sap\underline{~~~} \quad tat\underline{~~~} \quad lei\underline{~~~}
  \end{center}

  \item Leitura em voz alta com o professor:
  \begin{center}
  mãe \quad sapo \quad tatu \quad leite \quad pipa \quad bola \quad foca \quad navio
  \end{center}

\end{enumerate}

\end{exercicio}

%% PLANCHETA 0.3

\plancheta{Folha de Prática 0.3 --- Produção com ÃE}

\begin{exercicio}[Folha de Prática 0.3 --- Produção com ÃE]

\noindent\textbf{Nome:} \underline{\hspace{3.2cm}} \quad \textbf{Data:} \underline{\hspace{1.6cm}} \quad \textbf{Nota:} \underline{\hspace{0.9cm}}/15


\noindent\textbf{Comando:} Escrever um bilhete curto para a família usando palavras das vogais.


\begin{enumerate}

  \item Planeje: escolha 3 palavras da unidade para usar:
  \begin{center}
  1. \underline{\hspace{3cm}} \quad 2. \underline{\hspace{3cm}} \quad 3. \underline{\hspace{3cm}}
  \end{center}

  \item Escreva o texto (mínimo 3 frases):
  \begin{center}
  \underline{\hspace{12cm}}\\\underline{\hspace{12cm}}\\\underline{\hspace{12cm}}\\\underline{\hspace{12cm}}\\\underline{\hspace{12cm}}
  \end{center}

  \item Releia e corrija: marque palavras com a grafia ÃE que você usou.

  \item Ilustre o texto:
  \begin{center}
  \begin{tikzpicture}
    \draw[thick, dashed] (0,0) rectangle (12,6);
  \end{tikzpicture}
  \end{center}

\end{enumerate}

\end{exercicio}

%% PLANCHETA 0.4

\plancheta{Folha de Prática 0.4 --- Ditado e Avaliação da Unidade 0}

\begin{exercicio}[Folha de Prática 0.4 --- Ditado e Avaliação da Unidade 0]

\noindent\textbf{Nome:} \underline{\hspace{3.2cm}} \quad \textbf{Data:} \underline{\hspace{1.6cm}} \quad \textbf{Nota:} \underline{\hspace{0.9cm}}/20


\subsection*{PARTE 1: DITADO (10 itens)}

O professor dita as palavras e frases abaixo; o aluno escreve com letra cursiva.

\begin{enumerate}

  \item \underline{\hspace{10.5cm}}

  \item \underline{\hspace{10.5cm}}

  \item \underline{\hspace{10.5cm}}

  \item \underline{\hspace{10.5cm}}

  \item \underline{\hspace{10.5cm}}

  \item \underline{\hspace{10.5cm}}

\end{enumerate}

\subsection*{PARTE 2: AUTOAVALIAÇÃO DO ALUNO}

\begin{enumerate}[label=\alph*)]

  \item Eu li as palavras da unidade sem ajuda.  \quad ( ) Sim   ( ) Não

  \item Eu escrevi o ditado com as grafias certas.  \quad ( ) Sim   ( ) Não

  \item Eu separei as palavras em sílabas.  \quad ( ) Sim   ( ) Não

  \item Eu sei explicar a regra com as minhas palavras.  \quad ( ) Sim   ( ) Não

\end{enumerate}

\subsection*{PARTE 3: REGISTRO DO PROFESSOR (S/F/R/N)}

\begin{center}
\renewcommand{\arraystretch}{1.4}
\begin{tabular}{@{}p{9cm}cccc@{}}
\toprule
Aspecto observado & S & F & R & N \\
\midrule
Lê com precisão as palavras-alvo & & & & \\
Escreve com ortografia estável & & & & \\
Generaliza a regra em escrita espontânea & & & & \\
Mantém atenção durante o ditado & & & & \\
\bottomrule\end{tabular}\end{center}

\end{exercicio}

%% PLANCHETA 0.5

\plancheta{Folha de Prática 0.5 --- Leitura e Fluência da Unidade 0}

\begin{exercicio}[Folha de Prática 0.5 --- Leitura e Fluência da Unidade 0]

\noindent\textbf{Nome:} \underline{\hspace{3.2cm}} \quad \textbf{Data:} \underline{\hspace{1.6cm}} \quad \textbf{Nota:} \underline{\hspace{0.9cm}}/10


\subsection*{PARTE 1: LEITURA EM VOZ ALTA}

Leia o texto abaixo três vezes; a cada leitura, marque um círculo no cronômetro do professor.

\begin{quote}
A mãe lê para mim. O sapo pula na lagoa.

A pipa voa no céu. A foca nada no mar.
\end{quote}

\begin{center}
1ª leitura: \underline{\hspace{1.5cm}} \quad 2ª leitura: \underline{\hspace{1.5cm}} \quad 3ª leitura: \underline{\hspace{1.5cm}}
\end{center}

\subsection*{PARTE 2: COMPREENSÃO}

\begin{enumerate}

  \item Quem são os personagens do texto?
  \begin{center}
  \underline{\hspace{12cm}}
  \end{center}

  \item Qual é a grafia-alvo que mais aparece no texto? Escreva três exemplos:
  \begin{center}
  \underline{\hspace{3.5cm}} \quad \underline{\hspace{3.5cm}} \quad \underline{\hspace{3.5cm}}
  \end{center}

  \item Você gostou do texto? Por quê?
  \begin{center}
  \underline{\hspace{12cm}}
  \end{center}

\end{enumerate}

\subsection*{PARTE 3: DESENHO}

Desenhe uma cena do texto:

\begin{center}
\begin{tikzpicture}
  \draw[thick, dashed] (0,0) rectangle (12,5);
\end{tikzpicture}
\end{center}

\end{exercicio}

\begin{dica}
**Dica para o Professor:** Antes de avançar, garanta que a leitura de sílabas diretas esteja automática; use o espelho da Boquinha sempre que houver dúvida.
\end{dica}

%% ============================================
%% UNIDADE 1 -- O parágrafo
%% ============================================

\chapter{Unidade 1 --- O parágrafo}

\metadata{I-01}{Silábico-Alfabético Avançado}{EF04LP01, EF04LP02, EF04LP03, EF04LP05, EF04LP06, EF04LP07, EF04LP09}{D1}{EF04LP01, EF04LP02, EF04LP03, EF04LP05, EF04LP06, EF04LP07, EF04LP09}{Resolucao CNE/CEB 7/2010, art. 12}

\textbf{Regra:} CH representa o som /ch/: chave, chuva, chapéu. Em palavras como chiclete, o CH aparece no meio e no início.

\textbf{Palavra geradora do bloco:} \textbf{CHUVA} --- a chuva é cotidiana e permite sentir o som das gotas, perfeita para ancorar o dígrafo.

\section{Lição 1.1 --- Escrever parágrafos com ideia única}

\textbf{Foco da lição:} Escrever parágrafos com ideia única


\begin{boquinha}[Boquinha da Técnica]
Para dizer CH, os lábios se projetam, a língua toca o céu da boca e o ar escapa com atrito. Compare CHA com CA.
\end{boquinha}

\noindent\textbf{Formação guiada:}

\begin{center}
\resizebox{\textwidth}{!}{%
\begin{tabular}{ccccc}
  \sil{CH}{A} & \sil{CH}{E} & \sil{CH}{I} & \sil{CH}{O} & \sil{CH}{U} \\
  \end{tabular}}
\end{center}

\noindent\textbf{Palavras da unidade:} chave, chapéu, chuva, chão, cheiro, choro, fecho, chá.

\noindent\textbf{Frases para leitura oral:}
\begin{itemize}[leftmargin=1.6em]
  \item A chave abriu a porta.
  \item A chuva molhou o chão.
  \item O chapéu protege do sol.
  \item Sinto o cheiro do bolo.
\end{itemize}

\section{Lição 1.2 --- Leitura do texto curto}

\noindent\textbf{Leia o texto em voz alta, com atenção às palavras da unidade:}


\begin{quote}
A chave abriu a porta. A chuva molhou o chão.

O chapéu protege do sol. Sinto o cheiro do bolo.
\end{quote}

\textbf{Depois de ler, responda por escrito:}
\begin{enumerate}

  \item O que acontece no texto que você leu?
  \begin{center}
  \underline{\hspace{12cm}}
  \end{center}

  \item Quantas frases o texto tem?
  \begin{center}
  \underline{\hspace{12cm}}
  \end{center}

  \item Circule no texto as palavras com a grafia CH.

  \item Leia o texto de novo em voz alta, agora sem ajuda.

\end{enumerate}

\section{Lição 1.3 --- Escrita com a grafia CH}

\noindent\textbf{Regra da unidade:} CH representa o som /ch/: chave, chuva, chapéu. Em palavras como chiclete, o CH aparece no meio e no início.


\noindent\textbf{Ditado guiado:} o professor dita as palavras abaixo; o aluno escreve e depois corrige com o professor.

\begin{center}
  \underline{\hspace{2.3cm}} \quad \underline{\hspace{2.3cm}} \quad \underline{\hspace{2.3cm}} \quad \underline{\hspace{2.3cm}}

  \underline{\hspace{2.3cm}} \quad \underline{\hspace{2.3cm}} \quad \underline{\hspace{2.3cm}} \quad \underline{\hspace{2.3cm}}
\end{center}

\noindent\textbf{Complete as palavras com a grafia CH:}
\begin{center}
  \underline{~~~}ave \quad \underline{~~~}apéu \quad \underline{~~~}uva \quad \underline{~~~}ão \quad \underline{~~~}eiro
\end{center}

\noindent\textbf{Escreva a palavra que o professor mostrar (banco de palavras) e separe em sílabas:}

\begin{center}
  \underline{\hspace{2cm}} $=$ \underline{\hspace{2cm}} \quad \underline{\hspace{2cm}} $=$ \underline{\hspace{2cm}} \quad \underline{\hspace{2cm}} $=$ \underline{\hspace{2cm}}
  \end{center}

%% PLANCHETA 1.1

\plancheta{Folha de Prática 1.1 --- O parágrafo (Parte 1)}

\begin{exercicio}[Folha de Prática 1.1 --- O parágrafo (Parte 1)]

\noindent\textbf{Nome:} \underline{\hspace{3.2cm}} \quad \textbf{Data:} \underline{\hspace{1.6cm}} \quad \textbf{Nota:} \underline{\hspace{0.9cm}}/10


\subsection*{PARTE 1: RECONHECIMENTO}

\begin{enumerate}

  \item Forme as sílabas com os sons da unidade:

\begin{center}
\resizebox{\textwidth}{!}{%
\begin{tabular}{ccccc}
  \sil{CH}{A} & \sil{CH}{E} & \sil{CH}{I} & \sil{CH}{O} & \sil{CH}{U} \\
  \end{tabular}}
\end{center}

  \item Complete a tabela de sílabas da unidade:

\begin{center}
\resizebox{\textwidth}{!}{%
\begin{tabular}{ccccc}
  \sil{CH}{A} & \sil{CH}{E} & \sil{CH}{I} & \sil{CH}{O} & \sil{CH}{U} \\
  \end{tabular}}
\end{center}

  \item Sublinhe as palavras da unidade:

\begin{center}
  \uline{chave} \quad \uline{chapéu} \quad \uline{chuva} \quad \uline{chão} \quad \uline{cheiro} \quad \uline{choro} \quad fecho \quad chá \quad gato \quad sapo \quad rato \quad bola \quad menino
\end{center}

  \item Identifique o som-alvo em cada palavra:
  \begin{center}
  CHAVE $\rightarrow$ som-alvo: \underline{\hspace{1cm}} \quad CHAPÉU $\rightarrow$ som-alvo: \underline{\hspace{1cm}} \quad CHUVA $\rightarrow$ som-alvo: \underline{\hspace{1cm}}
  \end{center}

  \item Separe em sílabas:
  \begin{center}
  CHAVE $\rightarrow$ \underline{\hspace{1cm}} \quad CHAPÉU $\rightarrow$ \underline{\hspace{1cm}} \quad CHUVA $\rightarrow$ \underline{\hspace{1cm}}
  \end{center}

  \item Conte as sílabas:
  \begin{center}
  CHAVE $\rightarrow$ \underline{\hspace{0.9cm}} \quad CHAPÉU $\rightarrow$ \underline{\hspace{0.9cm}} \quad CHUVA $\rightarrow$ \underline{\hspace{0.9cm}}
  \end{center}

\end{enumerate}

\subsection*{PARTE 2: PRODUÇÃO GUIADA}

\begin{enumerate}[resume]

  \item Complete com a grafia da unidade:
  \begin{center}
  CH\underline{~~~} \quad CH\underline{~~~} \quad CH\underline{~~~} \quad CH\underline{~~~} \quad CH\underline{~~~}
  \end{center}

  \item Escreva 3 palavras com o som-alvo:
  \begin{center}
  1. \underline{\hspace{3cm}} \quad 2. \underline{\hspace{3cm}} \quad 3. \underline{\hspace{3cm}}
  \end{center}

  \item Copie as palavras:
  \begin{center}
  chave \quad chapéu \quad chuva \quad chão \quad cheiro
  \end{center}

  \item Separe em sílabas:
  \begin{center}
  CHAVE = \underline{\hspace{1cm}} \underline{\hspace{1cm}} \quad CHAPÉU = \underline{\hspace{1cm}} \underline{\hspace{1cm}} \quad CHUVA = \underline{\hspace{1cm}} \underline{\hspace{1cm}}
  \end{center}

  \item Escreva 2 frases usando palavras da unidade:
  \begin{enumerate}[label=\alph*)]
    \item \underline{\hspace{8cm}}
    \item \underline{\hspace{8cm}}
  \end{enumerate}

\end{enumerate}

\end{exercicio}

%% PLANCHETA 1.2

\plancheta{Folha de Prática 1.2 --- O parágrafo (Parte 2)}

\begin{exercicio}[Folha de Prática 1.2 --- O parágrafo (Parte 2)]

\noindent\textbf{Nome:} \underline{\hspace{3.2cm}} \quad \textbf{Data:} \underline{\hspace{1.6cm}} \quad \textbf{Nota:} \underline{\hspace{0.9cm}}/10


\subsection*{PARTE 3: INTEGRAÇÃO}

\begin{enumerate}[resume]

  \item Leia as palavras em voz alta:
  \begin{center}
  $\rightarrow$ chave ~ chapéu ~ chuva ~ chão ~ cheiro ~ choro
  \end{center}

  \item Leia as frases em voz alta:
  \begin{center}
  $\rightarrow$ A chave abriu a porta.\\
  $\rightarrow$ A chuva molhou o chão.\\
  $\rightarrow$ O chapéu protege do sol.\\
  $\rightarrow$ Sinto o cheiro do bolo.\\
  \end{center}

  \item Circule a opção correta:
  \begin{enumerate}[label=(\alph*)]

    \item O som da unidade é: ( ) oral ( ) nasal 

    \item A primeira sílaba de CHAVE é: ( ) CHA ( ) CHE

    \item As palavras da unidade aparecem: ( ) no início ( ) no meio ( ) no fim das palavras

  \end{enumerate}

  \item Leia e complete:
  \begin{center}
  cha\underline{~~~} \quad cha\underline{~~~} \quad chu\underline{~~~} \quad chã\underline{~~~}
  \end{center}

  \item Leitura em voz alta com o professor:
  \begin{center}
  chave \quad chapéu \quad chuva \quad chão \quad cheiro \quad choro \quad fecho \quad chá
  \end{center}

\end{enumerate}

\end{exercicio}

%% PLANCHETA 1.3

\plancheta{Folha de Prática 1.3 --- Produção com CH}

\begin{exercicio}[Folha de Prática 1.3 --- Produção com CH]

\noindent\textbf{Nome:} \underline{\hspace{3.2cm}} \quad \textbf{Data:} \underline{\hspace{1.6cm}} \quad \textbf{Nota:} \underline{\hspace{0.9cm}}/15


\noindent\textbf{Comando:} Escrever três frases sobre o que você faz em um dia de chuva, usando palavras com CH.


\begin{enumerate}

  \item Planeje: escolha 3 palavras da unidade para usar:
  \begin{center}
  1. \underline{\hspace{3cm}} \quad 2. \underline{\hspace{3cm}} \quad 3. \underline{\hspace{3cm}}
  \end{center}

  \item Escreva o texto (mínimo 3 frases):
  \begin{center}
  \underline{\hspace{12cm}}\\\underline{\hspace{12cm}}\\\underline{\hspace{12cm}}\\\underline{\hspace{12cm}}\\\underline{\hspace{12cm}}
  \end{center}

  \item Releia e corrija: marque palavras com a grafia CH que você usou.

  \item Ilustre o texto:
  \begin{center}
  \begin{tikzpicture}
    \draw[thick, dashed] (0,0) rectangle (12,6);
  \end{tikzpicture}
  \end{center}

\end{enumerate}

\end{exercicio}

%% PLANCHETA 1.4

\plancheta{Folha de Prática 1.4 --- Ditado e Avaliação da Unidade 1}

\begin{exercicio}[Folha de Prática 1.4 --- Ditado e Avaliação da Unidade 1]

\noindent\textbf{Nome:} \underline{\hspace{3.2cm}} \quad \textbf{Data:} \underline{\hspace{1.6cm}} \quad \textbf{Nota:} \underline{\hspace{0.9cm}}/20


\subsection*{PARTE 1: DITADO (10 itens)}

O professor dita as palavras e frases abaixo; o aluno escreve com letra cursiva.

\begin{enumerate}

  \item \underline{\hspace{10.5cm}}

  \item \underline{\hspace{10.5cm}}

  \item \underline{\hspace{10.5cm}}

  \item \underline{\hspace{10.5cm}}

  \item \underline{\hspace{10.5cm}}

  \item \underline{\hspace{10.5cm}}

\end{enumerate}

\subsection*{PARTE 2: AUTOAVALIAÇÃO DO ALUNO}

\begin{enumerate}[label=\alph*)]

  \item Eu li as palavras da unidade sem ajuda.  \quad ( ) Sim   ( ) Não

  \item Eu escrevi o ditado com as grafias certas.  \quad ( ) Sim   ( ) Não

  \item Eu separei as palavras em sílabas.  \quad ( ) Sim   ( ) Não

  \item Eu sei explicar a regra com as minhas palavras.  \quad ( ) Sim   ( ) Não

\end{enumerate}

\subsection*{PARTE 3: REGISTRO DO PROFESSOR (S/F/R/N)}

\begin{center}
\renewcommand{\arraystretch}{1.4}
\begin{tabular}{@{}p{9cm}cccc@{}}
\toprule
Aspecto observado & S & F & R & N \\
\midrule
Lê com precisão as palavras-alvo & & & & \\
Escreve com ortografia estável & & & & \\
Generaliza a regra em escrita espontânea & & & & \\
Mantém atenção durante o ditado & & & & \\
\bottomrule\end{tabular}\end{center}

\end{exercicio}

\begin{dica}
**Dica para o Professor:** Faça o aluno comparar CHA (africado) com CA (oclusivo) no espelho; a diferença articulatória é o coração da lição.
\end{dica}

%% ============================================
%% UNIDADE 2 -- Ponto final e vírgula
%% ============================================

\input{checkpoints/rastreio-u1}
\chapter{Unidade 2 --- Ponto final e vírgula}

\metadata{I-02}{Silábico-Alfabético Avançado}{EF04LP01, EF04LP02, EF04LP03, EF04LP05, EF04LP06, EF04LP07, EF04LP09}{D1}{EF04LP01, EF04LP02, EF04LP03, EF04LP05, EF04LP06, EF04LP07, EF04LP09}{Resolucao CNE/CEB 7/2010, art. 12}

\textbf{Regra:} LH representa o som /lh/: ilha, coelho, orelha. A língua encosta no céu da boca e o som sai lateralmente.

\textbf{Palavra geradora do bloco:} \textbf{COELHO} --- o coelho é afetivo e conhecido; ajuda a fixar o som lateral do LH.

\section{Lição 2.1 --- Usar ponto final e vírgula na enumeração}

\textbf{Foco da lição:} Usar ponto final e vírgula na enumeração


\begin{boquinha}[Boquinha da Técnica]
Para LH, a ponta da língua sobe ao céu da boca e o ar sai pelas laterais, como um 'L molhado'. Observe o sorriso dos lábios.
\end{boquinha}

\noindent\textbf{Formação guiada:}

\begin{center}
\resizebox{\textwidth}{!}{%
\begin{tabular}{ccccc}
  \sil{LH}{A} & \sil{LH}{E} & \sil{LH}{I} & \sil{LH}{O} & \sil{LH}{U} \\
  \end{tabular}}
\end{center}

\noindent\textbf{Palavras da unidade:} ilha, coelho, orelha, milho, filho, folha, velho, telhado.

\noindent\textbf{Frases para leitura oral:}
\begin{itemize}[leftmargin=1.6em]
  \item O coelho comeu a folha.
  \item O filho mora na ilha.
  \item A orelha do elefante é grande.
  \item O velho consertou o telhado.
\end{itemize}

\section{Lição 2.2 --- Leitura do texto curto}

\noindent\textbf{Leia o texto em voz alta, com atenção às palavras da unidade:}


\begin{quote}
O coelho comeu a folha. O filho mora na ilha.

A orelha do elefante é grande. O velho consertou o telhado.
\end{quote}

\textbf{Depois de ler, responda por escrito:}
\begin{enumerate}

  \item O que acontece no texto que você leu?
  \begin{center}
  \underline{\hspace{12cm}}
  \end{center}

  \item Quantas frases o texto tem?
  \begin{center}
  \underline{\hspace{12cm}}
  \end{center}

  \item Circule no texto as palavras com a grafia LH.

  \item Leia o texto de novo em voz alta, agora sem ajuda.

\end{enumerate}

\section{Lição 2.3 --- Escrita com a grafia LH}

\noindent\textbf{Regra da unidade:} LH representa o som /lh/: ilha, coelho, orelha. A língua encosta no céu da boca e o som sai lateralmente.


\noindent\textbf{Ditado guiado:} o professor dita as palavras abaixo; o aluno escreve e depois corrige com o professor.

\begin{center}
  \underline{\hspace{2.3cm}} \quad \underline{\hspace{2.3cm}} \quad \underline{\hspace{2.3cm}} \quad \underline{\hspace{2.3cm}}

  \underline{\hspace{2.3cm}} \quad \underline{\hspace{2.3cm}} \quad \underline{\hspace{2.3cm}} \quad \underline{\hspace{2.3cm}}
\end{center}

\noindent\textbf{Complete as palavras com a grafia LH:}
\begin{center}
  ilh\underline{~} \quad coelh\underline{~} \quad orelh\underline{~} \quad milh\underline{~} \quad filh\underline{~}
\end{center}

\noindent\textbf{Escreva a palavra que o professor mostrar (banco de palavras) e separe em sílabas:}

\begin{center}
  \underline{\hspace{2cm}} $=$ \underline{\hspace{2cm}} \quad \underline{\hspace{2cm}} $=$ \underline{\hspace{2cm}} \quad \underline{\hspace{2cm}} $=$ \underline{\hspace{2cm}}
  \end{center}

%% PLANCHETA 2.1

\plancheta{Folha de Prática 2.1 --- Ponto final e vírgula (Parte 1)}

\begin{exercicio}[Folha de Prática 2.1 --- Ponto final e vírgula (Parte 1)]

\noindent\textbf{Nome:} \underline{\hspace{3.2cm}} \quad \textbf{Data:} \underline{\hspace{1.6cm}} \quad \textbf{Nota:} \underline{\hspace{0.9cm}}/10


\subsection*{PARTE 1: RECONHECIMENTO}

\begin{enumerate}

  \item Forme as sílabas com os sons da unidade:

\begin{center}
\resizebox{\textwidth}{!}{%
\begin{tabular}{ccccc}
  \sil{LH}{A} & \sil{LH}{E} & \sil{LH}{I} & \sil{LH}{O} & \sil{LH}{U} \\
  \end{tabular}}
\end{center}

  \item Complete a tabela de sílabas da unidade:

\begin{center}
\resizebox{\textwidth}{!}{%
\begin{tabular}{ccccc}
  \sil{LH}{A} & \sil{LH}{E} & \sil{LH}{I} & \sil{LH}{O} & \sil{LH}{U} \\
  \end{tabular}}
\end{center}

  \item Sublinhe as palavras da unidade:

\begin{center}
  \uline{ilha} \quad \uline{coelho} \quad \uline{orelha} \quad \uline{milho} \quad \uline{filho} \quad \uline{folha} \quad velho \quad telhado \quad gato \quad sapo \quad rato \quad bola \quad menino
\end{center}

  \item Identifique o som-alvo em cada palavra:
  \begin{center}
  ILHA $\rightarrow$ som-alvo: \underline{\hspace{1cm}} \quad COELHO $\rightarrow$ som-alvo: \underline{\hspace{1cm}} \quad ORELHA $\rightarrow$ som-alvo: \underline{\hspace{1cm}}
  \end{center}

  \item Separe em sílabas:
  \begin{center}
  ILHA $\rightarrow$ \underline{\hspace{1cm}} \quad COELHO $\rightarrow$ \underline{\hspace{1cm}} \quad ORELHA $\rightarrow$ \underline{\hspace{1cm}}
  \end{center}

  \item Conte as sílabas:
  \begin{center}
  ILHA $\rightarrow$ \underline{\hspace{0.9cm}} \quad COELHO $\rightarrow$ \underline{\hspace{0.9cm}} \quad ORELHA $\rightarrow$ \underline{\hspace{0.9cm}}
  \end{center}

\end{enumerate}

\subsection*{PARTE 2: PRODUÇÃO GUIADA}

\begin{enumerate}[resume]

  \item Complete com a grafia da unidade:
  \begin{center}
  LH\underline{~~~} \quad LH\underline{~~~} \quad LH\underline{~~~} \quad LH\underline{~~~} \quad LH\underline{~~~}
  \end{center}

  \item Escreva 3 palavras com o som-alvo:
  \begin{center}
  1. \underline{\hspace{3cm}} \quad 2. \underline{\hspace{3cm}} \quad 3. \underline{\hspace{3cm}}
  \end{center}

  \item Copie as palavras:
  \begin{center}
  ilha \quad coelho \quad orelha \quad milho \quad filho
  \end{center}

  \item Separe em sílabas:
  \begin{center}
  ILHA = \underline{\hspace{1cm}} \underline{\hspace{1cm}} \quad COELHO = \underline{\hspace{1cm}} \underline{\hspace{1cm}} \quad ORELHA = \underline{\hspace{1cm}} \underline{\hspace{1cm}}
  \end{center}

  \item Escreva 2 frases usando palavras da unidade:
  \begin{enumerate}[label=\alph*)]
    \item \underline{\hspace{8cm}}
    \item \underline{\hspace{8cm}}
  \end{enumerate}

\end{enumerate}

\end{exercicio}

%% PLANCHETA 2.2

\plancheta{Folha de Prática 2.2 --- Ponto final e vírgula (Parte 2)}

\begin{exercicio}[Folha de Prática 2.2 --- Ponto final e vírgula (Parte 2)]

\noindent\textbf{Nome:} \underline{\hspace{3.2cm}} \quad \textbf{Data:} \underline{\hspace{1.6cm}} \quad \textbf{Nota:} \underline{\hspace{0.9cm}}/10


\subsection*{PARTE 3: INTEGRAÇÃO}

\begin{enumerate}[resume]

  \item Leia as palavras em voz alta:
  \begin{center}
  $\rightarrow$ ilha ~ coelho ~ orelha ~ milho ~ filho ~ folha
  \end{center}

  \item Leia as frases em voz alta:
  \begin{center}
  $\rightarrow$ O coelho comeu a folha.\\
  $\rightarrow$ O filho mora na ilha.\\
  $\rightarrow$ A orelha do elefante é grande.\\
  $\rightarrow$ O velho consertou o telhado.\\
  \end{center}

  \item Circule a opção correta:
  \begin{enumerate}[label=(\alph*)]

    \item O som da unidade é: ( ) oral ( ) nasal 

    \item A primeira sílaba de ILHA é: ( ) LHA ( ) LHE

    \item As palavras da unidade aparecem: ( ) no início ( ) no meio ( ) no fim das palavras

  \end{enumerate}

  \item Leia e complete:
  \begin{center}
  ilh\underline{~~~} \quad coe\underline{~~~} \quad ore\underline{~~~} \quad mil\underline{~~~}
  \end{center}

  \item Leitura em voz alta com o professor:
  \begin{center}
  ilha \quad coelho \quad orelha \quad milho \quad filho \quad folha \quad velho \quad telhado
  \end{center}

\end{enumerate}

\end{exercicio}

%% PLANCHETA 2.3

\plancheta{Folha de Prática 2.3 --- Produção com LH}

\begin{exercicio}[Folha de Prática 2.3 --- Produção com LH]

\noindent\textbf{Nome:} \underline{\hspace{3.2cm}} \quad \textbf{Data:} \underline{\hspace{1.6cm}} \quad \textbf{Nota:} \underline{\hspace{0.9cm}}/15


\noindent\textbf{Comando:} Listar os integrantes da família e escrever uma frase para cada um, usando LH quando possível.


\begin{enumerate}

  \item Planeje: escolha 3 palavras da unidade para usar:
  \begin{center}
  1. \underline{\hspace{3cm}} \quad 2. \underline{\hspace{3cm}} \quad 3. \underline{\hspace{3cm}}
  \end{center}

  \item Escreva o texto (mínimo 3 frases):
  \begin{center}
  \underline{\hspace{12cm}}\\\underline{\hspace{12cm}}\\\underline{\hspace{12cm}}\\\underline{\hspace{12cm}}\\\underline{\hspace{12cm}}
  \end{center}

  \item Releia e corrija: marque palavras com a grafia LH que você usou.

  \item Ilustre o texto:
  \begin{center}
  \begin{tikzpicture}
    \draw[thick, dashed] (0,0) rectangle (12,6);
  \end{tikzpicture}
  \end{center}

\end{enumerate}

\end{exercicio}

%% PLANCHETA 2.4

\plancheta{Folha de Prática 2.4 --- Ditado e Avaliação da Unidade 2}

\begin{exercicio}[Folha de Prática 2.4 --- Ditado e Avaliação da Unidade 2]

\noindent\textbf{Nome:} \underline{\hspace{3.2cm}} \quad \textbf{Data:} \underline{\hspace{1.6cm}} \quad \textbf{Nota:} \underline{\hspace{0.9cm}}/20


\subsection*{PARTE 1: DITADO (10 itens)}

O professor dita as palavras e frases abaixo; o aluno escreve com letra cursiva.

\begin{enumerate}

  \item \underline{\hspace{10.5cm}}

  \item \underline{\hspace{10.5cm}}

  \item \underline{\hspace{10.5cm}}

  \item \underline{\hspace{10.5cm}}

  \item \underline{\hspace{10.5cm}}

  \item \underline{\hspace{10.5cm}}

\end{enumerate}

\subsection*{PARTE 2: AUTOAVALIAÇÃO DO ALUNO}

\begin{enumerate}[label=\alph*)]

  \item Eu li as palavras da unidade sem ajuda.  \quad ( ) Sim   ( ) Não

  \item Eu escrevi o ditado com as grafias certas.  \quad ( ) Sim   ( ) Não

  \item Eu separei as palavras em sílabas.  \quad ( ) Sim   ( ) Não

  \item Eu sei explicar a regra com as minhas palavras.  \quad ( ) Sim   ( ) Não

\end{enumerate}

\subsection*{PARTE 3: REGISTRO DO PROFESSOR (S/F/R/N)}

\begin{center}
\renewcommand{\arraystretch}{1.4}
\begin{tabular}{@{}p{9cm}cccc@{}}
\toprule
Aspecto observado & S & F & R & N \\
\midrule
Lê com precisão as palavras-alvo & & & & \\
Escreve com ortografia estável & & & & \\
Generaliza a regra em escrita espontânea & & & & \\
Mantém atenção durante o ditado & & & & \\
\bottomrule\end{tabular}\end{center}

\end{exercicio}

\begin{dica}
**Dica para o Professor:** LH exige apoio prolongado da língua no palato; volte ao espelho sempre que o aluno produzir um L comum.
\end{dica}

%% ============================================
%% UNIDADE 3 -- Substantivos próprio e comum
%% ============================================

\chapter{Unidade 3 --- Substantivos próprio e comum}

\metadata{I-03}{Silábico-Alfabético Avançado}{EF04LP01, EF04LP02, EF04LP03, EF04LP05, EF04LP06, EF04LP07, EF04LP09}{D1}{EF04LP01, EF04LP02, EF04LP03, EF04LP05, EF04LP06, EF04LP07, EF04LP09}{Resolucao CNE/CEB 7/2010, art. 12}

\textbf{Regra:} NH representa o som /nh/: ninho, banho, sonho. A língua inteira toca o céu da boca e o ar sai pelo nariz.

\textbf{Palavra geradora do bloco:} \textbf{NINHO} --- o ninho une o afeto (aves) e a forma nasal do som; é a palavra geradora perfeita.

\section{Lição 3.1 --- Diferenciar substantivos próprio e comum}

\textbf{Foco da lição:} Diferenciar substantivos próprio e comum


\begin{boquinha}[Boquinha da Técnica]
Para NH, a boca fica quase fechada, a língua toca o palato e o ar sai pelo nariz, como um sorriso nasal.
\end{boquinha}

\noindent\textbf{Formação guiada:}

\begin{center}
\resizebox{\textwidth}{!}{%
\begin{tabular}{ccccc}
  \sil{NH}{A} & \sil{NH}{E} & \sil{NH}{I} & \sil{NH}{O} & \sil{NH}{U} \\
  \end{tabular}}
\end{center}

\noindent\textbf{Palavras da unidade:} ninho, banho, sonho, pinho, galinha, minhoca, caminho, vinho.

\noindent\textbf{Frases para leitura oral:}
\begin{itemize}[leftmargin=1.6em]
  \item O ninho tem três ovos.
  \item Tomo banho de manhã.
  \item A galinha cisca no caminho.
  \item Eu sonho com o mar.
\end{itemize}

\section{Lição 3.2 --- Leitura do texto curto}

\noindent\textbf{Leia o texto em voz alta, com atenção às palavras da unidade:}


\begin{quote}
O ninho tem três ovos. Tomo banho de manhã.

A galinha cisca no caminho. Eu sonho com o mar.
\end{quote}

\textbf{Depois de ler, responda por escrito:}
\begin{enumerate}

  \item O que acontece no texto que você leu?
  \begin{center}
  \underline{\hspace{12cm}}
  \end{center}

  \item Quantas frases o texto tem?
  \begin{center}
  \underline{\hspace{12cm}}
  \end{center}

  \item Circule no texto as palavras com a grafia NH.

  \item Leia o texto de novo em voz alta, agora sem ajuda.

\end{enumerate}

\section{Lição 3.3 --- Escrita com a grafia NH}

\noindent\textbf{Regra da unidade:} NH representa o som /nh/: ninho, banho, sonho. A língua inteira toca o céu da boca e o ar sai pelo nariz.


\noindent\textbf{Ditado guiado:} o professor dita as palavras abaixo; o aluno escreve e depois corrige com o professor.

\begin{center}
  \underline{\hspace{2.3cm}} \quad \underline{\hspace{2.3cm}} \quad \underline{\hspace{2.3cm}} \quad \underline{\hspace{2.3cm}}

  \underline{\hspace{2.3cm}} \quad \underline{\hspace{2.3cm}} \quad \underline{\hspace{2.3cm}} \quad \underline{\hspace{2.3cm}}
\end{center}

\noindent\textbf{Complete as palavras com a grafia NH:}
\begin{center}
  ninh\underline{~} \quad banh\underline{~} \quad sonh\underline{~} \quad pinh\underline{~} \quad galinh\underline{~}
\end{center}

\noindent\textbf{Escreva a palavra que o professor mostrar (banco de palavras) e separe em sílabas:}

\begin{center}
  \underline{\hspace{2cm}} $=$ \underline{\hspace{2cm}} \quad \underline{\hspace{2cm}} $=$ \underline{\hspace{2cm}} \quad \underline{\hspace{2cm}} $=$ \underline{\hspace{2cm}}
  \end{center}

%% PLANCHETA 3.1

\plancheta{Folha de Prática 3.1 --- Substantivos próprio e comum (Parte 1)}

\begin{exercicio}[Folha de Prática 3.1 --- Substantivos próprio e comum (Parte 1)]

\noindent\textbf{Nome:} \underline{\hspace{3.2cm}} \quad \textbf{Data:} \underline{\hspace{1.6cm}} \quad \textbf{Nota:} \underline{\hspace{0.9cm}}/10


\subsection*{PARTE 1: RECONHECIMENTO}

\begin{enumerate}

  \item Forme as sílabas com os sons da unidade:

\begin{center}
\resizebox{\textwidth}{!}{%
\begin{tabular}{ccccc}
  \sil{NH}{A} & \sil{NH}{E} & \sil{NH}{I} & \sil{NH}{O} & \sil{NH}{U} \\
  \end{tabular}}
\end{center}

  \item Complete a tabela de sílabas da unidade:

\begin{center}
\resizebox{\textwidth}{!}{%
\begin{tabular}{ccccc}
  \sil{NH}{A} & \sil{NH}{E} & \sil{NH}{I} & \sil{NH}{O} & \sil{NH}{U} \\
  \end{tabular}}
\end{center}

  \item Sublinhe as palavras da unidade:

\begin{center}
  \uline{ninho} \quad \uline{banho} \quad \uline{sonho} \quad \uline{pinho} \quad \uline{galinha} \quad \uline{minhoca} \quad caminho \quad vinho \quad gato \quad sapo \quad rato \quad bola \quad menino
\end{center}

  \item Identifique o som-alvo em cada palavra:
  \begin{center}
  NINHO $\rightarrow$ som-alvo: \underline{\hspace{1cm}} \quad BANHO $\rightarrow$ som-alvo: \underline{\hspace{1cm}} \quad SONHO $\rightarrow$ som-alvo: \underline{\hspace{1cm}}
  \end{center}

  \item Separe em sílabas:
  \begin{center}
  NINHO $\rightarrow$ \underline{\hspace{1cm}} \quad BANHO $\rightarrow$ \underline{\hspace{1cm}} \quad SONHO $\rightarrow$ \underline{\hspace{1cm}}
  \end{center}

  \item Conte as sílabas:
  \begin{center}
  NINHO $\rightarrow$ \underline{\hspace{0.9cm}} \quad BANHO $\rightarrow$ \underline{\hspace{0.9cm}} \quad SONHO $\rightarrow$ \underline{\hspace{0.9cm}}
  \end{center}

\end{enumerate}

\subsection*{PARTE 2: PRODUÇÃO GUIADA}

\begin{enumerate}[resume]

  \item Complete com a grafia da unidade:
  \begin{center}
  NH\underline{~~~} \quad NH\underline{~~~} \quad NH\underline{~~~} \quad NH\underline{~~~} \quad NH\underline{~~~}
  \end{center}

  \item Escreva 3 palavras com o som-alvo:
  \begin{center}
  1. \underline{\hspace{3cm}} \quad 2. \underline{\hspace{3cm}} \quad 3. \underline{\hspace{3cm}}
  \end{center}

  \item Copie as palavras:
  \begin{center}
  ninho \quad banho \quad sonho \quad pinho \quad galinha
  \end{center}

  \item Separe em sílabas:
  \begin{center}
  NINHO = \underline{\hspace{1cm}} \underline{\hspace{1cm}} \quad BANHO = \underline{\hspace{1cm}} \underline{\hspace{1cm}} \quad SONHO = \underline{\hspace{1cm}} \underline{\hspace{1cm}}
  \end{center}

  \item Escreva 2 frases usando palavras da unidade:
  \begin{enumerate}[label=\alph*)]
    \item \underline{\hspace{8cm}}
    \item \underline{\hspace{8cm}}
  \end{enumerate}

\end{enumerate}

\end{exercicio}

%% PLANCHETA 3.2

\plancheta{Folha de Prática 3.2 --- Substantivos próprio e comum (Parte 2)}

\begin{exercicio}[Folha de Prática 3.2 --- Substantivos próprio e comum (Parte 2)]

\noindent\textbf{Nome:} \underline{\hspace{3.2cm}} \quad \textbf{Data:} \underline{\hspace{1.6cm}} \quad \textbf{Nota:} \underline{\hspace{0.9cm}}/10


\subsection*{PARTE 3: INTEGRAÇÃO}

\begin{enumerate}[resume]

  \item Leia as palavras em voz alta:
  \begin{center}
  $\rightarrow$ ninho ~ banho ~ sonho ~ pinho ~ galinha ~ minhoca
  \end{center}

  \item Leia as frases em voz alta:
  \begin{center}
  $\rightarrow$ O ninho tem três ovos.\\
  $\rightarrow$ Tomo banho de manhã.\\
  $\rightarrow$ A galinha cisca no caminho.\\
  $\rightarrow$ Eu sonho com o mar.\\
  \end{center}

  \item Circule a opção correta:
  \begin{enumerate}[label=(\alph*)]

    \item O som da unidade é: ( ) oral ( ) nasal 

    \item A primeira sílaba de NINHO é: ( ) NHA ( ) NHE

    \item As palavras da unidade aparecem: ( ) no início ( ) no meio ( ) no fim das palavras

  \end{enumerate}

  \item Leia e complete:
  \begin{center}
  nin\underline{~~~} \quad ban\underline{~~~} \quad son\underline{~~~} \quad pin\underline{~~~}
  \end{center}

  \item Leitura em voz alta com o professor:
  \begin{center}
  ninho \quad banho \quad sonho \quad pinho \quad galinha \quad minhoca \quad caminho \quad vinho
  \end{center}

\end{enumerate}

\end{exercicio}

%% PLANCHETA 3.3

\plancheta{Folha de Prática 3.3 --- Produção com NH}

\begin{exercicio}[Folha de Prática 3.3 --- Produção com NH]

\noindent\textbf{Nome:} \underline{\hspace{3.2cm}} \quad \textbf{Data:} \underline{\hspace{1.6cm}} \quad \textbf{Nota:} \underline{\hspace{0.9cm}}/15


\noindent\textbf{Comando:} Contar e escrever uma história curta usando pelo menos 5 palavras com NH.


\begin{enumerate}

  \item Planeje: escolha 3 palavras da unidade para usar:
  \begin{center}
  1. \underline{\hspace{3cm}} \quad 2. \underline{\hspace{3cm}} \quad 3. \underline{\hspace{3cm}}
  \end{center}

  \item Escreva o texto (mínimo 3 frases):
  \begin{center}
  \underline{\hspace{12cm}}\\\underline{\hspace{12cm}}\\\underline{\hspace{12cm}}\\\underline{\hspace{12cm}}\\\underline{\hspace{12cm}}
  \end{center}

  \item Releia e corrija: marque palavras com a grafia NH que você usou.

  \item Ilustre o texto:
  \begin{center}
  \begin{tikzpicture}
    \draw[thick, dashed] (0,0) rectangle (12,6);
  \end{tikzpicture}
  \end{center}

\end{enumerate}

\end{exercicio}

%% PLANCHETA 3.4

\plancheta{Folha de Prática 3.4 --- Ditado e Avaliação da Unidade 3}

\begin{exercicio}[Folha de Prática 3.4 --- Ditado e Avaliação da Unidade 3]

\noindent\textbf{Nome:} \underline{\hspace{3.2cm}} \quad \textbf{Data:} \underline{\hspace{1.6cm}} \quad \textbf{Nota:} \underline{\hspace{0.9cm}}/20


\subsection*{PARTE 1: DITADO (10 itens)}

O professor dita as palavras e frases abaixo; o aluno escreve com letra cursiva.

\begin{enumerate}

  \item \underline{\hspace{10.5cm}}

  \item \underline{\hspace{10.5cm}}

  \item \underline{\hspace{10.5cm}}

  \item \underline{\hspace{10.5cm}}

  \item \underline{\hspace{10.5cm}}

  \item \underline{\hspace{10.5cm}}

\end{enumerate}

\subsection*{PARTE 2: AUTOAVALIAÇÃO DO ALUNO}

\begin{enumerate}[label=\alph*)]

  \item Eu li as palavras da unidade sem ajuda.  \quad ( ) Sim   ( ) Não

  \item Eu escrevi o ditado com as grafias certas.  \quad ( ) Sim   ( ) Não

  \item Eu separei as palavras em sílabas.  \quad ( ) Sim   ( ) Não

  \item Eu sei explicar a regra com as minhas palavras.  \quad ( ) Sim   ( ) Não

\end{enumerate}

\subsection*{PARTE 3: REGISTRO DO PROFESSOR (S/F/R/N)}

\begin{center}
\renewcommand{\arraystretch}{1.4}
\begin{tabular}{@{}p{9cm}cccc@{}}
\toprule
Aspecto observado & S & F & R & N \\
\midrule
Lê com precisão as palavras-alvo & & & & \\
Escreve com ortografia estável & & & & \\
Generaliza a regra em escrita espontânea & & & & \\
Mantém atenção durante o ditado & & & & \\
\bottomrule\end{tabular}\end{center}

\end{exercicio}

%% PLANCHETA 3.5

\plancheta{Folha de Prática 3.5 --- Leitura e Fluência da Unidade 3}

\begin{exercicio}[Folha de Prática 3.5 --- Leitura e Fluência da Unidade 3]

\noindent\textbf{Nome:} \underline{\hspace{3.2cm}} \quad \textbf{Data:} \underline{\hspace{1.6cm}} \quad \textbf{Nota:} \underline{\hspace{0.9cm}}/10


\subsection*{PARTE 1: LEITURA EM VOZ ALTA}

Leia o texto abaixo três vezes; a cada leitura, marque um círculo no cronômetro do professor.

\begin{quote}
O ninho tem três ovos. Tomo banho de manhã.

A galinha cisca no caminho. Eu sonho com o mar.
\end{quote}

\begin{center}
1ª leitura: \underline{\hspace{1.5cm}} \quad 2ª leitura: \underline{\hspace{1.5cm}} \quad 3ª leitura: \underline{\hspace{1.5cm}}
\end{center}

\subsection*{PARTE 2: COMPREENSÃO}

\begin{enumerate}

  \item Quem são os personagens do texto?
  \begin{center}
  \underline{\hspace{12cm}}
  \end{center}

  \item Qual é a grafia-alvo que mais aparece no texto? Escreva três exemplos:
  \begin{center}
  \underline{\hspace{3.5cm}} \quad \underline{\hspace{3.5cm}} \quad \underline{\hspace{3.5cm}}
  \end{center}

  \item Você gostou do texto? Por quê?
  \begin{center}
  \underline{\hspace{12cm}}
  \end{center}

\end{enumerate}

\subsection*{PARTE 3: DESENHO}

Desenhe uma cena do texto:

\begin{center}
\begin{tikzpicture}
  \draw[thick, dashed] (0,0) rectangle (12,5);
\end{tikzpicture}
\end{center}

\end{exercicio}

\begin{dica}
**Dica para o Professor:** NH e LH são parecidos no gesto; compare no espelho a saída do ar (nariz para NH, laterais para LH).
\end{dica}

%% ============================================
%% UNIDADE 4 -- Adjetivos na descrição
%% ============================================

\input{checkpoints/rastreio-u2}
\chapter{Unidade 4 --- Adjetivos na descrição}

\metadata{I-04}{Silábico-Alfabético Avançado}{EF04LP01, EF04LP02, EF04LP03, EF04LP05, EF04LP06, EF04LP07, EF04LP09}{D1}{EF04LP01, EF04LP02, EF04LP03, EF04LP05, EF04LP06, EF04LP07, EF04LP09}{Resolucao CNE/CEB 7/2010, art. 12}

\textbf{Regra:} QU representa o som /k/ em que e qui (queijo, quilo). GU representa /g/ em que e gui (guerra, guia). O U não é pronunciado.

\textbf{Palavra geradora do bloco:} \textbf{QUEIJO} --- o queijo é familiar e permite explorar QUE/QUI com um alimento do cotidiano.

\section{Lição 4.1 --- Usar adjetivos para descrever personagens}

\textbf{Foco da lição:} Usar adjetivos para descrever personagens


\begin{boquinha}[Boquinha da Técnica]
Em QUE e QUI, a boca faz o som de K com os lábios arredondados; em GUE e GUI, o som de G com a mesma postura.
\end{boquinha}

\noindent\textbf{Formação guiada:}

\begin{center}
\resizebox{\textwidth}{!}{%
\begin{tabular}{cccc}
  \sil{QU}{E} & \sil{QU}{I} & \sil{GU}{E} & \sil{GU}{I} \\
  \end{tabular}}
\end{center}

\noindent\textbf{Palavras da unidade:} queijo, quilo, quente, guerra, guia, guitarra, quero, água.

\noindent\textbf{Frases para leitura oral:}
\begin{itemize}[leftmargin=1.6em]
  \item O queijo está na mesa.
  \item Eu quero um quilo de pão.
  \item A água está quente.
  \item A guitarra toca na festa.
\end{itemize}

\section{Lição 4.2 --- Leitura do texto curto}

\noindent\textbf{Leia o texto em voz alta, com atenção às palavras da unidade:}


\begin{quote}
O queijo está na mesa. Eu quero um quilo de pão.

A água está quente. A guitarra toca na festa.
\end{quote}

\textbf{Depois de ler, responda por escrito:}
\begin{enumerate}

  \item O que acontece no texto que você leu?
  \begin{center}
  \underline{\hspace{12cm}}
  \end{center}

  \item Quantas frases o texto tem?
  \begin{center}
  \underline{\hspace{12cm}}
  \end{center}

  \item Circule no texto as palavras com a grafia QU.

  \item Leia o texto de novo em voz alta, agora sem ajuda.

\end{enumerate}

\section{Lição 4.3 --- Escrita com a grafia QU}

\noindent\textbf{Regra da unidade:} QU representa o som /k/ em que e qui (queijo, quilo). GU representa /g/ em que e gui (guerra, guia). O U não é pronunciado.


\noindent\textbf{Ditado guiado:} o professor dita as palavras abaixo; o aluno escreve e depois corrige com o professor.

\begin{center}
  \underline{\hspace{2.3cm}} \quad \underline{\hspace{2.3cm}} \quad \underline{\hspace{2.3cm}} \quad \underline{\hspace{2.3cm}}

  \underline{\hspace{2.3cm}} \quad \underline{\hspace{2.3cm}} \quad \underline{\hspace{2.3cm}} \quad \underline{\hspace{2.3cm}}
\end{center}

\noindent\textbf{Complete as palavras com a grafia QU:}
\begin{center}
  \underline{~~~}eijo \quad \underline{~~~}ilo \quad \underline{~~~}ente \quad guerr\underline{~} \quad gui\underline{~}
\end{center}

\noindent\textbf{Escreva a palavra que o professor mostrar (banco de palavras) e separe em sílabas:}

\begin{center}
  \underline{\hspace{2cm}} $=$ \underline{\hspace{2cm}} \quad \underline{\hspace{2cm}} $=$ \underline{\hspace{2cm}} \quad \underline{\hspace{2cm}} $=$ \underline{\hspace{2cm}}
  \end{center}

%% PLANCHETA 4.1

\plancheta{Folha de Prática 4.1 --- Adjetivos na descrição (Parte 1)}

\begin{exercicio}[Folha de Prática 4.1 --- Adjetivos na descrição (Parte 1)]

\noindent\textbf{Nome:} \underline{\hspace{3.2cm}} \quad \textbf{Data:} \underline{\hspace{1.6cm}} \quad \textbf{Nota:} \underline{\hspace{0.9cm}}/10


\subsection*{PARTE 1: RECONHECIMENTO}

\begin{enumerate}

  \item Forme as sílabas com os sons da unidade:

\begin{center}
\resizebox{\textwidth}{!}{%
\begin{tabular}{cccc}
  \sil{QU}{E} & \sil{QU}{I} & \sil{GU}{E} & \sil{GU}{I} \\
  \end{tabular}}
\end{center}

  \item Complete a tabela de sílabas da unidade:

\begin{center}
\resizebox{\textwidth}{!}{%
\begin{tabular}{ccccc}
  \sil{QU}{A} & \sil{QU}{E} & \sil{QU}{I} & \sil{QU}{O} & \sil{QU}{U} \\
  \end{tabular}}
\end{center}

  \item Sublinhe as palavras da unidade:

\begin{center}
  \uline{queijo} \quad \uline{quilo} \quad \uline{quente} \quad \uline{guerra} \quad \uline{guia} \quad \uline{guitarra} \quad quero \quad água \quad gato \quad sapo \quad rato \quad bola \quad menino
\end{center}

  \item Identifique o som-alvo em cada palavra:
  \begin{center}
  QUEIJO $\rightarrow$ som-alvo: \underline{\hspace{1cm}} \quad QUILO $\rightarrow$ som-alvo: \underline{\hspace{1cm}} \quad QUENTE $\rightarrow$ som-alvo: \underline{\hspace{1cm}}
  \end{center}

  \item Separe em sílabas:
  \begin{center}
  QUEIJO $\rightarrow$ \underline{\hspace{1cm}} \quad QUILO $\rightarrow$ \underline{\hspace{1cm}} \quad QUENTE $\rightarrow$ \underline{\hspace{1cm}}
  \end{center}

  \item Conte as sílabas:
  \begin{center}
  QUEIJO $\rightarrow$ \underline{\hspace{0.9cm}} \quad QUILO $\rightarrow$ \underline{\hspace{0.9cm}} \quad QUENTE $\rightarrow$ \underline{\hspace{0.9cm}}
  \end{center}

\end{enumerate}

\subsection*{PARTE 2: PRODUÇÃO GUIADA}

\begin{enumerate}[resume]

  \item Complete com a grafia da unidade:
  \begin{center}
  QU\underline{~~~} \quad QU\underline{~~~} \quad QU\underline{~~~} \quad QU\underline{~~~} \quad QU\underline{~~~}
  \end{center}

  \item Escreva 3 palavras com o som-alvo:
  \begin{center}
  1. \underline{\hspace{3cm}} \quad 2. \underline{\hspace{3cm}} \quad 3. \underline{\hspace{3cm}}
  \end{center}

  \item Copie as palavras:
  \begin{center}
  queijo \quad quilo \quad quente \quad guerra \quad guia
  \end{center}

  \item Separe em sílabas:
  \begin{center}
  QUEIJO = \underline{\hspace{1cm}} \underline{\hspace{1cm}} \quad QUILO = \underline{\hspace{1cm}} \underline{\hspace{1cm}} \quad QUENTE = \underline{\hspace{1cm}} \underline{\hspace{1cm}}
  \end{center}

  \item Escreva 2 frases usando palavras da unidade:
  \begin{enumerate}[label=\alph*)]
    \item \underline{\hspace{8cm}}
    \item \underline{\hspace{8cm}}
  \end{enumerate}

\end{enumerate}

\end{exercicio}

%% PLANCHETA 4.2

\plancheta{Folha de Prática 4.2 --- Adjetivos na descrição (Parte 2)}

\begin{exercicio}[Folha de Prática 4.2 --- Adjetivos na descrição (Parte 2)]

\noindent\textbf{Nome:} \underline{\hspace{3.2cm}} \quad \textbf{Data:} \underline{\hspace{1.6cm}} \quad \textbf{Nota:} \underline{\hspace{0.9cm}}/10


\subsection*{PARTE 3: INTEGRAÇÃO}

\begin{enumerate}[resume]

  \item Leia as palavras em voz alta:
  \begin{center}
  $\rightarrow$ queijo ~ quilo ~ quente ~ guerra ~ guia ~ guitarra
  \end{center}

  \item Leia as frases em voz alta:
  \begin{center}
  $\rightarrow$ O queijo está na mesa.\\
  $\rightarrow$ Eu quero um quilo de pão.\\
  $\rightarrow$ A água está quente.\\
  $\rightarrow$ A guitarra toca na festa.\\
  \end{center}

  \item Circule a opção correta:
  \begin{enumerate}[label=(\alph*)]

    \item O som da unidade é: ( ) oral ( ) nasal 

    \item A primeira sílaba de QUEIJO é: ( ) QUE ( ) QUI

    \item As palavras da unidade aparecem: ( ) no início ( ) no meio ( ) no fim das palavras

  \end{enumerate}

  \item Leia e complete:
  \begin{center}
  que\underline{~~~} \quad qui\underline{~~~} \quad que\underline{~~~} \quad gue\underline{~~~}
  \end{center}

  \item Leitura em voz alta com o professor:
  \begin{center}
  queijo \quad quilo \quad quente \quad guerra \quad guia \quad guitarra \quad quero \quad água
  \end{center}

\end{enumerate}

\end{exercicio}

%% PLANCHETA 4.3

\plancheta{Folha de Prática 4.3 --- Produção com QU}

\begin{exercicio}[Folha de Prática 4.3 --- Produção com QU]

\noindent\textbf{Nome:} \underline{\hspace{3.2cm}} \quad \textbf{Data:} \underline{\hspace{1.6cm}} \quad \textbf{Nota:} \underline{\hspace{0.9cm}}/15


\noindent\textbf{Comando:} Escrever um convite de aniversário usando QUE e QUI e palavras com GU.


\begin{enumerate}

  \item Planeje: escolha 3 palavras da unidade para usar:
  \begin{center}
  1. \underline{\hspace{3cm}} \quad 2. \underline{\hspace{3cm}} \quad 3. \underline{\hspace{3cm}}
  \end{center}

  \item Escreva o texto (mínimo 3 frases):
  \begin{center}
  \underline{\hspace{12cm}}\\\underline{\hspace{12cm}}\\\underline{\hspace{12cm}}\\\underline{\hspace{12cm}}\\\underline{\hspace{12cm}}
  \end{center}

  \item Releia e corrija: marque palavras com a grafia QU que você usou.

  \item Ilustre o texto:
  \begin{center}
  \begin{tikzpicture}
    \draw[thick, dashed] (0,0) rectangle (12,6);
  \end{tikzpicture}
  \end{center}

\end{enumerate}

\end{exercicio}

%% PLANCHETA 4.4

\plancheta{Folha de Prática 4.4 --- Ditado e Avaliação da Unidade 4}

\begin{exercicio}[Folha de Prática 4.4 --- Ditado e Avaliação da Unidade 4]

\noindent\textbf{Nome:} \underline{\hspace{3.2cm}} \quad \textbf{Data:} \underline{\hspace{1.6cm}} \quad \textbf{Nota:} \underline{\hspace{0.9cm}}/20


\subsection*{PARTE 1: DITADO (10 itens)}

O professor dita as palavras e frases abaixo; o aluno escreve com letra cursiva.

\begin{enumerate}

  \item \underline{\hspace{10.5cm}}

  \item \underline{\hspace{10.5cm}}

  \item \underline{\hspace{10.5cm}}

  \item \underline{\hspace{10.5cm}}

  \item \underline{\hspace{10.5cm}}

  \item \underline{\hspace{10.5cm}}

\end{enumerate}

\subsection*{PARTE 2: AUTOAVALIAÇÃO DO ALUNO}

\begin{enumerate}[label=\alph*)]

  \item Eu li as palavras da unidade sem ajuda.  \quad ( ) Sim   ( ) Não

  \item Eu escrevi o ditado com as grafias certas.  \quad ( ) Sim   ( ) Não

  \item Eu separei as palavras em sílabas.  \quad ( ) Sim   ( ) Não

  \item Eu sei explicar a regra com as minhas palavras.  \quad ( ) Sim   ( ) Não

\end{enumerate}

\subsection*{PARTE 3: REGISTRO DO PROFESSOR (S/F/R/N)}

\begin{center}
\renewcommand{\arraystretch}{1.4}
\begin{tabular}{@{}p{9cm}cccc@{}}
\toprule
Aspecto observado & S & F & R & N \\
\midrule
Lê com precisão as palavras-alvo & & & & \\
Escreve com ortografia estável & & & & \\
Generaliza a regra em escrita espontânea & & & & \\
Mantém atenção durante o ditado & & & & \\
\bottomrule\end{tabular}\end{center}

\end{exercicio}

\begin{dica}
**Dica para o Professor:** Não confunda QUER com KER: explore a correspondência na leitura e na escrita; o treino Kumon fixa.
\end{dica}

%% ============================================
%% UNIDADE 5 -- Verbos no texto
%% ============================================

\chapter{Unidade 5 --- Verbos no texto}

\metadata{I-05}{Silábico-Alfabético Avançado}{EF04LP01, EF04LP02, EF04LP03, EF04LP05, EF04LP06, EF04LP07, EF04LP09}{D1}{EF04LP01, EF04LP02, EF04LP03, EF04LP05, EF04LP06, EF04LP07, EF04LP09}{Resolucao CNE/CEB 7/2010, art. 12}

\textbf{Regra:} Nos encontros consonantais, duas consoantes aparecem juntas e cada uma mantém seu som: braço, cravo, prato.

\textbf{Palavra geradora do bloco:} \textbf{PRATO} --- o prato está em todas as casas; o encontro PR é estável e fácil de ancorar.

\section{Lição 5.1 --- Reconhecer verbos de ação no texto}

\textbf{Foco da lição:} Reconhecer verbos de ação no texto


\begin{boquinha}[Boquinha da Técnica]
Ao dizer BR, os lábios começam juntos (B) e a língua vibra (R). Cada consoante do encontro é articulada, sem vogal no meio.
\end{boquinha}

\noindent\textbf{Formação guiada:}

\begin{center}
\resizebox{\textwidth}{!}{%
\begin{tabular}{ccccc}
  \sil{BR}{A} & \sil{BR}{O} & \sil{CR}{A} & \sil{CR}{O} & \sil{DR}{A} \\
  \sil{FR}{A} & \sil{FR}{O} & \sil{GR}{A} & \sil{PR}{A} & \sil{PR}{O} \\
  \sil{TR}{A} & \sil{VR}{A} \\
  \end{tabular}}
\end{center}

\noindent\textbf{Palavras da unidade:} braço, prato, trem, fruta, grade, cravo, dragão, livro.

\noindent\textbf{Frases para leitura oral:}
\begin{itemize}[leftmargin=1.6em]
  \item O braço dói hoje.
  \item O prato tem fruta.
  \item O trem chega cedo.
  \item O dragão cuspiu fogo.
\end{itemize}

\section{Lição 5.2 --- Leitura do texto curto}

\noindent\textbf{Leia o texto em voz alta, com atenção às palavras da unidade:}


\begin{quote}
O braço dói hoje. O prato tem fruta.

O trem chega cedo. O dragão cuspiu fogo.
\end{quote}

\textbf{Depois de ler, responda por escrito:}
\begin{enumerate}

  \item O que acontece no texto que você leu?
  \begin{center}
  \underline{\hspace{12cm}}
  \end{center}

  \item Quantas frases o texto tem?
  \begin{center}
  \underline{\hspace{12cm}}
  \end{center}

  \item Circule no texto as palavras com a grafia BR.

  \item Leia o texto de novo em voz alta, agora sem ajuda.

\end{enumerate}

\section{Lição 5.3 --- Escrita com a grafia BR}

\noindent\textbf{Regra da unidade:} Nos encontros consonantais, duas consoantes aparecem juntas e cada uma mantém seu som: braço, cravo, prato.


\noindent\textbf{Ditado guiado:} o professor dita as palavras abaixo; o aluno escreve e depois corrige com o professor.

\begin{center}
  \underline{\hspace{2.3cm}} \quad \underline{\hspace{2.3cm}} \quad \underline{\hspace{2.3cm}} \quad \underline{\hspace{2.3cm}}

  \underline{\hspace{2.3cm}} \quad \underline{\hspace{2.3cm}} \quad \underline{\hspace{2.3cm}} \quad \underline{\hspace{2.3cm}}
\end{center}

\noindent\textbf{Complete as palavras com a grafia BR:}
\begin{center}
  \underline{~~~}aço \quad prat\underline{~} \quad tre\underline{~} \quad frut\underline{~} \quad grad\underline{~}
\end{center}

\noindent\textbf{Escreva a palavra que o professor mostrar (banco de palavras) e separe em sílabas:}

\begin{center}
  \underline{\hspace{2cm}} $=$ \underline{\hspace{2cm}} \quad \underline{\hspace{2cm}} $=$ \underline{\hspace{2cm}} \quad \underline{\hspace{2cm}} $=$ \underline{\hspace{2cm}}
  \end{center}

%% PLANCHETA 5.1

\plancheta{Folha de Prática 5.1 --- Verbos no texto (Parte 1)}

\begin{exercicio}[Folha de Prática 5.1 --- Verbos no texto (Parte 1)]

\noindent\textbf{Nome:} \underline{\hspace{3.2cm}} \quad \textbf{Data:} \underline{\hspace{1.6cm}} \quad \textbf{Nota:} \underline{\hspace{0.9cm}}/10


\subsection*{PARTE 1: RECONHECIMENTO}

\begin{enumerate}

  \item Forme as sílabas com os sons da unidade:

\begin{center}
\resizebox{\textwidth}{!}{%
\begin{tabular}{ccccc}
  \sil{BR}{A} & \sil{BR}{O} & \sil{CR}{A} & \sil{CR}{O} & \sil{DR}{A} \\
  \sil{FR}{A} & \sil{FR}{O} & \sil{GR}{A} & \sil{PR}{A} & \sil{PR}{O} \\
  \sil{TR}{A} & \sil{VR}{A} \\
  \end{tabular}}
\end{center}

  \item Complete a tabela de sílabas da unidade:

\begin{center}
\resizebox{\textwidth}{!}{%
\begin{tabular}{ccccc}
  \sil{BR}{A} & \sil{BR}{E} & \sil{BR}{I} & \sil{BR}{O} & \sil{BR}{U} \\
  \end{tabular}}
\end{center}

  \item Sublinhe as palavras da unidade:

\begin{center}
  \uline{braço} \quad \uline{prato} \quad \uline{trem} \quad \uline{fruta} \quad \uline{grade} \quad \uline{cravo} \quad dragão \quad livro \quad gato \quad sapo \quad rato \quad bola \quad menino
\end{center}

  \item Identifique o som-alvo em cada palavra:
  \begin{center}
  BRAÇO $\rightarrow$ som-alvo: \underline{\hspace{1cm}} \quad PRATO $\rightarrow$ som-alvo: \underline{\hspace{1cm}} \quad TREM $\rightarrow$ som-alvo: \underline{\hspace{1cm}}
  \end{center}

  \item Separe em sílabas:
  \begin{center}
  BRAÇO $\rightarrow$ \underline{\hspace{1cm}} \quad PRATO $\rightarrow$ \underline{\hspace{1cm}} \quad TREM $\rightarrow$ \underline{\hspace{1cm}}
  \end{center}

  \item Conte as sílabas:
  \begin{center}
  BRAÇO $\rightarrow$ \underline{\hspace{0.9cm}} \quad PRATO $\rightarrow$ \underline{\hspace{0.9cm}} \quad TREM $\rightarrow$ \underline{\hspace{0.9cm}}
  \end{center}

\end{enumerate}

\subsection*{PARTE 2: PRODUÇÃO GUIADA}

\begin{enumerate}[resume]

  \item Complete com a grafia da unidade:
  \begin{center}
  BR\underline{~~~} \quad BR\underline{~~~} \quad BR\underline{~~~} \quad BR\underline{~~~} \quad BR\underline{~~~}
  \end{center}

  \item Escreva 3 palavras com o som-alvo:
  \begin{center}
  1. \underline{\hspace{3cm}} \quad 2. \underline{\hspace{3cm}} \quad 3. \underline{\hspace{3cm}}
  \end{center}

  \item Copie as palavras:
  \begin{center}
  braço \quad prato \quad trem \quad fruta \quad grade
  \end{center}

  \item Separe em sílabas:
  \begin{center}
  BRAÇO = \underline{\hspace{1cm}} \underline{\hspace{1cm}} \quad PRATO = \underline{\hspace{1cm}} \underline{\hspace{1cm}} \quad TREM = \underline{\hspace{1cm}} \underline{\hspace{1cm}}
  \end{center}

  \item Escreva 2 frases usando palavras da unidade:
  \begin{enumerate}[label=\alph*)]
    \item \underline{\hspace{8cm}}
    \item \underline{\hspace{8cm}}
  \end{enumerate}

\end{enumerate}

\end{exercicio}

%% PLANCHETA 5.2

\plancheta{Folha de Prática 5.2 --- Verbos no texto (Parte 2)}

\begin{exercicio}[Folha de Prática 5.2 --- Verbos no texto (Parte 2)]

\noindent\textbf{Nome:} \underline{\hspace{3.2cm}} \quad \textbf{Data:} \underline{\hspace{1.6cm}} \quad \textbf{Nota:} \underline{\hspace{0.9cm}}/10


\subsection*{PARTE 3: INTEGRAÇÃO}

\begin{enumerate}[resume]

  \item Leia as palavras em voz alta:
  \begin{center}
  $\rightarrow$ braço ~ prato ~ trem ~ fruta ~ grade ~ cravo
  \end{center}

  \item Leia as frases em voz alta:
  \begin{center}
  $\rightarrow$ O braço dói hoje.\\
  $\rightarrow$ O prato tem fruta.\\
  $\rightarrow$ O trem chega cedo.\\
  $\rightarrow$ O dragão cuspiu fogo.\\
  \end{center}

  \item Circule a opção correta:
  \begin{enumerate}[label=(\alph*)]

    \item O som da unidade é: ( ) oral ( ) nasal ( ) vibrante

    \item A primeira sílaba de BRAÇO é: ( ) BRA ( ) BRO

    \item As palavras da unidade aparecem: ( ) no início ( ) no meio ( ) no fim das palavras

  \end{enumerate}

  \item Leia e complete:
  \begin{center}
  bra\underline{~~~} \quad pra\underline{~~~} \quad tre\underline{~~~} \quad fru\underline{~~~}
  \end{center}

  \item Leitura em voz alta com o professor:
  \begin{center}
  braço \quad prato \quad trem \quad fruta \quad grade \quad cravo \quad dragão \quad livro
  \end{center}

\end{enumerate}

\end{exercicio}

%% PLANCHETA 5.3

\plancheta{Folha de Prática 5.3 --- Produção com BR}

\begin{exercicio}[Folha de Prática 5.3 --- Produção com BR]

\noindent\textbf{Nome:} \underline{\hspace{3.2cm}} \quad \textbf{Data:} \underline{\hspace{1.6cm}} \quad \textbf{Nota:} \underline{\hspace{0.9cm}}/15


\noindent\textbf{Comando:} Escrever uma receita simples usando palavras com encontros consonantais (ex.: brigadeiro com BR).


\begin{enumerate}

  \item Planeje: escolha 3 palavras da unidade para usar:
  \begin{center}
  1. \underline{\hspace{3cm}} \quad 2. \underline{\hspace{3cm}} \quad 3. \underline{\hspace{3cm}}
  \end{center}

  \item Escreva o texto (mínimo 3 frases):
  \begin{center}
  \underline{\hspace{12cm}}\\\underline{\hspace{12cm}}\\\underline{\hspace{12cm}}\\\underline{\hspace{12cm}}\\\underline{\hspace{12cm}}
  \end{center}

  \item Releia e corrija: marque palavras com a grafia BR que você usou.

  \item Ilustre o texto:
  \begin{center}
  \begin{tikzpicture}
    \draw[thick, dashed] (0,0) rectangle (12,6);
  \end{tikzpicture}
  \end{center}

\end{enumerate}

\end{exercicio}

%% PLANCHETA 5.4

\plancheta{Folha de Prática 5.4 --- Ditado e Avaliação da Unidade 5}

\begin{exercicio}[Folha de Prática 5.4 --- Ditado e Avaliação da Unidade 5]

\noindent\textbf{Nome:} \underline{\hspace{3.2cm}} \quad \textbf{Data:} \underline{\hspace{1.6cm}} \quad \textbf{Nota:} \underline{\hspace{0.9cm}}/20


\subsection*{PARTE 1: DITADO (10 itens)}

O professor dita as palavras e frases abaixo; o aluno escreve com letra cursiva.

\begin{enumerate}

  \item \underline{\hspace{10.5cm}}

  \item \underline{\hspace{10.5cm}}

  \item \underline{\hspace{10.5cm}}

  \item \underline{\hspace{10.5cm}}

  \item \underline{\hspace{10.5cm}}

  \item \underline{\hspace{10.5cm}}

\end{enumerate}

\subsection*{PARTE 2: AUTOAVALIAÇÃO DO ALUNO}

\begin{enumerate}[label=\alph*)]

  \item Eu li as palavras da unidade sem ajuda.  \quad ( ) Sim   ( ) Não

  \item Eu escrevi o ditado com as grafias certas.  \quad ( ) Sim   ( ) Não

  \item Eu separei as palavras em sílabas.  \quad ( ) Sim   ( ) Não

  \item Eu sei explicar a regra com as minhas palavras.  \quad ( ) Sim   ( ) Não

\end{enumerate}

\subsection*{PARTE 3: REGISTRO DO PROFESSOR (S/F/R/N)}

\begin{center}
\renewcommand{\arraystretch}{1.4}
\begin{tabular}{@{}p{9cm}cccc@{}}
\toprule
Aspecto observado & S & F & R & N \\
\midrule
Lê com precisão as palavras-alvo & & & & \\
Escreve com ortografia estável & & & & \\
Generaliza a regra em escrita espontânea & & & & \\
Mantém atenção durante o ditado & & & & \\
\bottomrule\end{tabular}\end{center}

\end{exercicio}

\begin{dica}
**Dica para o Professor:** Peça leitura lenta dos encontros (B-R-A) e depois rápida (BRA); o automatismo vem do treino diário Kumon.
\end{dica}

%% ============================================
%% UNIDADE 6 -- Ortografia: S, Z e X
%% ============================================

\input{checkpoints/rastreio-u3}
\chapter{Unidade 6 --- Ortografia: S, Z e X}

\metadata{I-06}{Silábico-Alfabético Avançado}{EF04LP01, EF04LP02, EF04LP03, EF04LP05, EF04LP06, EF04LP07, EF04LP09}{D1}{EF04LP01, EF04LP02, EF04LP03, EF04LP05, EF04LP06, EF04LP07, EF04LP09}{Resolucao CNE/CEB 7/2010, art. 12}

\textbf{Regra:} BL, CL, FL, GL, PL e TL são encontros com L: blusa, claro, flor, globo, placa, atlas. As duas consoantes se pronunciam juntas.

\textbf{Palavra geradora do bloco:} \textbf{FLOR} --- a flor é visual e afetiva; o encontro FL aparece em um dos primeiros desenhos de toda criança.

\section{Lição 6.1 --- Revisar a ortografia de S, Z e X}

\textbf{Foco da lição:} Revisar a ortografia de S, Z e X


\begin{boquinha}[Boquinha da Técnica]
Em BL, os lábios fecham (B) e a língua sobe no L: o som desliza. Compare BLA com BA: o L é a segunda consoante.
\end{boquinha}

\noindent\textbf{Formação guiada:}

\begin{center}
\resizebox{\textwidth}{!}{%
\begin{tabular}{ccccc}
  \sil{BL}{A} & \sil{BL}{O} & \sil{CL}{A} & \sil{CL}{O} & \sil{FL}{A} \\
  \sil{FL}{O} & \sil{GL}{O} & \sil{PL}{A} & \sil{PL}{O} & \sil{TL}{A} \\
  \end{tabular}}
\end{center}

\noindent\textbf{Palavras da unidade:} blusa, claro, flor, globo, placa, atlas, planta, flauta.

\noindent\textbf{Frases para leitura oral:}
\begin{itemize}[leftmargin=1.6em]
  \item A blusa é azul.
  \item O céu está claro.
  \item A flor abriu no jardim.
  \item A placa indica o caminho.
\end{itemize}

\section{Lição 6.2 --- Leitura do texto curto}

\noindent\textbf{Leia o texto em voz alta, com atenção às palavras da unidade:}


\begin{quote}
A blusa é azul. O céu está claro.

A flor abriu no jardim. A placa indica o caminho.
\end{quote}

\textbf{Depois de ler, responda por escrito:}
\begin{enumerate}

  \item O que acontece no texto que você leu?
  \begin{center}
  \underline{\hspace{12cm}}
  \end{center}

  \item Quantas frases o texto tem?
  \begin{center}
  \underline{\hspace{12cm}}
  \end{center}

  \item Circule no texto as palavras com a grafia BL.

  \item Leia o texto de novo em voz alta, agora sem ajuda.

\end{enumerate}

\section{Lição 6.3 --- Escrita com a grafia BL}

\noindent\textbf{Regra da unidade:} BL, CL, FL, GL, PL e TL são encontros com L: blusa, claro, flor, globo, placa, atlas. As duas consoantes se pronunciam juntas.


\noindent\textbf{Ditado guiado:} o professor dita as palavras abaixo; o aluno escreve e depois corrige com o professor.

\begin{center}
  \underline{\hspace{2.3cm}} \quad \underline{\hspace{2.3cm}} \quad \underline{\hspace{2.3cm}} \quad \underline{\hspace{2.3cm}}

  \underline{\hspace{2.3cm}} \quad \underline{\hspace{2.3cm}} \quad \underline{\hspace{2.3cm}} \quad \underline{\hspace{2.3cm}}
\end{center}

\noindent\textbf{Complete as palavras com a grafia BL:}
\begin{center}
  \underline{~~~}usa \quad clar\underline{~} \quad flo\underline{~} \quad glob\underline{~} \quad plac\underline{~}
\end{center}

\noindent\textbf{Escreva a palavra que o professor mostrar (banco de palavras) e separe em sílabas:}

\begin{center}
  \underline{\hspace{2cm}} $=$ \underline{\hspace{2cm}} \quad \underline{\hspace{2cm}} $=$ \underline{\hspace{2cm}} \quad \underline{\hspace{2cm}} $=$ \underline{\hspace{2cm}}
  \end{center}

%% PLANCHETA 6.1

\plancheta{Folha de Prática 6.1 --- Ortografia: S, Z e X (Parte 1)}

\begin{exercicio}[Folha de Prática 6.1 --- Ortografia: S, Z e X (Parte 1)]

\noindent\textbf{Nome:} \underline{\hspace{3.2cm}} \quad \textbf{Data:} \underline{\hspace{1.6cm}} \quad \textbf{Nota:} \underline{\hspace{0.9cm}}/10


\subsection*{PARTE 1: RECONHECIMENTO}

\begin{enumerate}

  \item Forme as sílabas com os sons da unidade:

\begin{center}
\resizebox{\textwidth}{!}{%
\begin{tabular}{ccccc}
  \sil{BL}{A} & \sil{BL}{O} & \sil{CL}{A} & \sil{CL}{O} & \sil{FL}{A} \\
  \sil{FL}{O} & \sil{GL}{O} & \sil{PL}{A} & \sil{PL}{O} & \sil{TL}{A} \\
  \end{tabular}}
\end{center}

  \item Complete a tabela de sílabas da unidade:

\begin{center}
\resizebox{\textwidth}{!}{%
\begin{tabular}{ccccc}
  \sil{BL}{A} & \sil{BL}{E} & \sil{BL}{I} & \sil{BL}{O} & \sil{BL}{U} \\
  \end{tabular}}
\end{center}

  \item Sublinhe as palavras da unidade:

\begin{center}
  \uline{blusa} \quad \uline{claro} \quad \uline{flor} \quad \uline{globo} \quad \uline{placa} \quad \uline{atlas} \quad planta \quad flauta \quad gato \quad sapo \quad rato \quad bola \quad menino
\end{center}

  \item Identifique o som-alvo em cada palavra:
  \begin{center}
  BLUSA $\rightarrow$ som-alvo: \underline{\hspace{1cm}} \quad CLARO $\rightarrow$ som-alvo: \underline{\hspace{1cm}} \quad FLOR $\rightarrow$ som-alvo: \underline{\hspace{1cm}}
  \end{center}

  \item Separe em sílabas:
  \begin{center}
  BLUSA $\rightarrow$ \underline{\hspace{1cm}} \quad CLARO $\rightarrow$ \underline{\hspace{1cm}} \quad FLOR $\rightarrow$ \underline{\hspace{1cm}}
  \end{center}

  \item Conte as sílabas:
  \begin{center}
  BLUSA $\rightarrow$ \underline{\hspace{0.9cm}} \quad CLARO $\rightarrow$ \underline{\hspace{0.9cm}} \quad FLOR $\rightarrow$ \underline{\hspace{0.9cm}}
  \end{center}

\end{enumerate}

\subsection*{PARTE 2: PRODUÇÃO GUIADA}

\begin{enumerate}[resume]

  \item Complete com a grafia da unidade:
  \begin{center}
  BL\underline{~~~} \quad BL\underline{~~~} \quad BL\underline{~~~} \quad BL\underline{~~~} \quad BL\underline{~~~}
  \end{center}

  \item Escreva 3 palavras com o som-alvo:
  \begin{center}
  1. \underline{\hspace{3cm}} \quad 2. \underline{\hspace{3cm}} \quad 3. \underline{\hspace{3cm}}
  \end{center}

  \item Copie as palavras:
  \begin{center}
  blusa \quad claro \quad flor \quad globo \quad placa
  \end{center}

  \item Separe em sílabas:
  \begin{center}
  BLUSA = \underline{\hspace{1cm}} \underline{\hspace{1cm}} \quad CLARO = \underline{\hspace{1cm}} \underline{\hspace{1cm}} \quad FLOR = \underline{\hspace{1cm}} \underline{\hspace{1cm}}
  \end{center}

  \item Escreva 2 frases usando palavras da unidade:
  \begin{enumerate}[label=\alph*)]
    \item \underline{\hspace{8cm}}
    \item \underline{\hspace{8cm}}
  \end{enumerate}

\end{enumerate}

\end{exercicio}

%% PLANCHETA 6.2

\plancheta{Folha de Prática 6.2 --- Ortografia: S, Z e X (Parte 2)}

\begin{exercicio}[Folha de Prática 6.2 --- Ortografia: S, Z e X (Parte 2)]

\noindent\textbf{Nome:} \underline{\hspace{3.2cm}} \quad \textbf{Data:} \underline{\hspace{1.6cm}} \quad \textbf{Nota:} \underline{\hspace{0.9cm}}/10


\subsection*{PARTE 3: INTEGRAÇÃO}

\begin{enumerate}[resume]

  \item Leia as palavras em voz alta:
  \begin{center}
  $\rightarrow$ blusa ~ claro ~ flor ~ globo ~ placa ~ atlas
  \end{center}

  \item Leia as frases em voz alta:
  \begin{center}
  $\rightarrow$ A blusa é azul.\\
  $\rightarrow$ O céu está claro.\\
  $\rightarrow$ A flor abriu no jardim.\\
  $\rightarrow$ A placa indica o caminho.\\
  \end{center}

  \item Circule a opção correta:
  \begin{enumerate}[label=(\alph*)]

    \item O som da unidade é: ( ) oral ( ) nasal 

    \item A primeira sílaba de BLUSA é: ( ) BLA ( ) BLO

    \item As palavras da unidade aparecem: ( ) no início ( ) no meio ( ) no fim das palavras

  \end{enumerate}

  \item Leia e complete:
  \begin{center}
  blu\underline{~~~} \quad cla\underline{~~~} \quad flo\underline{~~~} \quad glo\underline{~~~}
  \end{center}

  \item Leitura em voz alta com o professor:
  \begin{center}
  blusa \quad claro \quad flor \quad globo \quad placa \quad atlas \quad planta \quad flauta
  \end{center}

\end{enumerate}

\end{exercicio}

%% PLANCHETA 6.3

\plancheta{Folha de Prática 6.3 --- Produção com BL}

\begin{exercicio}[Folha de Prática 6.3 --- Produção com BL]

\noindent\textbf{Nome:} \underline{\hspace{3.2cm}} \quad \textbf{Data:} \underline{\hspace{1.6cm}} \quad \textbf{Nota:} \underline{\hspace{0.9cm}}/15


\noindent\textbf{Comando:} Descrever um jardim usando palavras com FL e PL.


\begin{enumerate}

  \item Planeje: escolha 3 palavras da unidade para usar:
  \begin{center}
  1. \underline{\hspace{3cm}} \quad 2. \underline{\hspace{3cm}} \quad 3. \underline{\hspace{3cm}}
  \end{center}

  \item Escreva o texto (mínimo 3 frases):
  \begin{center}
  \underline{\hspace{12cm}}\\\underline{\hspace{12cm}}\\\underline{\hspace{12cm}}\\\underline{\hspace{12cm}}\\\underline{\hspace{12cm}}
  \end{center}

  \item Releia e corrija: marque palavras com a grafia BL que você usou.

  \item Ilustre o texto:
  \begin{center}
  \begin{tikzpicture}
    \draw[thick, dashed] (0,0) rectangle (12,6);
  \end{tikzpicture}
  \end{center}

\end{enumerate}

\end{exercicio}

%% PLANCHETA 6.4

\plancheta{Folha de Prática 6.4 --- Ditado e Avaliação da Unidade 6}

\begin{exercicio}[Folha de Prática 6.4 --- Ditado e Avaliação da Unidade 6]

\noindent\textbf{Nome:} \underline{\hspace{3.2cm}} \quad \textbf{Data:} \underline{\hspace{1.6cm}} \quad \textbf{Nota:} \underline{\hspace{0.9cm}}/20


\subsection*{PARTE 1: DITADO (10 itens)}

O professor dita as palavras e frases abaixo; o aluno escreve com letra cursiva.

\begin{enumerate}

  \item \underline{\hspace{10.5cm}}

  \item \underline{\hspace{10.5cm}}

  \item \underline{\hspace{10.5cm}}

  \item \underline{\hspace{10.5cm}}

  \item \underline{\hspace{10.5cm}}

  \item \underline{\hspace{10.5cm}}

\end{enumerate}

\subsection*{PARTE 2: AUTOAVALIAÇÃO DO ALUNO}

\begin{enumerate}[label=\alph*)]

  \item Eu li as palavras da unidade sem ajuda.  \quad ( ) Sim   ( ) Não

  \item Eu escrevi o ditado com as grafias certas.  \quad ( ) Sim   ( ) Não

  \item Eu separei as palavras em sílabas.  \quad ( ) Sim   ( ) Não

  \item Eu sei explicar a regra com as minhas palavras.  \quad ( ) Sim   ( ) Não

\end{enumerate}

\subsection*{PARTE 3: REGISTRO DO PROFESSOR (S/F/R/N)}

\begin{center}
\renewcommand{\arraystretch}{1.4}
\begin{tabular}{@{}p{9cm}cccc@{}}
\toprule
Aspecto observado & S & F & R & N \\
\midrule
Lê com precisão as palavras-alvo & & & & \\
Escreve com ortografia estável & & & & \\
Generaliza a regra em escrita espontânea & & & & \\
Mantém atenção durante o ditado & & & & \\
\bottomrule\end{tabular}\end{center}

\end{exercicio}

%% PLANCHETA 6.5

\plancheta{Folha de Prática 6.5 --- Leitura e Fluência da Unidade 6}

\begin{exercicio}[Folha de Prática 6.5 --- Leitura e Fluência da Unidade 6]

\noindent\textbf{Nome:} \underline{\hspace{3.2cm}} \quad \textbf{Data:} \underline{\hspace{1.6cm}} \quad \textbf{Nota:} \underline{\hspace{0.9cm}}/10


\subsection*{PARTE 1: LEITURA EM VOZ ALTA}

Leia o texto abaixo três vezes; a cada leitura, marque um círculo no cronômetro do professor.

\begin{quote}
A blusa é azul. O céu está claro.

A flor abriu no jardim. A placa indica o caminho.
\end{quote}

\begin{center}
1ª leitura: \underline{\hspace{1.5cm}} \quad 2ª leitura: \underline{\hspace{1.5cm}} \quad 3ª leitura: \underline{\hspace{1.5cm}}
\end{center}

\subsection*{PARTE 2: COMPREENSÃO}

\begin{enumerate}

  \item Quem são os personagens do texto?
  \begin{center}
  \underline{\hspace{12cm}}
  \end{center}

  \item Qual é a grafia-alvo que mais aparece no texto? Escreva três exemplos:
  \begin{center}
  \underline{\hspace{3.5cm}} \quad \underline{\hspace{3.5cm}} \quad \underline{\hspace{3.5cm}}
  \end{center}

  \item Você gostou do texto? Por quê?
  \begin{center}
  \underline{\hspace{12cm}}
  \end{center}

\end{enumerate}

\subsection*{PARTE 3: DESENHO}

Desenhe uma cena do texto:

\begin{center}
\begin{tikzpicture}
  \draw[thick, dashed] (0,0) rectangle (12,5);
\end{tikzpicture}
\end{center}

\end{exercicio}

\begin{dica}
**Dica para o Professor:** BL é o mais frequente; alunos que trocam BL por B precisam de repetição com espelho e do apoio do Registro.
\end{dica}

%% ============================================
%% UNIDADE 7 -- Ortografia: G e J
%% ============================================

\chapter{Unidade 7 --- Ortografia: G e J}

\metadata{I-07}{Silábico-Alfabético Avançado}{EF04LP01, EF04LP02, EF04LP03, EF04LP05, EF04LP06, EF04LP07, EF04LP09}{D1}{EF04LP01, EF04LP02, EF04LP03, EF04LP05, EF04LP06, EF04LP07, EF04LP09}{Resolucao CNE/CEB 7/2010, art. 12}

\textbf{Regra:} RR e SS aparecem ENTRE vogais: carro, terra, pasar, osso. Entre vogais, R simples tem som fraco e RR tem som forte.

\textbf{Palavra geradora do bloco:} \textbf{CARRO} --- o carro é brinquedo universal; a dupla RR fixa a regra de que entre vogais o R forte se escreve RR.

\section{Lição 7.1 --- Revisar a ortografia de G e J}

\textbf{Foco da lição:} Revisar a ortografia de G e J


\begin{boquinha}[Boquinha da Técnica]
O R forte (RR) faz a língua vibrar atrás dos dentes; o S forte (SS) produz o som de cobra, com o ar passando com atrito.
\end{boquinha}

\noindent\textbf{Formação guiada:}

\begin{center}
\resizebox{\textwidth}{!}{%
\begin{tabular}{ccccc}
  \sil{RR}{A} & \sil{RR}{O} & \sil{RR}{E} & \sil{SS}{A} & \sil{SS}{O} \\
  \sil{SS}{E} \\
  \end{tabular}}
\end{center}

\noindent\textbf{Palavras da unidade:} carro, terra, barro, sorriso, passar, osso, assado, morro.

\noindent\textbf{Frases para leitura oral:}
\begin{itemize}[leftmargin=1.6em]
  \item O carro andou na terra.
  \item O sorriso da menina é bonito.
  \item O osso é do cachorro.
  \item O frango assado está quente.
\end{itemize}

\section{Lição 7.2 --- Leitura do texto curto}

\noindent\textbf{Leia o texto em voz alta, com atenção às palavras da unidade:}


\begin{quote}
O carro andou na terra. O sorriso da menina é bonito.

O osso é do cachorro. O frango assado está quente.
\end{quote}

\textbf{Depois de ler, responda por escrito:}
\begin{enumerate}

  \item O que acontece no texto que você leu?
  \begin{center}
  \underline{\hspace{12cm}}
  \end{center}

  \item Quantas frases o texto tem?
  \begin{center}
  \underline{\hspace{12cm}}
  \end{center}

  \item Circule no texto as palavras com a grafia RR.

  \item Leia o texto de novo em voz alta, agora sem ajuda.

\end{enumerate}

\section{Lição 7.3 --- Escrita com a grafia RR}

\noindent\textbf{Regra da unidade:} RR e SS aparecem ENTRE vogais: carro, terra, pasar, osso. Entre vogais, R simples tem som fraco e RR tem som forte.


\noindent\textbf{Ditado guiado:} o professor dita as palavras abaixo; o aluno escreve e depois corrige com o professor.

\begin{center}
  \underline{\hspace{2.3cm}} \quad \underline{\hspace{2.3cm}} \quad \underline{\hspace{2.3cm}} \quad \underline{\hspace{2.3cm}}

  \underline{\hspace{2.3cm}} \quad \underline{\hspace{2.3cm}} \quad \underline{\hspace{2.3cm}} \quad \underline{\hspace{2.3cm}}
\end{center}

\noindent\textbf{Complete as palavras com a grafia RR:}
\begin{center}
  carr\underline{~} \quad terr\underline{~} \quad barr\underline{~} \quad sorris\underline{~} \quad passa\underline{~}
\end{center}

\noindent\textbf{Escreva a palavra que o professor mostrar (banco de palavras) e separe em sílabas:}

\begin{center}
  \underline{\hspace{2cm}} $=$ \underline{\hspace{2cm}} \quad \underline{\hspace{2cm}} $=$ \underline{\hspace{2cm}} \quad \underline{\hspace{2cm}} $=$ \underline{\hspace{2cm}}
  \end{center}

%% PLANCHETA 7.1

\plancheta{Folha de Prática 7.1 --- Ortografia: G e J (Parte 1)}

\begin{exercicio}[Folha de Prática 7.1 --- Ortografia: G e J (Parte 1)]

\noindent\textbf{Nome:} \underline{\hspace{3.2cm}} \quad \textbf{Data:} \underline{\hspace{1.6cm}} \quad \textbf{Nota:} \underline{\hspace{0.9cm}}/10


\subsection*{PARTE 1: RECONHECIMENTO}

\begin{enumerate}

  \item Forme as sílabas com os sons da unidade:

\begin{center}
\resizebox{\textwidth}{!}{%
\begin{tabular}{ccccc}
  \sil{RR}{A} & \sil{RR}{O} & \sil{RR}{E} & \sil{SS}{A} & \sil{SS}{O} \\
  \sil{SS}{E} \\
  \end{tabular}}
\end{center}

  \item Complete a tabela de sílabas da unidade:

\begin{center}
\resizebox{\textwidth}{!}{%
\begin{tabular}{ccccc}
  \sil{RR}{A} & \sil{RR}{E} & \sil{RR}{I} & \sil{RR}{O} & \sil{RR}{U} \\
  \end{tabular}}
\end{center}

  \item Sublinhe as palavras da unidade:

\begin{center}
  \uline{carro} \quad \uline{terra} \quad \uline{barro} \quad \uline{sorriso} \quad \uline{passar} \quad \uline{osso} \quad assado \quad morro \quad gato \quad sapo \quad rato \quad bola \quad menino
\end{center}

  \item Identifique o som-alvo em cada palavra:
  \begin{center}
  CARRO $\rightarrow$ som-alvo: \underline{\hspace{1cm}} \quad TERRA $\rightarrow$ som-alvo: \underline{\hspace{1cm}} \quad BARRO $\rightarrow$ som-alvo: \underline{\hspace{1cm}}
  \end{center}

  \item Separe em sílabas:
  \begin{center}
  CARRO $\rightarrow$ \underline{\hspace{1cm}} \quad TERRA $\rightarrow$ \underline{\hspace{1cm}} \quad BARRO $\rightarrow$ \underline{\hspace{1cm}}
  \end{center}

  \item Conte as sílabas:
  \begin{center}
  CARRO $\rightarrow$ \underline{\hspace{0.9cm}} \quad TERRA $\rightarrow$ \underline{\hspace{0.9cm}} \quad BARRO $\rightarrow$ \underline{\hspace{0.9cm}}
  \end{center}

\end{enumerate}

\subsection*{PARTE 2: PRODUÇÃO GUIADA}

\begin{enumerate}[resume]

  \item Complete com a grafia da unidade:
  \begin{center}
  RR\underline{~~~} \quad RR\underline{~~~} \quad RR\underline{~~~} \quad RR\underline{~~~} \quad RR\underline{~~~}
  \end{center}

  \item Escreva 3 palavras com o som-alvo:
  \begin{center}
  1. \underline{\hspace{3cm}} \quad 2. \underline{\hspace{3cm}} \quad 3. \underline{\hspace{3cm}}
  \end{center}

  \item Copie as palavras:
  \begin{center}
  carro \quad terra \quad barro \quad sorriso \quad passar
  \end{center}

  \item Separe em sílabas:
  \begin{center}
  CARRO = \underline{\hspace{1cm}} \underline{\hspace{1cm}} \quad TERRA = \underline{\hspace{1cm}} \underline{\hspace{1cm}} \quad BARRO = \underline{\hspace{1cm}} \underline{\hspace{1cm}}
  \end{center}

  \item Escreva 2 frases usando palavras da unidade:
  \begin{enumerate}[label=\alph*)]
    \item \underline{\hspace{8cm}}
    \item \underline{\hspace{8cm}}
  \end{enumerate}

\end{enumerate}

\end{exercicio}

%% PLANCHETA 7.2

\plancheta{Folha de Prática 7.2 --- Ortografia: G e J (Parte 2)}

\begin{exercicio}[Folha de Prática 7.2 --- Ortografia: G e J (Parte 2)]

\noindent\textbf{Nome:} \underline{\hspace{3.2cm}} \quad \textbf{Data:} \underline{\hspace{1.6cm}} \quad \textbf{Nota:} \underline{\hspace{0.9cm}}/10


\subsection*{PARTE 3: INTEGRAÇÃO}

\begin{enumerate}[resume]

  \item Leia as palavras em voz alta:
  \begin{center}
  $\rightarrow$ carro ~ terra ~ barro ~ sorriso ~ passar ~ osso
  \end{center}

  \item Leia as frases em voz alta:
  \begin{center}
  $\rightarrow$ O carro andou na terra.\\
  $\rightarrow$ O sorriso da menina é bonito.\\
  $\rightarrow$ O osso é do cachorro.\\
  $\rightarrow$ O frango assado está quente.\\
  \end{center}

  \item Circule a opção correta:
  \begin{enumerate}[label=(\alph*)]

    \item O som da unidade é: ( ) oral ( ) nasal ( ) vibrante

    \item A primeira sílaba de CARRO é: ( ) RRA ( ) RRO

    \item As palavras da unidade aparecem: ( ) no início ( ) no meio ( ) no fim das palavras

  \end{enumerate}

  \item Leia e complete:
  \begin{center}
  car\underline{~~~} \quad ter\underline{~~~} \quad bar\underline{~~~} \quad sor\underline{~~~}
  \end{center}

  \item Leitura em voz alta com o professor:
  \begin{center}
  carro \quad terra \quad barro \quad sorriso \quad passar \quad osso \quad assado \quad morro
  \end{center}

\end{enumerate}

\end{exercicio}

%% PLANCHETA 7.3

\plancheta{Folha de Prática 7.3 --- Produção com RR}

\begin{exercicio}[Folha de Prática 7.3 --- Produção com RR]

\noindent\textbf{Nome:} \underline{\hspace{3.2cm}} \quad \textbf{Data:} \underline{\hspace{1.6cm}} \quad \textbf{Nota:} \underline{\hspace{0.9cm}}/15


\noindent\textbf{Comando:} Escrever uma lista de lugares e objetos que tenham RR ou SS.


\begin{enumerate}

  \item Planeje: escolha 3 palavras da unidade para usar:
  \begin{center}
  1. \underline{\hspace{3cm}} \quad 2. \underline{\hspace{3cm}} \quad 3. \underline{\hspace{3cm}}
  \end{center}

  \item Escreva o texto (mínimo 3 frases):
  \begin{center}
  \underline{\hspace{12cm}}\\\underline{\hspace{12cm}}\\\underline{\hspace{12cm}}\\\underline{\hspace{12cm}}\\\underline{\hspace{12cm}}
  \end{center}

  \item Releia e corrija: marque palavras com a grafia RR que você usou.

  \item Ilustre o texto:
  \begin{center}
  \begin{tikzpicture}
    \draw[thick, dashed] (0,0) rectangle (12,6);
  \end{tikzpicture}
  \end{center}

\end{enumerate}

\end{exercicio}

%% PLANCHETA 7.4

\plancheta{Folha de Prática 7.4 --- Ditado e Avaliação da Unidade 7}

\begin{exercicio}[Folha de Prática 7.4 --- Ditado e Avaliação da Unidade 7]

\noindent\textbf{Nome:} \underline{\hspace{3.2cm}} \quad \textbf{Data:} \underline{\hspace{1.6cm}} \quad \textbf{Nota:} \underline{\hspace{0.9cm}}/20


\subsection*{PARTE 1: DITADO (10 itens)}

O professor dita as palavras e frases abaixo; o aluno escreve com letra cursiva.

\begin{enumerate}

  \item \underline{\hspace{10.5cm}}

  \item \underline{\hspace{10.5cm}}

  \item \underline{\hspace{10.5cm}}

  \item \underline{\hspace{10.5cm}}

  \item \underline{\hspace{10.5cm}}

  \item \underline{\hspace{10.5cm}}

\end{enumerate}

\subsection*{PARTE 2: AUTOAVALIAÇÃO DO ALUNO}

\begin{enumerate}[label=\alph*)]

  \item Eu li as palavras da unidade sem ajuda.  \quad ( ) Sim   ( ) Não

  \item Eu escrevi o ditado com as grafias certas.  \quad ( ) Sim   ( ) Não

  \item Eu separei as palavras em sílabas.  \quad ( ) Sim   ( ) Não

  \item Eu sei explicar a regra com as minhas palavras.  \quad ( ) Sim   ( ) Não

\end{enumerate}

\subsection*{PARTE 3: REGISTRO DO PROFESSOR (S/F/R/N)}

\begin{center}
\renewcommand{\arraystretch}{1.4}
\begin{tabular}{@{}p{9cm}cccc@{}}
\toprule
Aspecto observado & S & F & R & N \\
\midrule
Lê com precisão as palavras-alvo & & & & \\
Escreve com ortografia estável & & & & \\
Generaliza a regra em escrita espontânea & & & & \\
Mantém atenção durante o ditado & & & & \\
\bottomrule\end{tabular}\end{center}

\end{exercicio}

\begin{dica}
**Dica para o Professor:** Mostre o contraste R (muro) × RR (morro) no espelho; o som forte de RR usa mais vibração da língua.
\end{dica}

%% ============================================
%% UNIDADE 8 -- Ortografia: Ç, SS e C
%% ============================================

\input{checkpoints/rastreio-u4}
\chapter{Unidade 8 --- Ortografia: Ç, SS e C}

\metadata{I-08}{Silábico-Alfabético Avançado}{EF04LP01, EF04LP02, EF04LP03, EF04LP05, EF04LP06, EF04LP07, EF04LP09}{D1}{EF04LP01, EF04LP02, EF04LP03, EF04LP05, EF04LP06, EF04LP07, EF04LP09}{Resolucao CNE/CEB 7/2010, art. 12}

\textbf{Regra:} Entre vogais, S tem som de Z: casa, mesa, asa. Z representa o som de Z: zero, zebra. SS, entre vogais, tem som de S forte.

\textbf{Palavra geradora do bloco:} \textbf{CASA} --- a casa é a primeira palavra com S entre vogais que a criança reconhece; o som de Z nela surpreende e fixa a regra.

\section{Lição 8.1 --- Revisar a ortografia de Ç, SS e C}

\textbf{Foco da lição:} Revisar a ortografia de Ç, SS e C


\begin{boquinha}[Boquinha da Técnica]
O Z faz vibrar a voz com o ar passando; o S entre vogais vira Z; o SS mantém o som de cobra. Compare ZA, SA e SSA.
\end{boquinha}

\noindent\textbf{Formação guiada:}

\begin{center}
\resizebox{\textwidth}{!}{%
\begin{tabular}{ccccc}
  \sil{Z}{A} & \sil{Z}{E} & \sil{Z}{I} & \sil{Z}{O} & \sil{Z}{U} \\
  \sil{S}{A} & \sil{S}{E} & \sil{S}{I} \\
  \end{tabular}}
\end{center}

\noindent\textbf{Palavras da unidade:} casa, mesa, asa, razão, zero, zebra, poesia, visita.

\noindent\textbf{Frases para leitura oral:}
\begin{itemize}[leftmargin=1.6em]
  \item A casa é na rua da escola.
  \item A mesa tem quatro cadeiras.
  \item A zebra corre na savana.
  \item A poesia é bonita.
\end{itemize}

\section{Lição 8.2 --- Leitura do texto curto}

\noindent\textbf{Leia o texto em voz alta, com atenção às palavras da unidade:}


\begin{quote}
A casa é na rua da escola. A mesa tem quatro cadeiras.

A zebra corre na savana. A poesia é bonita.
\end{quote}

\textbf{Depois de ler, responda por escrito:}
\begin{enumerate}

  \item O que acontece no texto que você leu?
  \begin{center}
  \underline{\hspace{12cm}}
  \end{center}

  \item Quantas frases o texto tem?
  \begin{center}
  \underline{\hspace{12cm}}
  \end{center}

  \item Circule no texto as palavras com a grafia S/Z.

  \item Leia o texto de novo em voz alta, agora sem ajuda.

\end{enumerate}

\section{Lição 8.3 --- Escrita com a grafia S/Z}

\noindent\textbf{Regra da unidade:} Entre vogais, S tem som de Z: casa, mesa, asa. Z representa o som de Z: zero, zebra. SS, entre vogais, tem som de S forte.


\noindent\textbf{Ditado guiado:} o professor dita as palavras abaixo; o aluno escreve e depois corrige com o professor.

\begin{center}
  \underline{\hspace{2.3cm}} \quad \underline{\hspace{2.3cm}} \quad \underline{\hspace{2.3cm}} \quad \underline{\hspace{2.3cm}}

  \underline{\hspace{2.3cm}} \quad \underline{\hspace{2.3cm}} \quad \underline{\hspace{2.3cm}} \quad \underline{\hspace{2.3cm}}
\end{center}

\noindent\textbf{Complete as palavras com a grafia S/Z:}
\begin{center}
  cas\underline{~} \quad mes\underline{~} \quad as\underline{~} \quad razã\underline{~} \quad zer\underline{~}
\end{center}

\noindent\textbf{Escreva a palavra que o professor mostrar (banco de palavras) e separe em sílabas:}

\begin{center}
  \underline{\hspace{2cm}} $=$ \underline{\hspace{2cm}} \quad \underline{\hspace{2cm}} $=$ \underline{\hspace{2cm}} \quad \underline{\hspace{2cm}} $=$ \underline{\hspace{2cm}}
  \end{center}

%% PLANCHETA 8.1

\plancheta{Folha de Prática 8.1 --- Ortografia: Ç, SS e C (Parte 1)}

\begin{exercicio}[Folha de Prática 8.1 --- Ortografia: Ç, SS e C (Parte 1)]

\noindent\textbf{Nome:} \underline{\hspace{3.2cm}} \quad \textbf{Data:} \underline{\hspace{1.6cm}} \quad \textbf{Nota:} \underline{\hspace{0.9cm}}/10


\subsection*{PARTE 1: RECONHECIMENTO}

\begin{enumerate}

  \item Forme as sílabas com os sons da unidade:

\begin{center}
\resizebox{\textwidth}{!}{%
\begin{tabular}{ccccc}
  \sil{Z}{A} & \sil{Z}{E} & \sil{Z}{I} & \sil{Z}{O} & \sil{Z}{U} \\
  \sil{S}{A} & \sil{S}{E} & \sil{S}{I} \\
  \end{tabular}}
\end{center}

  \item Complete a tabela de sílabas da unidade:

\begin{center}
\resizebox{\textwidth}{!}{%
\begin{tabular}{ccccc}
  \sil{S}{A} & \sil{S}{E} & \sil{S}{I} & \sil{S}{O} & \sil{S}{U} \\
  \end{tabular}}
\end{center}

  \item Sublinhe as palavras da unidade:

\begin{center}
  \uline{casa} \quad \uline{mesa} \quad \uline{asa} \quad \uline{razão} \quad \uline{zero} \quad \uline{zebra} \quad poesia \quad visita \quad gato \quad sapo \quad rato \quad bola \quad menino
\end{center}

  \item Identifique o som-alvo em cada palavra:
  \begin{center}
  CASA $\rightarrow$ som-alvo: \underline{\hspace{1cm}} \quad MESA $\rightarrow$ som-alvo: \underline{\hspace{1cm}} \quad ASA $\rightarrow$ som-alvo: \underline{\hspace{1cm}}
  \end{center}

  \item Separe em sílabas:
  \begin{center}
  CASA $\rightarrow$ \underline{\hspace{1cm}} \quad MESA $\rightarrow$ \underline{\hspace{1cm}} \quad ASA $\rightarrow$ \underline{\hspace{1cm}}
  \end{center}

  \item Conte as sílabas:
  \begin{center}
  CASA $\rightarrow$ \underline{\hspace{0.9cm}} \quad MESA $\rightarrow$ \underline{\hspace{0.9cm}} \quad ASA $\rightarrow$ \underline{\hspace{0.9cm}}
  \end{center}

\end{enumerate}

\subsection*{PARTE 2: PRODUÇÃO GUIADA}

\begin{enumerate}[resume]

  \item Complete com a grafia da unidade:
  \begin{center}
  S/Z\underline{~~~} \quad S/Z\underline{~~~} \quad S/Z\underline{~~~} \quad S/Z\underline{~~~} \quad S/Z\underline{~~~}
  \end{center}

  \item Escreva 3 palavras com o som-alvo:
  \begin{center}
  1. \underline{\hspace{3cm}} \quad 2. \underline{\hspace{3cm}} \quad 3. \underline{\hspace{3cm}}
  \end{center}

  \item Copie as palavras:
  \begin{center}
  casa \quad mesa \quad asa \quad razão \quad zero
  \end{center}

  \item Separe em sílabas:
  \begin{center}
  CASA = \underline{\hspace{1cm}} \underline{\hspace{1cm}} \quad MESA = \underline{\hspace{1cm}} \underline{\hspace{1cm}} \quad ASA = \underline{\hspace{1cm}} \underline{\hspace{1cm}}
  \end{center}

  \item Escreva 2 frases usando palavras da unidade:
  \begin{enumerate}[label=\alph*)]
    \item \underline{\hspace{8cm}}
    \item \underline{\hspace{8cm}}
  \end{enumerate}

\end{enumerate}

\end{exercicio}

%% PLANCHETA 8.2

\plancheta{Folha de Prática 8.2 --- Ortografia: Ç, SS e C (Parte 2)}

\begin{exercicio}[Folha de Prática 8.2 --- Ortografia: Ç, SS e C (Parte 2)]

\noindent\textbf{Nome:} \underline{\hspace{3.2cm}} \quad \textbf{Data:} \underline{\hspace{1.6cm}} \quad \textbf{Nota:} \underline{\hspace{0.9cm}}/10


\subsection*{PARTE 3: INTEGRAÇÃO}

\begin{enumerate}[resume]

  \item Leia as palavras em voz alta:
  \begin{center}
  $\rightarrow$ casa ~ mesa ~ asa ~ razão ~ zero ~ zebra
  \end{center}

  \item Leia as frases em voz alta:
  \begin{center}
  $\rightarrow$ A casa é na rua da escola.\\
  $\rightarrow$ A mesa tem quatro cadeiras.\\
  $\rightarrow$ A zebra corre na savana.\\
  $\rightarrow$ A poesia é bonita.\\
  \end{center}

  \item Circule a opção correta:
  \begin{enumerate}[label=(\alph*)]

    \item O som da unidade é: ( ) oral ( ) nasal 

    \item A primeira sílaba de CASA é: ( ) SA ( ) ZA

    \item As palavras da unidade aparecem: ( ) no início ( ) no meio ( ) no fim das palavras

  \end{enumerate}

  \item Leia e complete:
  \begin{center}
  cas\underline{~~~} \quad mes\underline{~~~} \quad asa\underline{~~~} \quad raz\underline{~~~}
  \end{center}

  \item Leitura em voz alta com o professor:
  \begin{center}
  casa \quad mesa \quad asa \quad razão \quad zero \quad zebra \quad poesia \quad visita
  \end{center}

\end{enumerate}

\end{exercicio}

%% PLANCHETA 8.3

\plancheta{Folha de Prática 8.3 --- Produção com S/Z}

\begin{exercicio}[Folha de Prática 8.3 --- Produção com S/Z]

\noindent\textbf{Nome:} \underline{\hspace{3.2cm}} \quad \textbf{Data:} \underline{\hspace{1.6cm}} \quad \textbf{Nota:} \underline{\hspace{0.9cm}}/15


\noindent\textbf{Comando:} Escrever uma poesia curta usando palavras com S entre vogais.


\begin{enumerate}

  \item Planeje: escolha 3 palavras da unidade para usar:
  \begin{center}
  1. \underline{\hspace{3cm}} \quad 2. \underline{\hspace{3cm}} \quad 3. \underline{\hspace{3cm}}
  \end{center}

  \item Escreva o texto (mínimo 3 frases):
  \begin{center}
  \underline{\hspace{12cm}}\\\underline{\hspace{12cm}}\\\underline{\hspace{12cm}}\\\underline{\hspace{12cm}}\\\underline{\hspace{12cm}}
  \end{center}

  \item Releia e corrija: marque palavras com a grafia S/Z que você usou.

  \item Ilustre o texto:
  \begin{center}
  \begin{tikzpicture}
    \draw[thick, dashed] (0,0) rectangle (12,6);
  \end{tikzpicture}
  \end{center}

\end{enumerate}

\end{exercicio}

%% PLANCHETA 8.4

\plancheta{Folha de Prática 8.4 --- Ditado e Avaliação da Unidade 8}

\begin{exercicio}[Folha de Prática 8.4 --- Ditado e Avaliação da Unidade 8]

\noindent\textbf{Nome:} \underline{\hspace{3.2cm}} \quad \textbf{Data:} \underline{\hspace{1.6cm}} \quad \textbf{Nota:} \underline{\hspace{0.9cm}}/20


\subsection*{PARTE 1: DITADO (10 itens)}

O professor dita as palavras e frases abaixo; o aluno escreve com letra cursiva.

\begin{enumerate}

  \item \underline{\hspace{10.5cm}}

  \item \underline{\hspace{10.5cm}}

  \item \underline{\hspace{10.5cm}}

  \item \underline{\hspace{10.5cm}}

  \item \underline{\hspace{10.5cm}}

  \item \underline{\hspace{10.5cm}}

\end{enumerate}

\subsection*{PARTE 2: AUTOAVALIAÇÃO DO ALUNO}

\begin{enumerate}[label=\alph*)]

  \item Eu li as palavras da unidade sem ajuda.  \quad ( ) Sim   ( ) Não

  \item Eu escrevi o ditado com as grafias certas.  \quad ( ) Sim   ( ) Não

  \item Eu separei as palavras em sílabas.  \quad ( ) Sim   ( ) Não

  \item Eu sei explicar a regra com as minhas palavras.  \quad ( ) Sim   ( ) Não

\end{enumerate}

\subsection*{PARTE 3: REGISTRO DO PROFESSOR (S/F/R/N)}

\begin{center}
\renewcommand{\arraystretch}{1.4}
\begin{tabular}{@{}p{9cm}cccc@{}}
\toprule
Aspecto observado & S & F & R & N \\
\midrule
Lê com precisão as palavras-alvo & & & & \\
Escreve com ortografia estável & & & & \\
Generaliza a regra em escrita espontânea & & & & \\
Mantém atenção durante o ditado & & & & \\
\bottomrule\end{tabular}\end{center}

\end{exercicio}

\begin{dica}
**Dica para o Professor:** A regra 'S entre vogais soa como Z' é a mais frutífera do 2º ano; trabalhe em duplas de palavras (casa/caça).
\end{dica}

%% ============================================
%% UNIDADE 9 -- Dígrafos na escrita
%% ============================================

\chapter{Unidade 9 --- Dígrafos na escrita}

\metadata{I-09}{Silábico-Alfabético Avançado}{EF04LP01, EF04LP02, EF04LP03, EF04LP05, EF04LP06, EF04LP07, EF04LP09}{D1}{EF04LP01, EF04LP02, EF04LP03, EF04LP05, EF04LP06, EF04LP07, EF04LP09}{Resolucao CNE/CEB 7/2010, art. 12}

\textbf{Regra:} Ç representa o som de S e aparece antes de A, O, U: coração, cabeça, açúcar. Nunca antes de E ou I.

\textbf{Palavra geradora do bloco:} \textbf{CORAÇÃO} --- o coração é a palavra mais afetiva do mundo infantil; fixa ÇA, ÇO e ÇU com til.

\section{Lição 9.1 --- Consolidar a escrita de dígrafos}

\textbf{Foco da lição:} Consolidar a escrita de dígrafos


\begin{boquinha}[Boquinha da Técnica]
O cê-cedilha tem uma perninha que muda o som de K para S. Compare CA (K) com ÇA (S): a boca muda o gesto? O som muda.
\end{boquinha}

\noindent\textbf{Formação guiada:}

\begin{center}
\resizebox{\textwidth}{!}{%
\begin{tabular}{ccc}
  \textbf{ÇA} & \textbf{ÇO} & \textbf{ÇU} \\
  \end{tabular}}
\end{center}

\noindent\textbf{Palavras da unidade:} coração, caçarola, açúcar, poço, peça, laço, cabeça, maçã.

\noindent\textbf{Frases para leitura oral:}
\begin{itemize}[leftmargin=1.6em]
  \item O coração bate forte.
  \item O açúcar adoça o café.
  \item A cabeça dói um pouco.
  \item A menina ganhou um laço.
\end{itemize}

\section{Lição 9.2 --- Leitura do texto curto}

\noindent\textbf{Leia o texto em voz alta, com atenção às palavras da unidade:}


\begin{quote}
O coração bate forte. O açúcar adoça o café.

A cabeça dói um pouco. A menina ganhou um laço.
\end{quote}

\textbf{Depois de ler, responda por escrito:}
\begin{enumerate}

  \item O que acontece no texto que você leu?
  \begin{center}
  \underline{\hspace{12cm}}
  \end{center}

  \item Quantas frases o texto tem?
  \begin{center}
  \underline{\hspace{12cm}}
  \end{center}

  \item Circule no texto as palavras com a grafia Ç.

  \item Leia o texto de novo em voz alta, agora sem ajuda.

\end{enumerate}

\section{Lição 9.3 --- Escrita com a grafia Ç}

\noindent\textbf{Regra da unidade:} Ç representa o som de S e aparece antes de A, O, U: coração, cabeça, açúcar. Nunca antes de E ou I.


\noindent\textbf{Ditado guiado:} o professor dita as palavras abaixo; o aluno escreve e depois corrige com o professor.

\begin{center}
  \underline{\hspace{2.3cm}} \quad \underline{\hspace{2.3cm}} \quad \underline{\hspace{2.3cm}} \quad \underline{\hspace{2.3cm}}

  \underline{\hspace{2.3cm}} \quad \underline{\hspace{2.3cm}} \quad \underline{\hspace{2.3cm}} \quad \underline{\hspace{2.3cm}}
\end{center}

\noindent\textbf{Complete as palavras com a grafia Ç:}
\begin{center}
  coraçã\underline{~} \quad caçarol\underline{~} \quad açúca\underline{~} \quad poç\underline{~} \quad peç\underline{~}
\end{center}

\noindent\textbf{Escreva a palavra que o professor mostrar (banco de palavras) e separe em sílabas:}

\begin{center}
  \underline{\hspace{2cm}} $=$ \underline{\hspace{2cm}} \quad \underline{\hspace{2cm}} $=$ \underline{\hspace{2cm}} \quad \underline{\hspace{2cm}} $=$ \underline{\hspace{2cm}}
  \end{center}

%% PLANCHETA 9.1

\plancheta{Folha de Prática 9.1 --- Dígrafos na escrita (Parte 1)}

\begin{exercicio}[Folha de Prática 9.1 --- Dígrafos na escrita (Parte 1)]

\noindent\textbf{Nome:} \underline{\hspace{3.2cm}} \quad \textbf{Data:} \underline{\hspace{1.6cm}} \quad \textbf{Nota:} \underline{\hspace{0.9cm}}/10


\subsection*{PARTE 1: RECONHECIMENTO}

\begin{enumerate}

  \item Forme as sílabas com os sons da unidade:

\begin{center}
\resizebox{\textwidth}{!}{%
\begin{tabular}{ccc}
  \textbf{ÇA} & \textbf{ÇO} & \textbf{ÇU} \\
  \end{tabular}}
\end{center}

  \item Complete a tabela de sílabas da unidade:

\begin{center}
\resizebox{\textwidth}{!}{%
\begin{tabular}{ccccc}
  \sil{Ç}{A} & \sil{Ç}{E} & \sil{Ç}{I} & \sil{Ç}{O} & \sil{Ç}{U} \\
  \end{tabular}}
\end{center}

  \item Sublinhe as palavras da unidade:

\begin{center}
  \uline{coração} \quad \uline{caçarola} \quad \uline{açúcar} \quad \uline{poço} \quad \uline{peça} \quad \uline{laço} \quad cabeça \quad maçã \quad gato \quad sapo \quad rato \quad bola \quad menino
\end{center}

  \item Identifique o som-alvo em cada palavra:
  \begin{center}
  CORAÇÃO $\rightarrow$ som-alvo: \underline{\hspace{1cm}} \quad CAÇAROLA $\rightarrow$ som-alvo: \underline{\hspace{1cm}} \quad AÇÚCAR $\rightarrow$ som-alvo: \underline{\hspace{1cm}}
  \end{center}

  \item Separe em sílabas:
  \begin{center}
  CORAÇÃO $\rightarrow$ \underline{\hspace{1cm}} \quad CAÇAROLA $\rightarrow$ \underline{\hspace{1cm}} \quad AÇÚCAR $\rightarrow$ \underline{\hspace{1cm}}
  \end{center}

  \item Conte as sílabas:
  \begin{center}
  CORAÇÃO $\rightarrow$ \underline{\hspace{0.9cm}} \quad CAÇAROLA $\rightarrow$ \underline{\hspace{0.9cm}} \quad AÇÚCAR $\rightarrow$ \underline{\hspace{0.9cm}}
  \end{center}

\end{enumerate}

\subsection*{PARTE 2: PRODUÇÃO GUIADA}

\begin{enumerate}[resume]

  \item Complete com a grafia da unidade:
  \begin{center}
  Ç\underline{~~~} \quad Ç\underline{~~~} \quad Ç\underline{~~~} \quad Ç\underline{~~~} \quad Ç\underline{~~~}
  \end{center}

  \item Escreva 3 palavras com o som-alvo:
  \begin{center}
  1. \underline{\hspace{3cm}} \quad 2. \underline{\hspace{3cm}} \quad 3. \underline{\hspace{3cm}}
  \end{center}

  \item Copie as palavras:
  \begin{center}
  coração \quad caçarola \quad açúcar \quad poço \quad peça
  \end{center}

  \item Separe em sílabas:
  \begin{center}
  CORAÇÃO = \underline{\hspace{1cm}} \underline{\hspace{1cm}} \quad CAÇAROLA = \underline{\hspace{1cm}} \underline{\hspace{1cm}} \quad AÇÚCAR = \underline{\hspace{1cm}} \underline{\hspace{1cm}}
  \end{center}

  \item Escreva 2 frases usando palavras da unidade:
  \begin{enumerate}[label=\alph*)]
    \item \underline{\hspace{8cm}}
    \item \underline{\hspace{8cm}}
  \end{enumerate}

\end{enumerate}

\end{exercicio}

%% PLANCHETA 9.2

\plancheta{Folha de Prática 9.2 --- Dígrafos na escrita (Parte 2)}

\begin{exercicio}[Folha de Prática 9.2 --- Dígrafos na escrita (Parte 2)]

\noindent\textbf{Nome:} \underline{\hspace{3.2cm}} \quad \textbf{Data:} \underline{\hspace{1.6cm}} \quad \textbf{Nota:} \underline{\hspace{0.9cm}}/10


\subsection*{PARTE 3: INTEGRAÇÃO}

\begin{enumerate}[resume]

  \item Leia as palavras em voz alta:
  \begin{center}
  $\rightarrow$ coração ~ caçarola ~ açúcar ~ poço ~ peça ~ laço
  \end{center}

  \item Leia as frases em voz alta:
  \begin{center}
  $\rightarrow$ O coração bate forte.\\
  $\rightarrow$ O açúcar adoça o café.\\
  $\rightarrow$ A cabeça dói um pouco.\\
  $\rightarrow$ A menina ganhou um laço.\\
  \end{center}

  \item Circule a opção correta:
  \begin{enumerate}[label=(\alph*)]

    \item O som da unidade é: ( ) oral ( ) nasal 

    \item A primeira sílaba de CORAÇÃO é: ( ) ÇA ( ) ÇO

    \item As palavras da unidade aparecem: ( ) no início ( ) no meio ( ) no fim das palavras

  \end{enumerate}

  \item Leia e complete:
  \begin{center}
  cor\underline{~~~} \quad caç\underline{~~~} \quad açú\underline{~~~} \quad poç\underline{~~~}
  \end{center}

  \item Leitura em voz alta com o professor:
  \begin{center}
  coração \quad caçarola \quad açúcar \quad poço \quad peça \quad laço \quad cabeça \quad maçã
  \end{center}

\end{enumerate}

\end{exercicio}

%% PLANCHETA 9.3

\plancheta{Folha de Prática 9.3 --- Produção com Ç}

\begin{exercicio}[Folha de Prática 9.3 --- Produção com Ç]

\noindent\textbf{Nome:} \underline{\hspace{3.2cm}} \quad \textbf{Data:} \underline{\hspace{1.6cm}} \quad \textbf{Nota:} \underline{\hspace{0.9cm}}/15


\noindent\textbf{Comando:} Escrever três objetos de cozinha que tenham Ç (panelas, caçarola...) e usá-los em frases.


\begin{enumerate}

  \item Planeje: escolha 3 palavras da unidade para usar:
  \begin{center}
  1. \underline{\hspace{3cm}} \quad 2. \underline{\hspace{3cm}} \quad 3. \underline{\hspace{3cm}}
  \end{center}

  \item Escreva o texto (mínimo 3 frases):
  \begin{center}
  \underline{\hspace{12cm}}\\\underline{\hspace{12cm}}\\\underline{\hspace{12cm}}\\\underline{\hspace{12cm}}\\\underline{\hspace{12cm}}
  \end{center}

  \item Releia e corrija: marque palavras com a grafia Ç que você usou.

  \item Ilustre o texto:
  \begin{center}
  \begin{tikzpicture}
    \draw[thick, dashed] (0,0) rectangle (12,6);
  \end{tikzpicture}
  \end{center}

\end{enumerate}

\end{exercicio}

%% PLANCHETA 9.4

\plancheta{Folha de Prática 9.4 --- Ditado e Avaliação da Unidade 9}

\begin{exercicio}[Folha de Prática 9.4 --- Ditado e Avaliação da Unidade 9]

\noindent\textbf{Nome:} \underline{\hspace{3.2cm}} \quad \textbf{Data:} \underline{\hspace{1.6cm}} \quad \textbf{Nota:} \underline{\hspace{0.9cm}}/20


\subsection*{PARTE 1: DITADO (10 itens)}

O professor dita as palavras e frases abaixo; o aluno escreve com letra cursiva.

\begin{enumerate}

  \item \underline{\hspace{10.5cm}}

  \item \underline{\hspace{10.5cm}}

  \item \underline{\hspace{10.5cm}}

  \item \underline{\hspace{10.5cm}}

  \item \underline{\hspace{10.5cm}}

  \item \underline{\hspace{10.5cm}}

\end{enumerate}

\subsection*{PARTE 2: AUTOAVALIAÇÃO DO ALUNO}

\begin{enumerate}[label=\alph*)]

  \item Eu li as palavras da unidade sem ajuda.  \quad ( ) Sim   ( ) Não

  \item Eu escrevi o ditado com as grafias certas.  \quad ( ) Sim   ( ) Não

  \item Eu separei as palavras em sílabas.  \quad ( ) Sim   ( ) Não

  \item Eu sei explicar a regra com as minhas palavras.  \quad ( ) Sim   ( ) Não

\end{enumerate}

\subsection*{PARTE 3: REGISTRO DO PROFESSOR (S/F/R/N)}

\begin{center}
\renewcommand{\arraystretch}{1.4}
\begin{tabular}{@{}p{9cm}cccc@{}}
\toprule
Aspecto observado & S & F & R & N \\
\midrule
Lê com precisão as palavras-alvo & & & & \\
Escreve com ortografia estável & & & & \\
Generaliza a regra em escrita espontânea & & & & \\
Mantém atenção durante o ditado & & & & \\
\bottomrule\end{tabular}\end{center}

\end{exercicio}

%% PLANCHETA 9.5

\plancheta{Folha de Prática 9.5 --- Leitura e Fluência da Unidade 9}

\begin{exercicio}[Folha de Prática 9.5 --- Leitura e Fluência da Unidade 9]

\noindent\textbf{Nome:} \underline{\hspace{3.2cm}} \quad \textbf{Data:} \underline{\hspace{1.6cm}} \quad \textbf{Nota:} \underline{\hspace{0.9cm}}/10


\subsection*{PARTE 1: LEITURA EM VOZ ALTA}

Leia o texto abaixo três vezes; a cada leitura, marque um círculo no cronômetro do professor.

\begin{quote}
O coração bate forte. O açúcar adoça o café.

A cabeça dói um pouco. A menina ganhou um laço.
\end{quote}

\begin{center}
1ª leitura: \underline{\hspace{1.5cm}} \quad 2ª leitura: \underline{\hspace{1.5cm}} \quad 3ª leitura: \underline{\hspace{1.5cm}}
\end{center}

\subsection*{PARTE 2: COMPREENSÃO}

\begin{enumerate}

  \item Quem são os personagens do texto?
  \begin{center}
  \underline{\hspace{12cm}}
  \end{center}

  \item Qual é a grafia-alvo que mais aparece no texto? Escreva três exemplos:
  \begin{center}
  \underline{\hspace{3.5cm}} \quad \underline{\hspace{3.5cm}} \quad \underline{\hspace{3.5cm}}
  \end{center}

  \item Você gostou do texto? Por quê?
  \begin{center}
  \underline{\hspace{12cm}}
  \end{center}

\end{enumerate}

\subsection*{PARTE 3: DESENHO}

Desenhe uma cena do texto:

\begin{center}
\begin{tikzpicture}
  \draw[thick, dashed] (0,0) rectangle (12,5);
\end{tikzpicture}
\end{center}

\end{exercicio}

\begin{dica}
**Dica para o Professor:** Ç é um excelente porta de entrada para a ortografia: a regra (A, O, U) é simples e o som é previsível.
\end{dica}

%% ============================================
%% UNIDADE 10 -- Til e acentuação
%% ============================================

\input{checkpoints/rastreio-u5}
\chapter{Unidade 10 --- Til e acentuação}

\metadata{I-10}{Silábico-Alfabético Avançado}{EF04LP01, EF04LP02, EF04LP03, EF04LP05, EF04LP06, EF04LP07, EF04LP09}{D1}{EF04LP01, EF04LP02, EF04LP03, EF04LP05, EF04LP06, EF04LP07, EF04LP09}{Resolucao CNE/CEB 7/2010, art. 12}

\textbf{Regra:} AM e AN representam sílabas nasais: sambar, também, campo, dança. No final de verbos, escreve-se AM: eles cantam.

\textbf{Palavra geradora do bloco:} \textbf{TAMBÉM} --- também é uma palavra comum e funcional; a sílaba final BÉM mostra o AM nasal de forma clara.

\section{Lição 10.1 --- Acentuar corretamente palavras com til}

\textbf{Foco da lição:} Acentuar corretamente palavras com til


\begin{boquinha}[Boquinha da Técnica]
Ao dizer AM, os lábios fecham no M e o ar sai pelo nariz; ao dizer AN, a língua toca atrás dos dentes e o ar sai pelo nariz.
\end{boquinha}

\noindent\textbf{Formação guiada:}

\begin{center}
\resizebox{\textwidth}{!}{%
\begin{tabular}{ccccc}
  \textbf{MA} & \textbf{MB} & \textbf{BA} & \textbf{AM} & \textbf{AN} \\
  \textbf{CAN} & \textbf{SAN} & \textbf{TAM} \\
  \end{tabular}}
\end{center}

\noindent\textbf{Palavras da unidade:} também, sambar, canto, sangue, banco, campo, dança, lanche.

\noindent\textbf{Frases para leitura oral:}
\begin{itemize}[leftmargin=1.6em]
  \item Eu também quero lanche.
  \item O canto da sala é azul.
  \item O banco está na praça.
  \item A dança começa já.
\end{itemize}

\section{Lição 10.2 --- Leitura do texto curto}

\noindent\textbf{Leia o texto em voz alta, com atenção às palavras da unidade:}


\begin{quote}
Eu também quero lanche. O canto da sala é azul.

O banco está na praça. A dança começa já.
\end{quote}

\textbf{Depois de ler, responda por escrito:}
\begin{enumerate}

  \item O que acontece no texto que você leu?
  \begin{center}
  \underline{\hspace{12cm}}
  \end{center}

  \item Quantas frases o texto tem?
  \begin{center}
  \underline{\hspace{12cm}}
  \end{center}

  \item Circule no texto as palavras com a grafia AM/AN.

  \item Leia o texto de novo em voz alta, agora sem ajuda.

\end{enumerate}

\section{Lição 10.3 --- Escrita com a grafia AM/AN}

\noindent\textbf{Regra da unidade:} AM e AN representam sílabas nasais: sambar, também, campo, dança. No final de verbos, escreve-se AM: eles cantam.


\noindent\textbf{Ditado guiado:} o professor dita as palavras abaixo; o aluno escreve e depois corrige com o professor.

\begin{center}
  \underline{\hspace{2.3cm}} \quad \underline{\hspace{2.3cm}} \quad \underline{\hspace{2.3cm}} \quad \underline{\hspace{2.3cm}}

  \underline{\hspace{2.3cm}} \quad \underline{\hspace{2.3cm}} \quad \underline{\hspace{2.3cm}} \quad \underline{\hspace{2.3cm}}
\end{center}

\noindent\textbf{Complete as palavras com a grafia AM/AN:}
\begin{center}
  també\underline{~} \quad samba\underline{~} \quad cant\underline{~} \quad sangu\underline{~} \quad banc\underline{~}
\end{center}

\noindent\textbf{Escreva a palavra que o professor mostrar (banco de palavras) e separe em sílabas:}

\begin{center}
  \underline{\hspace{2cm}} $=$ \underline{\hspace{2cm}} \quad \underline{\hspace{2cm}} $=$ \underline{\hspace{2cm}} \quad \underline{\hspace{2cm}} $=$ \underline{\hspace{2cm}}
  \end{center}

%% PLANCHETA 10.1

\plancheta{Folha de Prática 10.1 --- Til e acentuação (Parte 1)}

\begin{exercicio}[Folha de Prática 10.1 --- Til e acentuação (Parte 1)]

\noindent\textbf{Nome:} \underline{\hspace{3.2cm}} \quad \textbf{Data:} \underline{\hspace{1.6cm}} \quad \textbf{Nota:} \underline{\hspace{0.9cm}}/10


\subsection*{PARTE 1: RECONHECIMENTO}

\begin{enumerate}

  \item Forme as sílabas com os sons da unidade:

\begin{center}
\resizebox{\textwidth}{!}{%
\begin{tabular}{ccccc}
  \textbf{MA} & \textbf{MB} & \textbf{BA} & \textbf{AM} & \textbf{AN} \\
  \textbf{CAN} & \textbf{SAN} & \textbf{TAM} \\
  \end{tabular}}
\end{center}

  \item Complete a tabela de sílabas da unidade:

\begin{center}
\resizebox{\textwidth}{!}{%
\begin{tabular}{ccccc}
  \sil{AM}{A} & \sil{AM}{E} & \sil{AM}{I} & \sil{AM}{O} & \sil{AM}{U} \\
  \end{tabular}}
\end{center}

  \item Sublinhe as palavras da unidade:

\begin{center}
  \uline{também} \quad \uline{sambar} \quad \uline{canto} \quad \uline{sangue} \quad \uline{banco} \quad \uline{campo} \quad dança \quad lanche \quad gato \quad sapo \quad rato \quad bola \quad menino
\end{center}

  \item Identifique o som-alvo em cada palavra:
  \begin{center}
  TAMBÉM $\rightarrow$ som-alvo: \underline{\hspace{1cm}} \quad SAMBAR $\rightarrow$ som-alvo: \underline{\hspace{1cm}} \quad CANTO $\rightarrow$ som-alvo: \underline{\hspace{1cm}}
  \end{center}

  \item Separe em sílabas:
  \begin{center}
  TAMBÉM $\rightarrow$ \underline{\hspace{1cm}} \quad SAMBAR $\rightarrow$ \underline{\hspace{1cm}} \quad CANTO $\rightarrow$ \underline{\hspace{1cm}}
  \end{center}

  \item Conte as sílabas:
  \begin{center}
  TAMBÉM $\rightarrow$ \underline{\hspace{0.9cm}} \quad SAMBAR $\rightarrow$ \underline{\hspace{0.9cm}} \quad CANTO $\rightarrow$ \underline{\hspace{0.9cm}}
  \end{center}

\end{enumerate}

\subsection*{PARTE 2: PRODUÇÃO GUIADA}

\begin{enumerate}[resume]

  \item Complete com a grafia da unidade:
  \begin{center}
  AM/AN\underline{~~~} \quad AM/AN\underline{~~~} \quad AM/AN\underline{~~~} \quad AM/AN\underline{~~~} \quad AM/AN\underline{~~~}
  \end{center}

  \item Escreva 3 palavras com o som-alvo:
  \begin{center}
  1. \underline{\hspace{3cm}} \quad 2. \underline{\hspace{3cm}} \quad 3. \underline{\hspace{3cm}}
  \end{center}

  \item Copie as palavras:
  \begin{center}
  também \quad sambar \quad canto \quad sangue \quad banco
  \end{center}

  \item Separe em sílabas:
  \begin{center}
  TAMBÉM = \underline{\hspace{1cm}} \underline{\hspace{1cm}} \quad SAMBAR = \underline{\hspace{1cm}} \underline{\hspace{1cm}} \quad CANTO = \underline{\hspace{1cm}} \underline{\hspace{1cm}}
  \end{center}

  \item Escreva 2 frases usando palavras da unidade:
  \begin{enumerate}[label=\alph*)]
    \item \underline{\hspace{8cm}}
    \item \underline{\hspace{8cm}}
  \end{enumerate}

\end{enumerate}

\end{exercicio}

%% PLANCHETA 10.2

\plancheta{Folha de Prática 10.2 --- Til e acentuação (Parte 2)}

\begin{exercicio}[Folha de Prática 10.2 --- Til e acentuação (Parte 2)]

\noindent\textbf{Nome:} \underline{\hspace{3.2cm}} \quad \textbf{Data:} \underline{\hspace{1.6cm}} \quad \textbf{Nota:} \underline{\hspace{0.9cm}}/10


\subsection*{PARTE 3: INTEGRAÇÃO}

\begin{enumerate}[resume]

  \item Leia as palavras em voz alta:
  \begin{center}
  $\rightarrow$ também ~ sambar ~ canto ~ sangue ~ banco ~ campo
  \end{center}

  \item Leia as frases em voz alta:
  \begin{center}
  $\rightarrow$ Eu também quero lanche.\\
  $\rightarrow$ O canto da sala é azul.\\
  $\rightarrow$ O banco está na praça.\\
  $\rightarrow$ A dança começa já.\\
  \end{center}

  \item Circule a opção correta:
  \begin{enumerate}[label=(\alph*)]

    \item O som da unidade é: ( ) oral ( ) nasal 

    \item A primeira sílaba de TAMBÉM é: ( ) AM ( ) AN

    \item As palavras da unidade aparecem: ( ) no início ( ) no meio ( ) no fim das palavras

  \end{enumerate}

  \item Leia e complete:
  \begin{center}
  tam\underline{~~~} \quad sam\underline{~~~} \quad can\underline{~~~} \quad san\underline{~~~}
  \end{center}

  \item Leitura em voz alta com o professor:
  \begin{center}
  também \quad sambar \quad canto \quad sangue \quad banco \quad campo \quad dança \quad lanche
  \end{center}

\end{enumerate}

\end{exercicio}

%% PLANCHETA 10.3

\plancheta{Folha de Prática 10.3 --- Produção com AM/AN}

\begin{exercicio}[Folha de Prática 10.3 --- Produção com AM/AN]

\noindent\textbf{Nome:} \underline{\hspace{3.2cm}} \quad \textbf{Data:} \underline{\hspace{1.6cm}} \quad \textbf{Nota:} \underline{\hspace{0.9cm}}/15


\noindent\textbf{Comando:} Escrever verbos no presente que terminam em AM (cantam, brincam) e usá-los em frases.


\begin{enumerate}

  \item Planeje: escolha 3 palavras da unidade para usar:
  \begin{center}
  1. \underline{\hspace{3cm}} \quad 2. \underline{\hspace{3cm}} \quad 3. \underline{\hspace{3cm}}
  \end{center}

  \item Escreva o texto (mínimo 3 frases):
  \begin{center}
  \underline{\hspace{12cm}}\\\underline{\hspace{12cm}}\\\underline{\hspace{12cm}}\\\underline{\hspace{12cm}}\\\underline{\hspace{12cm}}
  \end{center}

  \item Releia e corrija: marque palavras com a grafia AM/AN que você usou.

  \item Ilustre o texto:
  \begin{center}
  \begin{tikzpicture}
    \draw[thick, dashed] (0,0) rectangle (12,6);
  \end{tikzpicture}
  \end{center}

\end{enumerate}

\end{exercicio}

%% PLANCHETA 10.4

\plancheta{Folha de Prática 10.4 --- Ditado e Avaliação da Unidade 10}

\begin{exercicio}[Folha de Prática 10.4 --- Ditado e Avaliação da Unidade 10]

\noindent\textbf{Nome:} \underline{\hspace{3.2cm}} \quad \textbf{Data:} \underline{\hspace{1.6cm}} \quad \textbf{Nota:} \underline{\hspace{0.9cm}}/20


\subsection*{PARTE 1: DITADO (10 itens)}

O professor dita as palavras e frases abaixo; o aluno escreve com letra cursiva.

\begin{enumerate}

  \item \underline{\hspace{10.5cm}}

  \item \underline{\hspace{10.5cm}}

  \item \underline{\hspace{10.5cm}}

  \item \underline{\hspace{10.5cm}}

  \item \underline{\hspace{10.5cm}}

  \item \underline{\hspace{10.5cm}}

\end{enumerate}

\subsection*{PARTE 2: AUTOAVALIAÇÃO DO ALUNO}

\begin{enumerate}[label=\alph*)]

  \item Eu li as palavras da unidade sem ajuda.  \quad ( ) Sim   ( ) Não

  \item Eu escrevi o ditado com as grafias certas.  \quad ( ) Sim   ( ) Não

  \item Eu separei as palavras em sílabas.  \quad ( ) Sim   ( ) Não

  \item Eu sei explicar a regra com as minhas palavras.  \quad ( ) Sim   ( ) Não

\end{enumerate}

\subsection*{PARTE 3: REGISTRO DO PROFESSOR (S/F/R/N)}

\begin{center}
\renewcommand{\arraystretch}{1.4}
\begin{tabular}{@{}p{9cm}cccc@{}}
\toprule
Aspecto observado & S & F & R & N \\
\midrule
Lê com precisão as palavras-alvo & & & & \\
Escreve com ortografia estável & & & & \\
Generaliza a regra em escrita espontânea & & & & \\
Mantém atenção durante o ditado & & & & \\
\bottomrule\end{tabular}\end{center}

\end{exercicio}

\begin{dica}
**Dica para o Professor:** AM/AN é a porta para a morfologia: 'eles cantam' com AM ensina que a escrita do verbo não depende do som apenas.
\end{dica}

%% ============================================
%% UNIDADE 11 -- Discurso direto
%% ============================================

\chapter{Unidade 11 --- Discurso direto}

\metadata{I-11}{Silábico-Alfabético Avançado}{EF04LP01, EF04LP02, EF04LP03, EF04LP05, EF04LP06, EF04LP07, EF04LP09}{D1}{EF04LP01, EF04LP02, EF04LP03, EF04LP05, EF04LP06, EF04LP07, EF04LP09}{Resolucao CNE/CEB 7/2010, art. 12}

\textbf{Regra:} X tem quatro sons possíveis: CH (xadrez), Z (exame), KS (táxi) e S (texto). O contexto decide; a leitura oral ajuda.

\textbf{Palavra geradora do bloco:} \textbf{XADREZ} --- o xadrez é lúdico e permite explorar o X com som de CH em uma palavra de jogo.

\section{Lição 11.1 --- Pontuar a fala dos personagens (discurso direto)}

\textbf{Foco da lição:} Pontuar a fala dos personagens (discurso direto)


\begin{boquinha}[Boquinha da Técnica]
O X não tem boca própria: ele 'empresta' o som. Ao ler, experimente os quatro sons até a palavra fazer sentido.
\end{boquinha}

\noindent\textbf{Formação guiada:}

\begin{center}
\resizebox{\textwidth}{!}{%
\begin{tabular}{ccccc}
  \textbf{XE} & \textbf{XI} & \textbf{XA} & \textbf{XO} & \textbf{XU} \\
  \end{tabular}}
\end{center}

\noindent\textbf{Palavras da unidade:} xadrez, peixe, caixa, exame, táxi, xícara, enxada, luxo.

\noindent\textbf{Frases para leitura oral:}
\begin{itemize}[leftmargin=1.6em]
  \item O peixe nada no aquário.
  \item A caixa tem brinquedos.
  \item O táxi parou na esquina.
  \item A xícara está cheia.
\end{itemize}

\section{Lição 11.2 --- Leitura do texto curto}

\noindent\textbf{Leia o texto em voz alta, com atenção às palavras da unidade:}


\begin{quote}
O peixe nada no aquário. A caixa tem brinquedos.

O táxi parou na esquina. A xícara está cheia.
\end{quote}

\textbf{Depois de ler, responda por escrito:}
\begin{enumerate}

  \item O que acontece no texto que você leu?
  \begin{center}
  \underline{\hspace{12cm}}
  \end{center}

  \item Quantas frases o texto tem?
  \begin{center}
  \underline{\hspace{12cm}}
  \end{center}

  \item Circule no texto as palavras com a grafia X.

  \item Leia o texto de novo em voz alta, agora sem ajuda.

\end{enumerate}

\section{Lição 11.3 --- Escrita com a grafia X}

\noindent\textbf{Regra da unidade:} X tem quatro sons possíveis: CH (xadrez), Z (exame), KS (táxi) e S (texto). O contexto decide; a leitura oral ajuda.


\noindent\textbf{Ditado guiado:} o professor dita as palavras abaixo; o aluno escreve e depois corrige com o professor.

\begin{center}
  \underline{\hspace{2.3cm}} \quad \underline{\hspace{2.3cm}} \quad \underline{\hspace{2.3cm}} \quad \underline{\hspace{2.3cm}}

  \underline{\hspace{2.3cm}} \quad \underline{\hspace{2.3cm}} \quad \underline{\hspace{2.3cm}} \quad \underline{\hspace{2.3cm}}
\end{center}

\noindent\textbf{Complete as palavras com a grafia X:}
\begin{center}
  xadre\underline{~} \quad peix\underline{~} \quad caix\underline{~} \quad exam\underline{~} \quad táx\underline{~}
\end{center}

\noindent\textbf{Escreva a palavra que o professor mostrar (banco de palavras) e separe em sílabas:}

\begin{center}
  \underline{\hspace{2cm}} $=$ \underline{\hspace{2cm}} \quad \underline{\hspace{2cm}} $=$ \underline{\hspace{2cm}} \quad \underline{\hspace{2cm}} $=$ \underline{\hspace{2cm}}
  \end{center}

%% PLANCHETA 11.1

\plancheta{Folha de Prática 11.1 --- Discurso direto (Parte 1)}

\begin{exercicio}[Folha de Prática 11.1 --- Discurso direto (Parte 1)]

\noindent\textbf{Nome:} \underline{\hspace{3.2cm}} \quad \textbf{Data:} \underline{\hspace{1.6cm}} \quad \textbf{Nota:} \underline{\hspace{0.9cm}}/10


\subsection*{PARTE 1: RECONHECIMENTO}

\begin{enumerate}

  \item Forme as sílabas com os sons da unidade:

\begin{center}
\resizebox{\textwidth}{!}{%
\begin{tabular}{ccccc}
  \textbf{XE} & \textbf{XI} & \textbf{XA} & \textbf{XO} & \textbf{XU} \\
  \end{tabular}}
\end{center}

  \item Complete a tabela de sílabas da unidade:

\begin{center}
\resizebox{\textwidth}{!}{%
\begin{tabular}{ccccc}
  \sil{X}{A} & \sil{X}{E} & \sil{X}{I} & \sil{X}{O} & \sil{X}{U} \\
  \end{tabular}}
\end{center}

  \item Sublinhe as palavras da unidade:

\begin{center}
  \uline{xadrez} \quad \uline{peixe} \quad \uline{caixa} \quad \uline{exame} \quad \uline{táxi} \quad \uline{xícara} \quad enxada \quad luxo \quad gato \quad sapo \quad rato \quad bola \quad menino
\end{center}

  \item Identifique o som-alvo em cada palavra:
  \begin{center}
  XADREZ $\rightarrow$ som-alvo: \underline{\hspace{1cm}} \quad PEIXE $\rightarrow$ som-alvo: \underline{\hspace{1cm}} \quad CAIXA $\rightarrow$ som-alvo: \underline{\hspace{1cm}}
  \end{center}

  \item Separe em sílabas:
  \begin{center}
  XADREZ $\rightarrow$ \underline{\hspace{1cm}} \quad PEIXE $\rightarrow$ \underline{\hspace{1cm}} \quad CAIXA $\rightarrow$ \underline{\hspace{1cm}}
  \end{center}

  \item Conte as sílabas:
  \begin{center}
  XADREZ $\rightarrow$ \underline{\hspace{0.9cm}} \quad PEIXE $\rightarrow$ \underline{\hspace{0.9cm}} \quad CAIXA $\rightarrow$ \underline{\hspace{0.9cm}}
  \end{center}

\end{enumerate}

\subsection*{PARTE 2: PRODUÇÃO GUIADA}

\begin{enumerate}[resume]

  \item Complete com a grafia da unidade:
  \begin{center}
  X\underline{~~~} \quad X\underline{~~~} \quad X\underline{~~~} \quad X\underline{~~~} \quad X\underline{~~~}
  \end{center}

  \item Escreva 3 palavras com o som-alvo:
  \begin{center}
  1. \underline{\hspace{3cm}} \quad 2. \underline{\hspace{3cm}} \quad 3. \underline{\hspace{3cm}}
  \end{center}

  \item Copie as palavras:
  \begin{center}
  xadrez \quad peixe \quad caixa \quad exame \quad táxi
  \end{center}

  \item Separe em sílabas:
  \begin{center}
  XADREZ = \underline{\hspace{1cm}} \underline{\hspace{1cm}} \quad PEIXE = \underline{\hspace{1cm}} \underline{\hspace{1cm}} \quad CAIXA = \underline{\hspace{1cm}} \underline{\hspace{1cm}}
  \end{center}

  \item Escreva 2 frases usando palavras da unidade:
  \begin{enumerate}[label=\alph*)]
    \item \underline{\hspace{8cm}}
    \item \underline{\hspace{8cm}}
  \end{enumerate}

\end{enumerate}

\end{exercicio}

%% PLANCHETA 11.2

\plancheta{Folha de Prática 11.2 --- Discurso direto (Parte 2)}

\begin{exercicio}[Folha de Prática 11.2 --- Discurso direto (Parte 2)]

\noindent\textbf{Nome:} \underline{\hspace{3.2cm}} \quad \textbf{Data:} \underline{\hspace{1.6cm}} \quad \textbf{Nota:} \underline{\hspace{0.9cm}}/10


\subsection*{PARTE 3: INTEGRAÇÃO}

\begin{enumerate}[resume]

  \item Leia as palavras em voz alta:
  \begin{center}
  $\rightarrow$ xadrez ~ peixe ~ caixa ~ exame ~ táxi ~ xícara
  \end{center}

  \item Leia as frases em voz alta:
  \begin{center}
  $\rightarrow$ O peixe nada no aquário.\\
  $\rightarrow$ A caixa tem brinquedos.\\
  $\rightarrow$ O táxi parou na esquina.\\
  $\rightarrow$ A xícara está cheia.\\
  \end{center}

  \item Circule a opção correta:
  \begin{enumerate}[label=(\alph*)]

    \item O som da unidade é: ( ) oral ( ) nasal 

    \item A primeira sílaba de XADREZ é: ( ) XA ( ) XE

    \item As palavras da unidade aparecem: ( ) no início ( ) no meio ( ) no fim das palavras

  \end{enumerate}

  \item Leia e complete:
  \begin{center}
  xad\underline{~~~} \quad pei\underline{~~~} \quad cai\underline{~~~} \quad exa\underline{~~~}
  \end{center}

  \item Leitura em voz alta com o professor:
  \begin{center}
  xadrez \quad peixe \quad caixa \quad exame \quad táxi \quad xícara \quad enxada \quad luxo
  \end{center}

\end{enumerate}

\end{exercicio}

%% PLANCHETA 11.3

\plancheta{Folha de Prática 11.3 --- Produção com X}

\begin{exercicio}[Folha de Prática 11.3 --- Produção com X]

\noindent\textbf{Nome:} \underline{\hspace{3.2cm}} \quad \textbf{Data:} \underline{\hspace{1.6cm}} \quad \textbf{Nota:} \underline{\hspace{0.9cm}}/15


\noindent\textbf{Comando:} Escrever uma lista de palavras com X e marcar com qual som cada uma se pronuncia.


\begin{enumerate}

  \item Planeje: escolha 3 palavras da unidade para usar:
  \begin{center}
  1. \underline{\hspace{3cm}} \quad 2. \underline{\hspace{3cm}} \quad 3. \underline{\hspace{3cm}}
  \end{center}

  \item Escreva o texto (mínimo 3 frases):
  \begin{center}
  \underline{\hspace{12cm}}\\\underline{\hspace{12cm}}\\\underline{\hspace{12cm}}\\\underline{\hspace{12cm}}\\\underline{\hspace{12cm}}
  \end{center}

  \item Releia e corrija: marque palavras com a grafia X que você usou.

  \item Ilustre o texto:
  \begin{center}
  \begin{tikzpicture}
    \draw[thick, dashed] (0,0) rectangle (12,6);
  \end{tikzpicture}
  \end{center}

\end{enumerate}

\end{exercicio}

%% PLANCHETA 11.4

\plancheta{Folha de Prática 11.4 --- Ditado e Avaliação da Unidade 11}

\begin{exercicio}[Folha de Prática 11.4 --- Ditado e Avaliação da Unidade 11]

\noindent\textbf{Nome:} \underline{\hspace{3.2cm}} \quad \textbf{Data:} \underline{\hspace{1.6cm}} \quad \textbf{Nota:} \underline{\hspace{0.9cm}}/20


\subsection*{PARTE 1: DITADO (10 itens)}

O professor dita as palavras e frases abaixo; o aluno escreve com letra cursiva.

\begin{enumerate}

  \item \underline{\hspace{10.5cm}}

  \item \underline{\hspace{10.5cm}}

  \item \underline{\hspace{10.5cm}}

  \item \underline{\hspace{10.5cm}}

  \item \underline{\hspace{10.5cm}}

  \item \underline{\hspace{10.5cm}}

\end{enumerate}

\subsection*{PARTE 2: AUTOAVALIAÇÃO DO ALUNO}

\begin{enumerate}[label=\alph*)]

  \item Eu li as palavras da unidade sem ajuda.  \quad ( ) Sim   ( ) Não

  \item Eu escrevi o ditado com as grafias certas.  \quad ( ) Sim   ( ) Não

  \item Eu separei as palavras em sílabas.  \quad ( ) Sim   ( ) Não

  \item Eu sei explicar a regra com as minhas palavras.  \quad ( ) Sim   ( ) Não

\end{enumerate}

\subsection*{PARTE 3: REGISTRO DO PROFESSOR (S/F/R/N)}

\begin{center}
\renewcommand{\arraystretch}{1.4}
\begin{tabular}{@{}p{9cm}cccc@{}}
\toprule
Aspecto observado & S & F & R & N \\
\midrule
Lê com precisão as palavras-alvo & & & & \\
Escreve com ortografia estável & & & & \\
Generaliza a regra em escrita espontânea & & & & \\
Mantém atenção durante o ditado & & & & \\
\bottomrule\end{tabular}\end{center}

\end{exercicio}

\begin{dica}
**Dica para o Professor:** Não force a regra antes do tempo: a leitura por tentativa (e erro corrigido) é o caminho; o X é polissêmico.
\end{dica}

%% ============================================
%% UNIDADE 12 -- Discurso indireto
%% ============================================

\chapter{Unidade 12 --- Discurso indireto}

\metadata{I-12}{Silábico-Alfabético Avançado}{EF04LP01, EF04LP02, EF04LP03, EF04LP05, EF04LP06, EF04LP07, EF04LP09}{D1}{EF04LP01, EF04LP02, EF04LP03, EF04LP05, EF04LP06, EF04LP07, EF04LP09}{Resolucao CNE/CEB 7/2010, art. 12}

\textbf{Regra:} Nesta unidade, o aluno lê um texto curto com fluência, localiza informações e produz uma resposta escrita.

\textbf{Palavra geradora do bloco:} \textbf{BIBLIOTECA} --- a biblioteca reúne os aprendizados: ler com prazer é o objetivo final do 2º ano.

\section{Lição 12.1 --- Recontar a fala no discurso indireto}

\textbf{Foco da lição:} Recontar a fala no discurso indireto


\begin{boquinha}[Boquinha da Técnica]
A leitura fluente usa todas as Boquinhas do ano: o aluno percebe que já conhece os sons e lê com ritmo.
\end{boquinha}

\noindent\textbf{Formação guiada:}

\begin{center}
\resizebox{\textwidth}{!}{%
\begin{tabular}{cccc}
  \textbf{LEI} & \textbf{TEX} & \textbf{TO} & \textbf{FLU} \\
  \end{tabular}}
\end{center}

\noindent\textbf{Palavras da unidade:} biblioteca, leitura, estante, revista, história, capítulo, autor, página.

\noindent\textbf{Frases para leitura oral:}
\begin{itemize}[leftmargin=1.6em]
  \item A biblioteca da escola é grande.
  \item Eu leio um capítulo por dia.
  \item A história é de um coelho.
  \item O autor assinou o livro.
\end{itemize}

\section{Lição 12.2 --- Leitura do texto curto}

\noindent\textbf{Leia o texto em voz alta, com atenção às palavras da unidade:}


\begin{quote}
A biblioteca da escola é grande. Eu leio um capítulo por dia.

A história é de um coelho. O autor assinou o livro.
\end{quote}

\textbf{Depois de ler, responda por escrito:}
\begin{enumerate}

  \item O que acontece no texto que você leu?
  \begin{center}
  \underline{\hspace{12cm}}
  \end{center}

  \item Quantas frases o texto tem?
  \begin{center}
  \underline{\hspace{12cm}}
  \end{center}

  \item Circule no texto as palavras com a grafia LEITURA.

  \item Leia o texto de novo em voz alta, agora sem ajuda.

\end{enumerate}

\section{Lição 12.3 --- Escrita com a grafia LEITURA}

\noindent\textbf{Regra da unidade:} Nesta unidade, o aluno lê um texto curto com fluência, localiza informações e produz uma resposta escrita.


\noindent\textbf{Ditado guiado:} o professor dita as palavras abaixo; o aluno escreve e depois corrige com o professor.

\begin{center}
  \underline{\hspace{2.3cm}} \quad \underline{\hspace{2.3cm}} \quad \underline{\hspace{2.3cm}} \quad \underline{\hspace{2.3cm}}

  \underline{\hspace{2.3cm}} \quad \underline{\hspace{2.3cm}} \quad \underline{\hspace{2.3cm}} \quad \underline{\hspace{2.3cm}}
\end{center}

\noindent\textbf{Complete as palavras com a grafia LEITURA:}
\begin{center}
  bibliotec\underline{~} \quad leitur\underline{~} \quad estant\underline{~} \quad revist\underline{~} \quad históri\underline{~}
\end{center}

\noindent\textbf{Escreva a palavra que o professor mostrar (banco de palavras) e separe em sílabas:}

\begin{center}
  \underline{\hspace{2cm}} $=$ \underline{\hspace{2cm}} \quad \underline{\hspace{2cm}} $=$ \underline{\hspace{2cm}} \quad \underline{\hspace{2cm}} $=$ \underline{\hspace{2cm}}
  \end{center}

%% PLANCHETA 12.1

\plancheta{Folha de Prática 12.1 --- Discurso indireto (Parte 1)}

\begin{exercicio}[Folha de Prática 12.1 --- Discurso indireto (Parte 1)]

\noindent\textbf{Nome:} \underline{\hspace{3.2cm}} \quad \textbf{Data:} \underline{\hspace{1.6cm}} \quad \textbf{Nota:} \underline{\hspace{0.9cm}}/10


\subsection*{PARTE 1: RECONHECIMENTO}

\begin{enumerate}

  \item Forme as sílabas com os sons da unidade:

\begin{center}
\resizebox{\textwidth}{!}{%
\begin{tabular}{cccc}
  \textbf{LEI} & \textbf{TEX} & \textbf{TO} & \textbf{FLU} \\
  \end{tabular}}
\end{center}

  \item Sublinhe as palavras da unidade:

\begin{center}
  \uline{biblioteca} \quad \uline{leitura} \quad \uline{estante} \quad \uline{revista} \quad \uline{história} \quad \uline{capítulo} \quad autor \quad página \quad gato \quad sapo \quad rato \quad bola \quad menino
\end{center}

  \item Identifique o som-alvo em cada palavra:
  \begin{center}
  BIBLIOTECA $\rightarrow$ som-alvo: \underline{\hspace{1cm}} \quad LEITURA $\rightarrow$ som-alvo: \underline{\hspace{1cm}} \quad ESTANTE $\rightarrow$ som-alvo: \underline{\hspace{1cm}}
  \end{center}

  \item Separe em sílabas:
  \begin{center}
  BIBLIOTECA $\rightarrow$ \underline{\hspace{1cm}} \quad LEITURA $\rightarrow$ \underline{\hspace{1cm}} \quad ESTANTE $\rightarrow$ \underline{\hspace{1cm}}
  \end{center}

  \item Conte as sílabas:
  \begin{center}
  BIBLIOTECA $\rightarrow$ \underline{\hspace{0.9cm}} \quad LEITURA $\rightarrow$ \underline{\hspace{0.9cm}} \quad ESTANTE $\rightarrow$ \underline{\hspace{0.9cm}}
  \end{center}

\end{enumerate}

\subsection*{PARTE 2: PRODUÇÃO GUIADA}

\begin{enumerate}[resume]

  \item Complete com a grafia da unidade:
  \begin{center}
  LEITURA\underline{~~~} \quad LEITURA\underline{~~~} \quad LEITURA\underline{~~~} \quad LEITURA\underline{~~~} \quad LEITURA\underline{~~~}
  \end{center}

  \item Escreva 3 palavras com o som-alvo:
  \begin{center}
  1. \underline{\hspace{3cm}} \quad 2. \underline{\hspace{3cm}} \quad 3. \underline{\hspace{3cm}}
  \end{center}

  \item Copie as palavras:
  \begin{center}
  biblioteca \quad leitura \quad estante \quad revista \quad história
  \end{center}

  \item Separe em sílabas:
  \begin{center}
  BIBLIOTECA = \underline{\hspace{1cm}} \underline{\hspace{1cm}} \quad LEITURA = \underline{\hspace{1cm}} \underline{\hspace{1cm}} \quad ESTANTE = \underline{\hspace{1cm}} \underline{\hspace{1cm}}
  \end{center}

  \item Escreva 2 frases usando palavras da unidade:
  \begin{enumerate}[label=\alph*)]
    \item \underline{\hspace{8cm}}
    \item \underline{\hspace{8cm}}
  \end{enumerate}

\end{enumerate}

\end{exercicio}

%% PLANCHETA 12.2

\plancheta{Folha de Prática 12.2 --- Discurso indireto (Parte 2)}

\begin{exercicio}[Folha de Prática 12.2 --- Discurso indireto (Parte 2)]

\noindent\textbf{Nome:} \underline{\hspace{3.2cm}} \quad \textbf{Data:} \underline{\hspace{1.6cm}} \quad \textbf{Nota:} \underline{\hspace{0.9cm}}/10


\subsection*{PARTE 3: INTEGRAÇÃO}

\begin{enumerate}[resume]

  \item Leia as palavras em voz alta:
  \begin{center}
  $\rightarrow$ biblioteca ~ leitura ~ estante ~ revista ~ história ~ capítulo
  \end{center}

  \item Leia as frases em voz alta:
  \begin{center}
  $\rightarrow$ A biblioteca da escola é grande.\\
  $\rightarrow$ Eu leio um capítulo por dia.\\
  $\rightarrow$ A história é de um coelho.\\
  $\rightarrow$ O autor assinou o livro.\\
  \end{center}

  \item Circule a opção correta:
  \begin{enumerate}[label=(\alph*)]

    \item O som da unidade é: ( ) oral ( ) nasal ( ) vibrante

    \item A primeira sílaba de BIBLIOTECA é: ( ) LEI ( ) TEX

    \item As palavras da unidade aparecem: ( ) no início ( ) no meio ( ) no fim das palavras

  \end{enumerate}

  \item Leia e complete:
  \begin{center}
  bib\underline{~~~} \quad lei\underline{~~~} \quad est\underline{~~~} \quad rev\underline{~~~}
  \end{center}

  \item Leitura em voz alta com o professor:
  \begin{center}
  biblioteca \quad leitura \quad estante \quad revista \quad história \quad capítulo \quad autor \quad página
  \end{center}

\end{enumerate}

\end{exercicio}

%% PLANCHETA 12.3

\plancheta{Folha de Prática 12.3 --- Produção com LEITURA}

\begin{exercicio}[Folha de Prática 12.3 --- Produção com LEITURA]

\noindent\textbf{Nome:} \underline{\hspace{3.2cm}} \quad \textbf{Data:} \underline{\hspace{1.6cm}} \quad \textbf{Nota:} \underline{\hspace{0.9cm}}/15


\noindent\textbf{Comando:} Produzir um pequeno texto de opinião: qual foi o livro favorito do ano e por quê?


\begin{enumerate}

  \item Planeje: escolha 3 palavras da unidade para usar:
  \begin{center}
  1. \underline{\hspace{3cm}} \quad 2. \underline{\hspace{3cm}} \quad 3. \underline{\hspace{3cm}}
  \end{center}

  \item Escreva o texto (mínimo 3 frases):
  \begin{center}
  \underline{\hspace{12cm}}\\\underline{\hspace{12cm}}\\\underline{\hspace{12cm}}\\\underline{\hspace{12cm}}\\\underline{\hspace{12cm}}
  \end{center}

  \item Releia e corrija: marque palavras com a grafia LEITURA que você usou.

  \item Ilustre o texto:
  \begin{center}
  \begin{tikzpicture}
    \draw[thick, dashed] (0,0) rectangle (12,6);
  \end{tikzpicture}
  \end{center}

\end{enumerate}

\end{exercicio}

%% PLANCHETA 12.4

\plancheta{Folha de Prática 12.4 --- Ditado e Avaliação da Unidade 12}

\begin{exercicio}[Folha de Prática 12.4 --- Ditado e Avaliação da Unidade 12]

\noindent\textbf{Nome:} \underline{\hspace{3.2cm}} \quad \textbf{Data:} \underline{\hspace{1.6cm}} \quad \textbf{Nota:} \underline{\hspace{0.9cm}}/20


\subsection*{PARTE 1: DITADO (10 itens)}

O professor dita as palavras e frases abaixo; o aluno escreve com letra cursiva.

\begin{enumerate}

  \item \underline{\hspace{10.5cm}}

  \item \underline{\hspace{10.5cm}}

  \item \underline{\hspace{10.5cm}}

  \item \underline{\hspace{10.5cm}}

  \item \underline{\hspace{10.5cm}}

  \item \underline{\hspace{10.5cm}}

\end{enumerate}

\subsection*{PARTE 2: AUTOAVALIAÇÃO DO ALUNO}

\begin{enumerate}[label=\alph*)]

  \item Eu li as palavras da unidade sem ajuda.  \quad ( ) Sim   ( ) Não

  \item Eu escrevi o ditado com as grafias certas.  \quad ( ) Sim   ( ) Não

  \item Eu separei as palavras em sílabas.  \quad ( ) Sim   ( ) Não

  \item Eu sei explicar a regra com as minhas palavras.  \quad ( ) Sim   ( ) Não

\end{enumerate}

\subsection*{PARTE 3: REGISTRO DO PROFESSOR (S/F/R/N)}

\begin{center}
\renewcommand{\arraystretch}{1.4}
\begin{tabular}{@{}p{9cm}cccc@{}}
\toprule
Aspecto observado & S & F & R & N \\
\midrule
Lê com precisão as palavras-alvo & & & & \\
Escreve com ortografia estável & & & & \\
Generaliza a regra em escrita espontânea & & & & \\
Mantém atenção durante o ditado & & & & \\
\bottomrule\end{tabular}\end{center}

\end{exercicio}

%% PLANCHETA 12.5

\plancheta{Folha de Prática 12.5 --- Leitura e Fluência da Unidade 12}

\begin{exercicio}[Folha de Prática 12.5 --- Leitura e Fluência da Unidade 12]

\noindent\textbf{Nome:} \underline{\hspace{3.2cm}} \quad \textbf{Data:} \underline{\hspace{1.6cm}} \quad \textbf{Nota:} \underline{\hspace{0.9cm}}/10


\subsection*{PARTE 1: LEITURA EM VOZ ALTA}

Leia o texto abaixo três vezes; a cada leitura, marque um círculo no cronômetro do professor.

\begin{quote}
A biblioteca da escola é grande. Eu leio um capítulo por dia.

A história é de um coelho. O autor assinou o livro.
\end{quote}

\begin{center}
1ª leitura: \underline{\hspace{1.5cm}} \quad 2ª leitura: \underline{\hspace{1.5cm}} \quad 3ª leitura: \underline{\hspace{1.5cm}}
\end{center}

\subsection*{PARTE 2: COMPREENSÃO}

\begin{enumerate}

  \item Quem são os personagens do texto?
  \begin{center}
  \underline{\hspace{12cm}}
  \end{center}

  \item Qual é a grafia-alvo que mais aparece no texto? Escreva três exemplos:
  \begin{center}
  \underline{\hspace{3.5cm}} \quad \underline{\hspace{3.5cm}} \quad \underline{\hspace{3.5cm}}
  \end{center}

  \item Você gostou do texto? Por quê?
  \begin{center}
  \underline{\hspace{12cm}}
  \end{center}

\end{enumerate}

\subsection*{PARTE 3: DESENHO}

Desenhe uma cena do texto:

\begin{center}
\begin{tikzpicture}
  \draw[thick, dashed] (0,0) rectangle (12,5);
\end{tikzpicture}
\end{center}

\end{exercicio}

\begin{dica}
**Dica para o Professor:** Use a fluência como meta: rum, não velocidade. A entonação correta indica compreensão; registre no Registro de Progresso.
\end{dica}

%% ============================================
%% UNIDADE 13 -- Rima e poema
%% ============================================

\chapter{Unidade 13 --- Rima e poema}

\metadata{I-13}{Silábico-Alfabético Avançado}{EF04LP01, EF04LP02, EF04LP03, EF04LP05, EF04LP06, EF04LP07, EF04LP09}{D1}{EF04LP01, EF04LP02, EF04LP03, EF04LP05, EF04LP06, EF04LP07, EF04LP09}{Resolucao CNE/CEB 7/2010, art. 12}

\textbf{Regra:} CH, LH e NH são dígrafos: duas letras, um som. A revisão lembra a boca de cada um e fortalece a leitura e a escrita.

\textbf{Palavra geradora do bloco:} \textbf{COELHO} --- o coelho junta LH e o afeto; revisitá-lo amarra a revisão dos três dígrafos.

\section{Lição 13.1 --- Ler e produzir poemas com rima}

\textbf{Foco da lição:} Ler e produzir poemas com rima


\begin{boquinha}[Boquinha da Técnica]
Na revisão, compare as três Boquinhas: CH tem atrito de ar, LH tem língua lateral, NH tem ar pelo nariz. Diga as três palavras-chave: CHUVA, COELHO, NINHO.
\end{boquinha}

\noindent\textbf{Formação guiada:}

\begin{center}
\resizebox{\textwidth}{!}{%
\begin{tabular}{ccccc}
  \sil{CH}{A} & \sil{CH}{O} & \sil{LH}{A} & \sil{LH}{O} & \sil{NH}{A} \\
  \sil{NH}{O} & \sil{CH}{E} & \sil{NH}{E} \\
  \end{tabular}}
\end{center}

\noindent\textbf{Palavras da unidade:} chuva, coelho, ninho, chave, galinha, folha, sonho, chapéu.

\noindent\textbf{Frases para leitura oral:}
\begin{itemize}[leftmargin=1.6em]
  \item A chuva molhou o ninho.
  \item O coelho escondeu-se na folha.
  \item A galinha sonha com o milho.
  \item O chapéu do velho é preto.
\end{itemize}

\section{Lição 13.2 --- Leitura do texto curto}

\noindent\textbf{Leia o texto em voz alta, com atenção às palavras da unidade:}


\begin{quote}
A chuva molhou o ninho. O coelho escondeu-se na folha.

A galinha sonha com o milho. O chapéu do velho é preto.
\end{quote}

\textbf{Depois de ler, responda por escrito:}
\begin{enumerate}

  \item O que acontece no texto que você leu?
  \begin{center}
  \underline{\hspace{12cm}}
  \end{center}

  \item Quantas frases o texto tem?
  \begin{center}
  \underline{\hspace{12cm}}
  \end{center}

  \item Circule no texto as palavras com a grafia CH/LH/NH.

  \item Leia o texto de novo em voz alta, agora sem ajuda.

\end{enumerate}

\section{Lição 13.3 --- Escrita com a grafia CH/LH/NH}

\noindent\textbf{Regra da unidade:} CH, LH e NH são dígrafos: duas letras, um som. A revisão lembra a boca de cada um e fortalece a leitura e a escrita.


\noindent\textbf{Ditado guiado:} o professor dita as palavras abaixo; o aluno escreve e depois corrige com o professor.

\begin{center}
  \underline{\hspace{2.3cm}} \quad \underline{\hspace{2.3cm}} \quad \underline{\hspace{2.3cm}} \quad \underline{\hspace{2.3cm}}

  \underline{\hspace{2.3cm}} \quad \underline{\hspace{2.3cm}} \quad \underline{\hspace{2.3cm}} \quad \underline{\hspace{2.3cm}}
\end{center}

\noindent\textbf{Complete as palavras com a grafia CH/LH/NH:}
\begin{center}
  \underline{~~~}uva \quad coelh\underline{~} \quad ninh\underline{~} \quad \underline{~~~}ave \quad galinh\underline{~}
\end{center}

\noindent\textbf{Escreva a palavra que o professor mostrar (banco de palavras) e separe em sílabas:}

\begin{center}
  \underline{\hspace{2cm}} $=$ \underline{\hspace{2cm}} \quad \underline{\hspace{2cm}} $=$ \underline{\hspace{2cm}} \quad \underline{\hspace{2cm}} $=$ \underline{\hspace{2cm}}
  \end{center}

%% PLANCHETA 13.1

\plancheta{Folha de Prática 13.1 --- Rima e poema (Parte 1)}

\begin{exercicio}[Folha de Prática 13.1 --- Rima e poema (Parte 1)]

\noindent\textbf{Nome:} \underline{\hspace{3.2cm}} \quad \textbf{Data:} \underline{\hspace{1.6cm}} \quad \textbf{Nota:} \underline{\hspace{0.9cm}}/10


\subsection*{PARTE 1: RECONHECIMENTO}

\begin{enumerate}

  \item Forme as sílabas com os sons da unidade:

\begin{center}
\resizebox{\textwidth}{!}{%
\begin{tabular}{ccccc}
  \sil{CH}{A} & \sil{CH}{O} & \sil{LH}{A} & \sil{LH}{O} & \sil{NH}{A} \\
  \sil{NH}{O} & \sil{CH}{E} & \sil{NH}{E} \\
  \end{tabular}}
\end{center}

  \item Complete a tabela de sílabas da unidade:

\begin{center}
\resizebox{\textwidth}{!}{%
\begin{tabular}{ccccc}
  \sil{CH}{A} & \sil{CH}{E} & \sil{CH}{I} & \sil{CH}{O} & \sil{CH}{U} \\
  \end{tabular}}
\end{center}

  \item Sublinhe as palavras da unidade:

\begin{center}
  \uline{chuva} \quad \uline{coelho} \quad \uline{ninho} \quad \uline{chave} \quad \uline{galinha} \quad \uline{folha} \quad sonho \quad chapéu \quad gato \quad sapo \quad rato \quad bola \quad menino
\end{center}

  \item Identifique o som-alvo em cada palavra:
  \begin{center}
  CHUVA $\rightarrow$ som-alvo: \underline{\hspace{1cm}} \quad COELHO $\rightarrow$ som-alvo: \underline{\hspace{1cm}} \quad NINHO $\rightarrow$ som-alvo: \underline{\hspace{1cm}}
  \end{center}

  \item Separe em sílabas:
  \begin{center}
  CHUVA $\rightarrow$ \underline{\hspace{1cm}} \quad COELHO $\rightarrow$ \underline{\hspace{1cm}} \quad NINHO $\rightarrow$ \underline{\hspace{1cm}}
  \end{center}

  \item Conte as sílabas:
  \begin{center}
  CHUVA $\rightarrow$ \underline{\hspace{0.9cm}} \quad COELHO $\rightarrow$ \underline{\hspace{0.9cm}} \quad NINHO $\rightarrow$ \underline{\hspace{0.9cm}}
  \end{center}

\end{enumerate}

\subsection*{PARTE 2: PRODUÇÃO GUIADA}

\begin{enumerate}[resume]

  \item Complete com a grafia da unidade:
  \begin{center}
  CH/LH/NH\underline{~~~} \quad CH/LH/NH\underline{~~~} \quad CH/LH/NH\underline{~~~} \quad CH/LH/NH\underline{~~~} \quad CH/LH/NH\underline{~~~}
  \end{center}

  \item Escreva 3 palavras com o som-alvo:
  \begin{center}
  1. \underline{\hspace{3cm}} \quad 2. \underline{\hspace{3cm}} \quad 3. \underline{\hspace{3cm}}
  \end{center}

  \item Copie as palavras:
  \begin{center}
  chuva \quad coelho \quad ninho \quad chave \quad galinha
  \end{center}

  \item Separe em sílabas:
  \begin{center}
  CHUVA = \underline{\hspace{1cm}} \underline{\hspace{1cm}} \quad COELHO = \underline{\hspace{1cm}} \underline{\hspace{1cm}} \quad NINHO = \underline{\hspace{1cm}} \underline{\hspace{1cm}}
  \end{center}

  \item Escreva 2 frases usando palavras da unidade:
  \begin{enumerate}[label=\alph*)]
    \item \underline{\hspace{8cm}}
    \item \underline{\hspace{8cm}}
  \end{enumerate}

\end{enumerate}

\end{exercicio}

%% PLANCHETA 13.2

\plancheta{Folha de Prática 13.2 --- Rima e poema (Parte 2)}

\begin{exercicio}[Folha de Prática 13.2 --- Rima e poema (Parte 2)]

\noindent\textbf{Nome:} \underline{\hspace{3.2cm}} \quad \textbf{Data:} \underline{\hspace{1.6cm}} \quad \textbf{Nota:} \underline{\hspace{0.9cm}}/10


\subsection*{PARTE 3: INTEGRAÇÃO}

\begin{enumerate}[resume]

  \item Leia as palavras em voz alta:
  \begin{center}
  $\rightarrow$ chuva ~ coelho ~ ninho ~ chave ~ galinha ~ folha
  \end{center}

  \item Leia as frases em voz alta:
  \begin{center}
  $\rightarrow$ A chuva molhou o ninho.\\
  $\rightarrow$ O coelho escondeu-se na folha.\\
  $\rightarrow$ A galinha sonha com o milho.\\
  $\rightarrow$ O chapéu do velho é preto.\\
  \end{center}

  \item Circule a opção correta:
  \begin{enumerate}[label=(\alph*)]

    \item O som da unidade é: ( ) oral ( ) nasal 

    \item A primeira sílaba de CHUVA é: ( ) CHA ( ) LHA

    \item As palavras da unidade aparecem: ( ) no início ( ) no meio ( ) no fim das palavras

  \end{enumerate}

  \item Leia e complete:
  \begin{center}
  chu\underline{~~~} \quad coe\underline{~~~} \quad nin\underline{~~~} \quad cha\underline{~~~}
  \end{center}

  \item Leitura em voz alta com o professor:
  \begin{center}
  chuva \quad coelho \quad ninho \quad chave \quad galinha \quad folha \quad sonho \quad chapéu
  \end{center}

\end{enumerate}

\end{exercicio}

%% PLANCHETA 13.3

\plancheta{Folha de Prática 13.3 --- Produção com CH/LH/NH}

\begin{exercicio}[Folha de Prática 13.3 --- Produção com CH/LH/NH]

\noindent\textbf{Nome:} \underline{\hspace{3.2cm}} \quad \textbf{Data:} \underline{\hspace{1.6cm}} \quad \textbf{Nota:} \underline{\hspace{0.9cm}}/15


\noindent\textbf{Comando:} Escrever um pequeno conto de 4 frases usando pelo menos um CH, um LH e um NH.


\begin{enumerate}

  \item Planeje: escolha 3 palavras da unidade para usar:
  \begin{center}
  1. \underline{\hspace{3cm}} \quad 2. \underline{\hspace{3cm}} \quad 3. \underline{\hspace{3cm}}
  \end{center}

  \item Escreva o texto (mínimo 3 frases):
  \begin{center}
  \underline{\hspace{12cm}}\\\underline{\hspace{12cm}}\\\underline{\hspace{12cm}}\\\underline{\hspace{12cm}}\\\underline{\hspace{12cm}}
  \end{center}

  \item Releia e corrija: marque palavras com a grafia CH/LH/NH que você usou.

  \item Ilustre o texto:
  \begin{center}
  \begin{tikzpicture}
    \draw[thick, dashed] (0,0) rectangle (12,6);
  \end{tikzpicture}
  \end{center}

\end{enumerate}

\end{exercicio}

%% PLANCHETA 13.4

\plancheta{Folha de Prática 13.4 --- Ditado e Avaliação da Unidade 13}

\begin{exercicio}[Folha de Prática 13.4 --- Ditado e Avaliação da Unidade 13]

\noindent\textbf{Nome:} \underline{\hspace{3.2cm}} \quad \textbf{Data:} \underline{\hspace{1.6cm}} \quad \textbf{Nota:} \underline{\hspace{0.9cm}}/20


\subsection*{PARTE 1: DITADO (10 itens)}

O professor dita as palavras e frases abaixo; o aluno escreve com letra cursiva.

\begin{enumerate}

  \item \underline{\hspace{10.5cm}}

  \item \underline{\hspace{10.5cm}}

  \item \underline{\hspace{10.5cm}}

  \item \underline{\hspace{10.5cm}}

  \item \underline{\hspace{10.5cm}}

  \item \underline{\hspace{10.5cm}}

\end{enumerate}

\subsection*{PARTE 2: AUTOAVALIAÇÃO DO ALUNO}

\begin{enumerate}[label=\alph*)]

  \item Eu li as palavras da unidade sem ajuda.  \quad ( ) Sim   ( ) Não

  \item Eu escrevi o ditado com as grafias certas.  \quad ( ) Sim   ( ) Não

  \item Eu separei as palavras em sílabas.  \quad ( ) Sim   ( ) Não

  \item Eu sei explicar a regra com as minhas palavras.  \quad ( ) Sim   ( ) Não

\end{enumerate}

\subsection*{PARTE 3: REGISTRO DO PROFESSOR (S/F/R/N)}

\begin{center}
\renewcommand{\arraystretch}{1.4}
\begin{tabular}{@{}p{9cm}cccc@{}}
\toprule
Aspecto observado & S & F & R & N \\
\midrule
Lê com precisão as palavras-alvo & & & & \\
Escreve com ortografia estável & & & & \\
Generaliza a regra em escrita espontânea & & & & \\
Mantém atenção durante o ditado & & & & \\
\bottomrule\end{tabular}\end{center}

\end{exercicio}

\begin{dica}
**Dica para o Professor:** Use o quadro de contraste das três Boquinhas: a revisão é o momento de separar definitivamente os três sons.
\end{dica}

%% ============================================
%% UNIDADE 14 -- Produção: a notícia
%% ============================================

\chapter{Unidade 14 --- Produção: a notícia}

\metadata{I-14}{Silábico-Alfabético Avançado}{EF04LP01, EF04LP02, EF04LP03, EF04LP05, EF04LP06, EF04LP07, EF04LP09}{D1}{EF04LP01, EF04LP02, EF04LP03, EF04LP05, EF04LP06, EF04LP07, EF04LP09}{Resolucao CNE/CEB 7/2010, art. 12}

\textbf{Regra:} Os encontros consonantais mantêm os dois sons juntos: BRA, CRA, DRA... e BLA, CLA, FLA. A revisão automatiza a leitura.

\textbf{Palavra geradora do bloco:} \textbf{TREM} --- o trem junta TR e movimento; é o encontro preferido das crianças.

\section{Lição 14.1 --- Produzir notícia com lead e fatos}

\textbf{Foco da lição:} Produzir notícia com lead e fatos


\begin{boquinha}[Boquinha da Técnica]
Na revisão, leia os encontros em cadeia: BRA BRE BRI BRO BRU. Depois os de L: BLA BLE BLI BLO BLU. Sinta a diferença entre o R e o L.
\end{boquinha}

\noindent\textbf{Formação guiada:}

\begin{center}
\resizebox{\textwidth}{!}{%
\begin{tabular}{ccccc}
  \sil{BR}{A} & \sil{PR}{A} & \sil{TR}{A} & \sil{GR}{A} & \sil{FR}{A} \\
  \sil{CR}{A} & \sil{DR}{A} & \sil{VR}{A} & \sil{BL}{A} & \sil{CL}{A} \\
  \sil{FL}{A} & \sil{GL}{A} & \sil{PL}{A} \\
  \end{tabular}}
\end{center}

\noindent\textbf{Palavras da unidade:} trem, prato, braço, grade, fruta, cravo, dragão, flor, placa, livro.

\noindent\textbf{Frases para leitura oral:}
\begin{itemize}[leftmargin=1.6em]
  \item O trem passou pela grade.
  \item O prato tem fruta e flor.
  \item O dragão voou sobre a placa.
  \item O livro tem um cravo na capa.
\end{itemize}

\section{Lição 14.2 --- Leitura do texto curto}

\noindent\textbf{Leia o texto em voz alta, com atenção às palavras da unidade:}


\begin{quote}
O trem passou pela grade. O prato tem fruta e flor.

O dragão voou sobre a placa. O livro tem um cravo na capa.
\end{quote}

\textbf{Depois de ler, responda por escrito:}
\begin{enumerate}

  \item O que acontece no texto que você leu?
  \begin{center}
  \underline{\hspace{12cm}}
  \end{center}

  \item Quantas frases o texto tem?
  \begin{center}
  \underline{\hspace{12cm}}
  \end{center}

  \item Circule no texto as palavras com a grafia ENCONTROS.

  \item Leia o texto de novo em voz alta, agora sem ajuda.

\end{enumerate}

\section{Lição 14.3 --- Escrita com a grafia ENCONTROS}

\noindent\textbf{Regra da unidade:} Os encontros consonantais mantêm os dois sons juntos: BRA, CRA, DRA... e BLA, CLA, FLA. A revisão automatiza a leitura.


\noindent\textbf{Ditado guiado:} o professor dita as palavras abaixo; o aluno escreve e depois corrige com o professor.

\begin{center}
  \underline{\hspace{2.3cm}} \quad \underline{\hspace{2.3cm}} \quad \underline{\hspace{2.3cm}} \quad \underline{\hspace{2.3cm}}

  \underline{\hspace{2.3cm}} \quad \underline{\hspace{2.3cm}} \quad \underline{\hspace{2.3cm}} \quad \underline{\hspace{2.3cm}}
\end{center}

\noindent\textbf{Complete as palavras com a grafia ENCONTROS:}
\begin{center}
  tre\underline{~} \quad prat\underline{~} \quad braç\underline{~} \quad grad\underline{~} \quad frut\underline{~}
\end{center}

\noindent\textbf{Escreva a palavra que o professor mostrar (banco de palavras) e separe em sílabas:}

\begin{center}
  \underline{\hspace{2cm}} $=$ \underline{\hspace{2cm}} \quad \underline{\hspace{2cm}} $=$ \underline{\hspace{2cm}} \quad \underline{\hspace{2cm}} $=$ \underline{\hspace{2cm}}
  \end{center}

%% PLANCHETA 14.1

\plancheta{Folha de Prática 14.1 --- Produção: a notícia (Parte 1)}

\begin{exercicio}[Folha de Prática 14.1 --- Produção: a notícia (Parte 1)]

\noindent\textbf{Nome:} \underline{\hspace{3.2cm}} \quad \textbf{Data:} \underline{\hspace{1.6cm}} \quad \textbf{Nota:} \underline{\hspace{0.9cm}}/10


\subsection*{PARTE 1: RECONHECIMENTO}

\begin{enumerate}

  \item Forme as sílabas com os sons da unidade:

\begin{center}
\resizebox{\textwidth}{!}{%
\begin{tabular}{ccccc}
  \sil{BR}{A} & \sil{PR}{A} & \sil{TR}{A} & \sil{GR}{A} & \sil{FR}{A} \\
  \sil{CR}{A} & \sil{DR}{A} & \sil{VR}{A} & \sil{BL}{A} & \sil{CL}{A} \\
  \sil{FL}{A} & \sil{GL}{A} & \sil{PL}{A} \\
  \end{tabular}}
\end{center}

  \item Sublinhe as palavras da unidade:

\begin{center}
  \uline{trem} \quad \uline{prato} \quad \uline{braço} \quad \uline{grade} \quad \uline{fruta} \quad \uline{cravo} \quad dragão \quad flor \quad placa \quad livro \quad gato \quad sapo \quad rato \quad bola \quad menino
\end{center}

  \item Identifique o som-alvo em cada palavra:
  \begin{center}
  TREM $\rightarrow$ som-alvo: \underline{\hspace{1cm}} \quad PRATO $\rightarrow$ som-alvo: \underline{\hspace{1cm}} \quad BRAÇO $\rightarrow$ som-alvo: \underline{\hspace{1cm}}
  \end{center}

  \item Separe em sílabas:
  \begin{center}
  TREM $\rightarrow$ \underline{\hspace{1cm}} \quad PRATO $\rightarrow$ \underline{\hspace{1cm}} \quad BRAÇO $\rightarrow$ \underline{\hspace{1cm}}
  \end{center}

  \item Conte as sílabas:
  \begin{center}
  TREM $\rightarrow$ \underline{\hspace{0.9cm}} \quad PRATO $\rightarrow$ \underline{\hspace{0.9cm}} \quad BRAÇO $\rightarrow$ \underline{\hspace{0.9cm}}
  \end{center}

\end{enumerate}

\subsection*{PARTE 2: PRODUÇÃO GUIADA}

\begin{enumerate}[resume]

  \item Complete com a grafia da unidade:
  \begin{center}
  ENCONTROS\underline{~~~} \quad ENCONTROS\underline{~~~} \quad ENCONTROS\underline{~~~} \quad ENCONTROS\underline{~~~} \quad ENCONTROS\underline{~~~}
  \end{center}

  \item Escreva 3 palavras com o som-alvo:
  \begin{center}
  1. \underline{\hspace{3cm}} \quad 2. \underline{\hspace{3cm}} \quad 3. \underline{\hspace{3cm}}
  \end{center}

  \item Copie as palavras:
  \begin{center}
  trem \quad prato \quad braço \quad grade \quad fruta
  \end{center}

  \item Separe em sílabas:
  \begin{center}
  TREM = \underline{\hspace{1cm}} \underline{\hspace{1cm}} \quad PRATO = \underline{\hspace{1cm}} \underline{\hspace{1cm}} \quad BRAÇO = \underline{\hspace{1cm}} \underline{\hspace{1cm}}
  \end{center}

  \item Escreva 2 frases usando palavras da unidade:
  \begin{enumerate}[label=\alph*)]
    \item \underline{\hspace{8cm}}
    \item \underline{\hspace{8cm}}
  \end{enumerate}

\end{enumerate}

\end{exercicio}

%% PLANCHETA 14.2

\plancheta{Folha de Prática 14.2 --- Produção: a notícia (Parte 2)}

\begin{exercicio}[Folha de Prática 14.2 --- Produção: a notícia (Parte 2)]

\noindent\textbf{Nome:} \underline{\hspace{3.2cm}} \quad \textbf{Data:} \underline{\hspace{1.6cm}} \quad \textbf{Nota:} \underline{\hspace{0.9cm}}/10


\subsection*{PARTE 3: INTEGRAÇÃO}

\begin{enumerate}[resume]

  \item Leia as palavras em voz alta:
  \begin{center}
  $\rightarrow$ trem ~ prato ~ braço ~ grade ~ fruta ~ cravo
  \end{center}

  \item Leia as frases em voz alta:
  \begin{center}
  $\rightarrow$ O trem passou pela grade.\\
  $\rightarrow$ O prato tem fruta e flor.\\
  $\rightarrow$ O dragão voou sobre a placa.\\
  $\rightarrow$ O livro tem um cravo na capa.\\
  \end{center}

  \item Circule a opção correta:
  \begin{enumerate}[label=(\alph*)]

    \item O som da unidade é: ( ) oral ( ) nasal ( ) vibrante

    \item A primeira sílaba de TREM é: ( ) TRA ( ) PLA

    \item As palavras da unidade aparecem: ( ) no início ( ) no meio ( ) no fim das palavras

  \end{enumerate}

  \item Leia e complete:
  \begin{center}
  tre\underline{~~~} \quad pra\underline{~~~} \quad bra\underline{~~~} \quad gra\underline{~~~}
  \end{center}

  \item Leitura em voz alta com o professor:
  \begin{center}
  trem \quad prato \quad braço \quad grade \quad fruta \quad cravo \quad dragão \quad flor
  \end{center}

\end{enumerate}

\end{exercicio}

%% PLANCHETA 14.3

\plancheta{Folha de Prática 14.3 --- Produção com ENCONTROS}

\begin{exercicio}[Folha de Prática 14.3 --- Produção com ENCONTROS]

\noindent\textbf{Nome:} \underline{\hspace{3.2cm}} \quad \textbf{Data:} \underline{\hspace{1.6cm}} \quad \textbf{Nota:} \underline{\hspace{0.9cm}}/15


\noindent\textbf{Comando:} Escrever uma lista de 6 palavras com encontros consonantais e usá-las em um bilhete.


\begin{enumerate}

  \item Planeje: escolha 3 palavras da unidade para usar:
  \begin{center}
  1. \underline{\hspace{3cm}} \quad 2. \underline{\hspace{3cm}} \quad 3. \underline{\hspace{3cm}}
  \end{center}

  \item Escreva o texto (mínimo 3 frases):
  \begin{center}
  \underline{\hspace{12cm}}\\\underline{\hspace{12cm}}\\\underline{\hspace{12cm}}\\\underline{\hspace{12cm}}\\\underline{\hspace{12cm}}
  \end{center}

  \item Releia e corrija: marque palavras com a grafia ENCONTROS que você usou.

  \item Ilustre o texto:
  \begin{center}
  \begin{tikzpicture}
    \draw[thick, dashed] (0,0) rectangle (12,6);
  \end{tikzpicture}
  \end{center}

\end{enumerate}

\end{exercicio}

%% PLANCHETA 14.4

\plancheta{Folha de Prática 14.4 --- Ditado e Avaliação da Unidade 14}

\begin{exercicio}[Folha de Prática 14.4 --- Ditado e Avaliação da Unidade 14]

\noindent\textbf{Nome:} \underline{\hspace{3.2cm}} \quad \textbf{Data:} \underline{\hspace{1.6cm}} \quad \textbf{Nota:} \underline{\hspace{0.9cm}}/20


\subsection*{PARTE 1: DITADO (10 itens)}

O professor dita as palavras e frases abaixo; o aluno escreve com letra cursiva.

\begin{enumerate}

  \item \underline{\hspace{10.5cm}}

  \item \underline{\hspace{10.5cm}}

  \item \underline{\hspace{10.5cm}}

  \item \underline{\hspace{10.5cm}}

  \item \underline{\hspace{10.5cm}}

  \item \underline{\hspace{10.5cm}}

\end{enumerate}

\subsection*{PARTE 2: AUTOAVALIAÇÃO DO ALUNO}

\begin{enumerate}[label=\alph*)]

  \item Eu li as palavras da unidade sem ajuda.  \quad ( ) Sim   ( ) Não

  \item Eu escrevi o ditado com as grafias certas.  \quad ( ) Sim   ( ) Não

  \item Eu separei as palavras em sílabas.  \quad ( ) Sim   ( ) Não

  \item Eu sei explicar a regra com as minhas palavras.  \quad ( ) Sim   ( ) Não

\end{enumerate}

\subsection*{PARTE 3: REGISTRO DO PROFESSOR (S/F/R/N)}

\begin{center}
\renewcommand{\arraystretch}{1.4}
\begin{tabular}{@{}p{9cm}cccc@{}}
\toprule
Aspecto observado & S & F & R & N \\
\midrule
Lê com precisão as palavras-alvo & & & & \\
Escreve com ortografia estável & & & & \\
Generaliza a regra em escrita espontânea & & & & \\
Mantém atenção durante o ditado & & & & \\
\bottomrule\end{tabular}\end{center}

\end{exercicio}

\begin{dica}
**Dica para o Professor:** A revisão dos encontros é o momento do automatismo: leitura em cadeia e jogos de carrinho (BRA-BRE-BRI) fixam o padrão.
\end{dica}

%% ============================================
%% UNIDADE 15 -- Revisão ortográfica
%% ============================================

\chapter{Unidade 15 --- Revisão ortográfica}

\metadata{I-15}{Silábico-Alfabético Avançado}{EF04LP01, EF04LP02, EF04LP03, EF04LP05, EF04LP06, EF04LP07, EF04LP09}{D1}{EF04LP01, EF04LP02, EF04LP03, EF04LP05, EF04LP06, EF04LP07, EF04LP09}{Resolucao CNE/CEB 7/2010, art. 12}

\textbf{Regra:} Entre vogais: R forte vira RR, S forte vira SS, S vira Z, e Ç aparece antes de A, O, U. A revisão transforma regra em hábito.

\textbf{Palavra geradora do bloco:} \textbf{CARROÇA} --- carroça junta RR e Ç em uma palavra só: a revisão ortográfica completa.

\section{Lição 15.1 --- Revisar a própria escrita com checklist}

\textbf{Foco da lição:} Revisar a própria escrita com checklist


\begin{boquinha}[Boquinha da Técnica]
Na revisão, compare duplas: CARA/CARRO, CASA/CAÇA, PASSA/PAÇA. A boca mantém o mesmo gesto; a escrita decide pela regra.
\end{boquinha}

\noindent\textbf{Formação guiada:}

\begin{center}
\resizebox{\textwidth}{!}{%
\begin{tabular}{ccccc}
  \sil{RR}{A} & \sil{RR}{O} & \sil{SS}{A} & \sil{SS}{O} & \textbf{ZA} \\
  \textbf{CA} & \textbf{ÇA} & \textbf{ÇO} \\
  \end{tabular}}
\end{center}

\noindent\textbf{Palavras da unidade:} carroça, serra, passa, assado, casa, caça, laço, poço.

\noindent\textbf{Frases para leitura oral:}
\begin{itemize}[leftmargin=1.6em]
  \item A carroça sobe a serra.
  \item O assado passa do ponto.
  \item A casa está perto do poço.
  \item O laço é da festa.
\end{itemize}

\section{Lição 15.2 --- Leitura do texto curto}

\noindent\textbf{Leia o texto em voz alta, com atenção às palavras da unidade:}


\begin{quote}
A carroça sobe a serra. O assado passa do ponto.

A casa está perto do poço. O laço é da festa.
\end{quote}

\textbf{Depois de ler, responda por escrito:}
\begin{enumerate}

  \item O que acontece no texto que você leu?
  \begin{center}
  \underline{\hspace{12cm}}
  \end{center}

  \item Quantas frases o texto tem?
  \begin{center}
  \underline{\hspace{12cm}}
  \end{center}

  \item Circule no texto as palavras com a grafia RR/SS/SZ/Ç.

  \item Leia o texto de novo em voz alta, agora sem ajuda.

\end{enumerate}

\section{Lição 15.3 --- Escrita com a grafia RR/SS/SZ/Ç}

\noindent\textbf{Regra da unidade:} Entre vogais: R forte vira RR, S forte vira SS, S vira Z, e Ç aparece antes de A, O, U. A revisão transforma regra em hábito.


\noindent\textbf{Ditado guiado:} o professor dita as palavras abaixo; o aluno escreve e depois corrige com o professor.

\begin{center}
  \underline{\hspace{2.3cm}} \quad \underline{\hspace{2.3cm}} \quad \underline{\hspace{2.3cm}} \quad \underline{\hspace{2.3cm}}

  \underline{\hspace{2.3cm}} \quad \underline{\hspace{2.3cm}} \quad \underline{\hspace{2.3cm}} \quad \underline{\hspace{2.3cm}}
\end{center}

\noindent\textbf{Complete as palavras com a grafia RR/SS/SZ/Ç:}
\begin{center}
  carroç\underline{~} \quad serr\underline{~} \quad pass\underline{~} \quad assad\underline{~} \quad cas\underline{~}
\end{center}

\noindent\textbf{Escreva a palavra que o professor mostrar (banco de palavras) e separe em sílabas:}

\begin{center}
  \underline{\hspace{2cm}} $=$ \underline{\hspace{2cm}} \quad \underline{\hspace{2cm}} $=$ \underline{\hspace{2cm}} \quad \underline{\hspace{2cm}} $=$ \underline{\hspace{2cm}}
  \end{center}

%% PLANCHETA 15.1

\plancheta{Folha de Prática 15.1 --- Revisão ortográfica (Parte 1)}

\begin{exercicio}[Folha de Prática 15.1 --- Revisão ortográfica (Parte 1)]

\noindent\textbf{Nome:} \underline{\hspace{3.2cm}} \quad \textbf{Data:} \underline{\hspace{1.6cm}} \quad \textbf{Nota:} \underline{\hspace{0.9cm}}/10


\subsection*{PARTE 1: RECONHECIMENTO}

\begin{enumerate}

  \item Forme as sílabas com os sons da unidade:

\begin{center}
\resizebox{\textwidth}{!}{%
\begin{tabular}{ccccc}
  \sil{RR}{A} & \sil{RR}{O} & \sil{SS}{A} & \sil{SS}{O} & \textbf{ZA} \\
  \textbf{CA} & \textbf{ÇA} & \textbf{ÇO} \\
  \end{tabular}}
\end{center}

  \item Complete a tabela de sílabas da unidade:

\begin{center}
\resizebox{\textwidth}{!}{%
\begin{tabular}{ccccc}
  \sil{RR}{A} & \sil{RR}{E} & \sil{RR}{I} & \sil{RR}{O} & \sil{RR}{U} \\
  \end{tabular}}
\end{center}

  \item Sublinhe as palavras da unidade:

\begin{center}
  \uline{carroça} \quad \uline{serra} \quad \uline{passa} \quad \uline{assado} \quad \uline{casa} \quad \uline{caça} \quad laço \quad poço \quad gato \quad sapo \quad rato \quad bola \quad menino
\end{center}

  \item Identifique o som-alvo em cada palavra:
  \begin{center}
  CARROÇA $\rightarrow$ som-alvo: \underline{\hspace{1cm}} \quad SERRA $\rightarrow$ som-alvo: \underline{\hspace{1cm}} \quad PASSA $\rightarrow$ som-alvo: \underline{\hspace{1cm}}
  \end{center}

  \item Separe em sílabas:
  \begin{center}
  CARROÇA $\rightarrow$ \underline{\hspace{1cm}} \quad SERRA $\rightarrow$ \underline{\hspace{1cm}} \quad PASSA $\rightarrow$ \underline{\hspace{1cm}}
  \end{center}

  \item Conte as sílabas:
  \begin{center}
  CARROÇA $\rightarrow$ \underline{\hspace{0.9cm}} \quad SERRA $\rightarrow$ \underline{\hspace{0.9cm}} \quad PASSA $\rightarrow$ \underline{\hspace{0.9cm}}
  \end{center}

\end{enumerate}

\subsection*{PARTE 2: PRODUÇÃO GUIADA}

\begin{enumerate}[resume]

  \item Complete com a grafia da unidade:
  \begin{center}
  RR/SS/SZ/Ç\underline{~~~} \quad RR/SS/SZ/Ç\underline{~~~} \quad RR/SS/SZ/Ç\underline{~~~} \quad RR/SS/SZ/Ç\underline{~~~} \quad RR/SS/SZ/Ç\underline{~~~}
  \end{center}

  \item Escreva 3 palavras com o som-alvo:
  \begin{center}
  1. \underline{\hspace{3cm}} \quad 2. \underline{\hspace{3cm}} \quad 3. \underline{\hspace{3cm}}
  \end{center}

  \item Copie as palavras:
  \begin{center}
  carroça \quad serra \quad passa \quad assado \quad casa
  \end{center}

  \item Separe em sílabas:
  \begin{center}
  CARROÇA = \underline{\hspace{1cm}} \underline{\hspace{1cm}} \quad SERRA = \underline{\hspace{1cm}} \underline{\hspace{1cm}} \quad PASSA = \underline{\hspace{1cm}} \underline{\hspace{1cm}}
  \end{center}

  \item Escreva 2 frases usando palavras da unidade:
  \begin{enumerate}[label=\alph*)]
    \item \underline{\hspace{8cm}}
    \item \underline{\hspace{8cm}}
  \end{enumerate}

\end{enumerate}

\end{exercicio}

%% PLANCHETA 15.2

\plancheta{Folha de Prática 15.2 --- Revisão ortográfica (Parte 2)}

\begin{exercicio}[Folha de Prática 15.2 --- Revisão ortográfica (Parte 2)]

\noindent\textbf{Nome:} \underline{\hspace{3.2cm}} \quad \textbf{Data:} \underline{\hspace{1.6cm}} \quad \textbf{Nota:} \underline{\hspace{0.9cm}}/10


\subsection*{PARTE 3: INTEGRAÇÃO}

\begin{enumerate}[resume]

  \item Leia as palavras em voz alta:
  \begin{center}
  $\rightarrow$ carroça ~ serra ~ passa ~ assado ~ casa ~ caça
  \end{center}

  \item Leia as frases em voz alta:
  \begin{center}
  $\rightarrow$ A carroça sobe a serra.\\
  $\rightarrow$ O assado passa do ponto.\\
  $\rightarrow$ A casa está perto do poço.\\
  $\rightarrow$ O laço é da festa.\\
  \end{center}

  \item Circule a opção correta:
  \begin{enumerate}[label=(\alph*)]

    \item O som da unidade é: ( ) oral ( ) nasal ( ) vibrante

    \item A primeira sílaba de CARROÇA é: ( ) RRA ( ) CASA

    \item As palavras da unidade aparecem: ( ) no início ( ) no meio ( ) no fim das palavras

  \end{enumerate}

  \item Leia e complete:
  \begin{center}
  car\underline{~~~} \quad ser\underline{~~~} \quad pas\underline{~~~} \quad ass\underline{~~~}
  \end{center}

  \item Leitura em voz alta com o professor:
  \begin{center}
  carroça \quad serra \quad passa \quad assado \quad casa \quad caça \quad laço \quad poço
  \end{center}

\end{enumerate}

\end{exercicio}

%% PLANCHETA 15.3

\plancheta{Folha de Prática 15.3 --- Produção com RR/SS/SZ/Ç}

\begin{exercicio}[Folha de Prática 15.3 --- Produção com RR/SS/SZ/Ç]

\noindent\textbf{Nome:} \underline{\hspace{3.2cm}} \quad \textbf{Data:} \underline{\hspace{1.6cm}} \quad \textbf{Nota:} \underline{\hspace{0.9cm}}/15


\noindent\textbf{Comando:} Reescrever 4 frases corrigindo os erros de ortografia (ditado de revisão).


\begin{enumerate}

  \item Planeje: escolha 3 palavras da unidade para usar:
  \begin{center}
  1. \underline{\hspace{3cm}} \quad 2. \underline{\hspace{3cm}} \quad 3. \underline{\hspace{3cm}}
  \end{center}

  \item Escreva o texto (mínimo 3 frases):
  \begin{center}
  \underline{\hspace{12cm}}\\\underline{\hspace{12cm}}\\\underline{\hspace{12cm}}\\\underline{\hspace{12cm}}\\\underline{\hspace{12cm}}
  \end{center}

  \item Releia e corrija: marque palavras com a grafia RR/SS/SZ/Ç que você usou.

  \item Ilustre o texto:
  \begin{center}
  \begin{tikzpicture}
    \draw[thick, dashed] (0,0) rectangle (12,6);
  \end{tikzpicture}
  \end{center}

\end{enumerate}

\end{exercicio}

%% PLANCHETA 15.4

\plancheta{Folha de Prática 15.4 --- Ditado e Avaliação da Unidade 15}

\begin{exercicio}[Folha de Prática 15.4 --- Ditado e Avaliação da Unidade 15]

\noindent\textbf{Nome:} \underline{\hspace{3.2cm}} \quad \textbf{Data:} \underline{\hspace{1.6cm}} \quad \textbf{Nota:} \underline{\hspace{0.9cm}}/20


\subsection*{PARTE 1: DITADO (10 itens)}

O professor dita as palavras e frases abaixo; o aluno escreve com letra cursiva.

\begin{enumerate}

  \item \underline{\hspace{10.5cm}}

  \item \underline{\hspace{10.5cm}}

  \item \underline{\hspace{10.5cm}}

  \item \underline{\hspace{10.5cm}}

  \item \underline{\hspace{10.5cm}}

  \item \underline{\hspace{10.5cm}}

\end{enumerate}

\subsection*{PARTE 2: AUTOAVALIAÇÃO DO ALUNO}

\begin{enumerate}[label=\alph*)]

  \item Eu li as palavras da unidade sem ajuda.  \quad ( ) Sim   ( ) Não

  \item Eu escrevi o ditado com as grafias certas.  \quad ( ) Sim   ( ) Não

  \item Eu separei as palavras em sílabas.  \quad ( ) Sim   ( ) Não

  \item Eu sei explicar a regra com as minhas palavras.  \quad ( ) Sim   ( ) Não

\end{enumerate}

\subsection*{PARTE 3: REGISTRO DO PROFESSOR (S/F/R/N)}

\begin{center}
\renewcommand{\arraystretch}{1.4}
\begin{tabular}{@{}p{9cm}cccc@{}}
\toprule
Aspecto observado & S & F & R & N \\
\midrule
Lê com precisão as palavras-alvo & & & & \\
Escreve com ortografia estável & & & & \\
Generaliza a regra em escrita espontânea & & & & \\
Mantém atenção durante o ditado & & & & \\
\bottomrule\end{tabular}\end{center}

\end{exercicio}

%% PLANCHETA 15.5

\plancheta{Folha de Prática 15.5 --- Leitura e Fluência da Unidade 15}

\begin{exercicio}[Folha de Prática 15.5 --- Leitura e Fluência da Unidade 15]

\noindent\textbf{Nome:} \underline{\hspace{3.2cm}} \quad \textbf{Data:} \underline{\hspace{1.6cm}} \quad \textbf{Nota:} \underline{\hspace{0.9cm}}/10


\subsection*{PARTE 1: LEITURA EM VOZ ALTA}

Leia o texto abaixo três vezes; a cada leitura, marque um círculo no cronômetro do professor.

\begin{quote}
A carroça sobe a serra. O assado passa do ponto.

A casa está perto do poço. O laço é da festa.
\end{quote}

\begin{center}
1ª leitura: \underline{\hspace{1.5cm}} \quad 2ª leitura: \underline{\hspace{1.5cm}} \quad 3ª leitura: \underline{\hspace{1.5cm}}
\end{center}

\subsection*{PARTE 2: COMPREENSÃO}

\begin{enumerate}

  \item Quem são os personagens do texto?
  \begin{center}
  \underline{\hspace{12cm}}
  \end{center}

  \item Qual é a grafia-alvo que mais aparece no texto? Escreva três exemplos:
  \begin{center}
  \underline{\hspace{3.5cm}} \quad \underline{\hspace{3.5cm}} \quad \underline{\hspace{3.5cm}}
  \end{center}

  \item Você gostou do texto? Por quê?
  \begin{center}
  \underline{\hspace{12cm}}
  \end{center}

\end{enumerate}

\subsection*{PARTE 3: DESENHO}

Desenhe uma cena do texto:

\begin{center}
\begin{tikzpicture}
  \draw[thick, dashed] (0,0) rectangle (12,5);
\end{tikzpicture}
\end{center}

\end{exercicio}

\begin{dica}
**Dica para o Professor:** Na revisão, o aluno consulta a regra e justifica a escolha: 'R entre vogais, som forte, escrevo RR'. O hábito vem da justificativa.
\end{dica}

%% ============================================
%% UNIDADE 16 -- Interpretação: ideias principais
%% ============================================

\chapter{Unidade 16 --- Interpretação: ideias principais}

\metadata{I-16}{Silábico-Alfabético Avançado}{EF04LP01, EF04LP02, EF04LP03, EF04LP05, EF04LP06, EF04LP07, EF04LP09}{D1}{EF04LP01, EF04LP02, EF04LP03, EF04LP05, EF04LP06, EF04LP07, EF04LP09}{Resolucao CNE/CEB 7/2010, art. 12}

\textbf{Regra:} Poemas e cantigas têm ritmo, rima e repetição. Ler com expressividade mostra que o aluno compreende o texto.

\textbf{Palavra geradora do bloco:} \textbf{CANTIGA} --- a cantiga é a memória oral do povo; une leitura, ritmo e alegria.

\section{Lição 16.1 --- Localizar a ideia principal de cada parágrafo}

\textbf{Foco da lição:} Localizar a ideia principal de cada parágrafo


\begin{boquinha}[Boquinha da Técnica]
Na leitura de poema, a boca acompanha a música do verso: subir e descer a voz, pausar nas rimas. Ler como quem canta.
\end{boquinha}

\noindent\textbf{Formação guiada:}

\begin{center}
\resizebox{\textwidth}{!}{%
\begin{tabular}{cccc}
  \textbf{CAN} & \textbf{TIGA} & \textbf{RIMA} & \textbf{POE} \\
  \end{tabular}}
\end{center}

\noindent\textbf{Palavras da unidade:} cantiga, poema, rima, estrofe, verso, cantor, melodia, refrão.

\noindent\textbf{Frases para leitura oral:}
\begin{itemize}[leftmargin=1.6em]
  \item A cantiga tem uma rima.
  \item O poema fala do vento.
  \item O refrão repete a melodia.
  \item O cantor sorri no verso.
\end{itemize}

\section{Lição 16.2 --- Leitura do texto curto}

\noindent\textbf{Leia o texto em voz alta, com atenção às palavras da unidade:}


\begin{quote}
A cantiga tem uma rima. O poema fala do vento.

O refrão repete a melodia. O cantor sorri no verso.
\end{quote}

\textbf{Depois de ler, responda por escrito:}
\begin{enumerate}

  \item O que acontece no texto que você leu?
  \begin{center}
  \underline{\hspace{12cm}}
  \end{center}

  \item Quantas frases o texto tem?
  \begin{center}
  \underline{\hspace{12cm}}
  \end{center}

  \item Circule no texto as palavras com a grafia POEMA.

  \item Leia o texto de novo em voz alta, agora sem ajuda.

\end{enumerate}

\section{Lição 16.3 --- Escrita com a grafia POEMA}

\noindent\textbf{Regra da unidade:} Poemas e cantigas têm ritmo, rima e repetição. Ler com expressividade mostra que o aluno compreende o texto.


\noindent\textbf{Ditado guiado:} o professor dita as palavras abaixo; o aluno escreve e depois corrige com o professor.

\begin{center}
  \underline{\hspace{2.3cm}} \quad \underline{\hspace{2.3cm}} \quad \underline{\hspace{2.3cm}} \quad \underline{\hspace{2.3cm}}

  \underline{\hspace{2.3cm}} \quad \underline{\hspace{2.3cm}} \quad \underline{\hspace{2.3cm}} \quad \underline{\hspace{2.3cm}}
\end{center}

\noindent\textbf{Complete as palavras com a grafia POEMA:}
\begin{center}
  cantig\underline{~} \quad poem\underline{~} \quad rim\underline{~} \quad estrof\underline{~} \quad vers\underline{~}
\end{center}

\noindent\textbf{Escreva a palavra que o professor mostrar (banco de palavras) e separe em sílabas:}

\begin{center}
  \underline{\hspace{2cm}} $=$ \underline{\hspace{2cm}} \quad \underline{\hspace{2cm}} $=$ \underline{\hspace{2cm}} \quad \underline{\hspace{2cm}} $=$ \underline{\hspace{2cm}}
  \end{center}

%% PLANCHETA 16.1

\plancheta{Folha de Prática 16.1 --- Interpretação: ideias principais (Parte 1)}

\begin{exercicio}[Folha de Prática 16.1 --- Interpretação: ideias principais (Parte 1)]

\noindent\textbf{Nome:} \underline{\hspace{3.2cm}} \quad \textbf{Data:} \underline{\hspace{1.6cm}} \quad \textbf{Nota:} \underline{\hspace{0.9cm}}/10


\subsection*{PARTE 1: RECONHECIMENTO}

\begin{enumerate}

  \item Forme as sílabas com os sons da unidade:

\begin{center}
\resizebox{\textwidth}{!}{%
\begin{tabular}{cccc}
  \textbf{CAN} & \textbf{TIGA} & \textbf{RIMA} & \textbf{POE} \\
  \end{tabular}}
\end{center}

  \item Sublinhe as palavras da unidade:

\begin{center}
  \uline{cantiga} \quad \uline{poema} \quad \uline{rima} \quad \uline{estrofe} \quad \uline{verso} \quad \uline{cantor} \quad melodia \quad refrão \quad gato \quad sapo \quad rato \quad bola \quad menino
\end{center}

  \item Identifique o som-alvo em cada palavra:
  \begin{center}
  CANTIGA $\rightarrow$ som-alvo: \underline{\hspace{1cm}} \quad POEMA $\rightarrow$ som-alvo: \underline{\hspace{1cm}} \quad RIMA $\rightarrow$ som-alvo: \underline{\hspace{1cm}}
  \end{center}

  \item Separe em sílabas:
  \begin{center}
  CANTIGA $\rightarrow$ \underline{\hspace{1cm}} \quad POEMA $\rightarrow$ \underline{\hspace{1cm}} \quad RIMA $\rightarrow$ \underline{\hspace{1cm}}
  \end{center}

  \item Conte as sílabas:
  \begin{center}
  CANTIGA $\rightarrow$ \underline{\hspace{0.9cm}} \quad POEMA $\rightarrow$ \underline{\hspace{0.9cm}} \quad RIMA $\rightarrow$ \underline{\hspace{0.9cm}}
  \end{center}

\end{enumerate}

\subsection*{PARTE 2: PRODUÇÃO GUIADA}

\begin{enumerate}[resume]

  \item Complete com a grafia da unidade:
  \begin{center}
  POEMA\underline{~~~} \quad POEMA\underline{~~~} \quad POEMA\underline{~~~} \quad POEMA\underline{~~~} \quad POEMA\underline{~~~}
  \end{center}

  \item Escreva 3 palavras com o som-alvo:
  \begin{center}
  1. \underline{\hspace{3cm}} \quad 2. \underline{\hspace{3cm}} \quad 3. \underline{\hspace{3cm}}
  \end{center}

  \item Copie as palavras:
  \begin{center}
  cantiga \quad poema \quad rima \quad estrofe \quad verso
  \end{center}

  \item Separe em sílabas:
  \begin{center}
  CANTIGA = \underline{\hspace{1cm}} \underline{\hspace{1cm}} \quad POEMA = \underline{\hspace{1cm}} \underline{\hspace{1cm}} \quad RIMA = \underline{\hspace{1cm}} \underline{\hspace{1cm}}
  \end{center}

  \item Escreva 2 frases usando palavras da unidade:
  \begin{enumerate}[label=\alph*)]
    \item \underline{\hspace{8cm}}
    \item \underline{\hspace{8cm}}
  \end{enumerate}

\end{enumerate}

\end{exercicio}

%% PLANCHETA 16.2

\plancheta{Folha de Prática 16.2 --- Interpretação: ideias principais (Parte 2)}

\begin{exercicio}[Folha de Prática 16.2 --- Interpretação: ideias principais (Parte 2)]

\noindent\textbf{Nome:} \underline{\hspace{3.2cm}} \quad \textbf{Data:} \underline{\hspace{1.6cm}} \quad \textbf{Nota:} \underline{\hspace{0.9cm}}/10


\subsection*{PARTE 3: INTEGRAÇÃO}

\begin{enumerate}[resume]

  \item Leia as palavras em voz alta:
  \begin{center}
  $\rightarrow$ cantiga ~ poema ~ rima ~ estrofe ~ verso ~ cantor
  \end{center}

  \item Leia as frases em voz alta:
  \begin{center}
  $\rightarrow$ A cantiga tem uma rima.\\
  $\rightarrow$ O poema fala do vento.\\
  $\rightarrow$ O refrão repete a melodia.\\
  $\rightarrow$ O cantor sorri no verso.\\
  \end{center}

  \item Circule a opção correta:
  \begin{enumerate}[label=(\alph*)]

    \item O som da unidade é: ( ) oral ( ) nasal 

    \item A primeira sílaba de CANTIGA é: ( ) CAN ( ) RIMA

    \item As palavras da unidade aparecem: ( ) no início ( ) no meio ( ) no fim das palavras

  \end{enumerate}

  \item Leia e complete:
  \begin{center}
  can\underline{~~~} \quad poe\underline{~~~} \quad rim\underline{~~~} \quad est\underline{~~~}
  \end{center}

  \item Leitura em voz alta com o professor:
  \begin{center}
  cantiga \quad poema \quad rima \quad estrofe \quad verso \quad cantor \quad melodia \quad refrão
  \end{center}

\end{enumerate}

\end{exercicio}

%% PLANCHETA 16.3

\plancheta{Folha de Prática 16.3 --- Produção com POEMA}

\begin{exercicio}[Folha de Prática 16.3 --- Produção com POEMA]

\noindent\textbf{Nome:} \underline{\hspace{3.2cm}} \quad \textbf{Data:} \underline{\hspace{1.6cm}} \quad \textbf{Nota:} \underline{\hspace{0.9cm}}/15


\noindent\textbf{Comando:} Recitar um poema curto para a turma e escrever o próprio poema com 2 versos e rima.


\begin{enumerate}

  \item Planeje: escolha 3 palavras da unidade para usar:
  \begin{center}
  1. \underline{\hspace{3cm}} \quad 2. \underline{\hspace{3cm}} \quad 3. \underline{\hspace{3cm}}
  \end{center}

  \item Escreva o texto (mínimo 3 frases):
  \begin{center}
  \underline{\hspace{12cm}}\\\underline{\hspace{12cm}}\\\underline{\hspace{12cm}}\\\underline{\hspace{12cm}}\\\underline{\hspace{12cm}}
  \end{center}

  \item Releia e corrija: marque palavras com a grafia POEMA que você usou.

  \item Ilustre o texto:
  \begin{center}
  \begin{tikzpicture}
    \draw[thick, dashed] (0,0) rectangle (12,6);
  \end{tikzpicture}
  \end{center}

\end{enumerate}

\end{exercicio}

%% PLANCHETA 16.4

\plancheta{Folha de Prática 16.4 --- Ditado e Avaliação da Unidade 16}

\begin{exercicio}[Folha de Prática 16.4 --- Ditado e Avaliação da Unidade 16]

\noindent\textbf{Nome:} \underline{\hspace{3.2cm}} \quad \textbf{Data:} \underline{\hspace{1.6cm}} \quad \textbf{Nota:} \underline{\hspace{0.9cm}}/20


\subsection*{PARTE 1: DITADO (10 itens)}

O professor dita as palavras e frases abaixo; o aluno escreve com letra cursiva.

\begin{enumerate}

  \item \underline{\hspace{10.5cm}}

  \item \underline{\hspace{10.5cm}}

  \item \underline{\hspace{10.5cm}}

  \item \underline{\hspace{10.5cm}}

  \item \underline{\hspace{10.5cm}}

  \item \underline{\hspace{10.5cm}}

\end{enumerate}

\subsection*{PARTE 2: AUTOAVALIAÇÃO DO ALUNO}

\begin{enumerate}[label=\alph*)]

  \item Eu li as palavras da unidade sem ajuda.  \quad ( ) Sim   ( ) Não

  \item Eu escrevi o ditado com as grafias certas.  \quad ( ) Sim   ( ) Não

  \item Eu separei as palavras em sílabas.  \quad ( ) Sim   ( ) Não

  \item Eu sei explicar a regra com as minhas palavras.  \quad ( ) Sim   ( ) Não

\end{enumerate}

\subsection*{PARTE 3: REGISTRO DO PROFESSOR (S/F/R/N)}

\begin{center}
\renewcommand{\arraystretch}{1.4}
\begin{tabular}{@{}p{9cm}cccc@{}}
\toprule
Aspecto observado & S & F & R & N \\
\midrule
Lê com precisão as palavras-alvo & & & & \\
Escreve com ortografia estável & & & & \\
Generaliza a regra em escrita espontânea & & & & \\
Mantém atenção durante o ditado & & & & \\
\bottomrule\end{tabular}\end{center}

\end{exercicio}

\begin{dica}
**Dica para o Professor:** A leitura expressiva é a ponte para a fluência: grave a recitação da criança e deixe que ela se ouça; o progresso fica evidente.
\end{dica}

%% ============================================
%% UNIDADE 18 -- Produção: a carta
%% ============================================

\chapter{Unidade 18 --- Produção: a carta}

\metadata{I-18}{Silábico-Alfabético Avançado}{EF04LP01, EF04LP02, EF04LP03, EF04LP05, EF04LP06, EF04LP07, EF04LP09}{D1}{EF04LP01, EF04LP02, EF04LP03, EF04LP05, EF04LP06, EF04LP07, EF04LP09}{Resolucao CNE/CEB 7/2010, art. 12}

\textbf{Regra:} A fábula é uma história curta com animais que agem como pessoas e termina com uma moral, um ensinamento.

\textbf{Palavra geradora do bloco:} \textbf{LEÃO} --- o leão é rei das fábulas; ajuda a trabalhar o encontro final e a moral.

\section{Lição 18.1 --- Produzir cartas com remetente e destinatário}

\textbf{Foco da lição:} Produzir cartas com remetente e destinatário


\begin{boquinha}[Boquinha da Técnica]
Ao ler a fábula, a boca dá voz aos animais: voz grossa para o leão, leve para o rato. A moral pede leitura pausada.
\end{boquinha}

\noindent\textbf{Formação guiada:}

\begin{center}
\resizebox{\textwidth}{!}{%
\begin{tabular}{cccc}
  \textbf{LE} & \textbf{ÃO} & \textbf{RA} & \textbf{TO} \\
  \end{tabular}}
\end{center}

\noindent\textbf{Palavras da unidade:} leão, rato, formiga, cigarra, moral, fábula, corrida, sabedoria.

\noindent\textbf{Frases para leitura oral:}
\begin{itemize}[leftmargin=1.6em]
  \item O leão poupou o rato.
  \item A formiga trabalhou no verão.
  \item A cigarra cantou no frio.
  \item A moral ensina a lição.
\end{itemize}

\section{Lição 18.2 --- Leitura do texto curto}

\noindent\textbf{Leia o texto em voz alta, com atenção às palavras da unidade:}


\begin{quote}
O leão poupou o rato. A formiga trabalhou no verão.

A cigarra cantou no frio. A moral ensina a lição.
\end{quote}

\textbf{Depois de ler, responda por escrito:}
\begin{enumerate}

  \item O que acontece no texto que você leu?
  \begin{center}
  \underline{\hspace{12cm}}
  \end{center}

  \item Quantas frases o texto tem?
  \begin{center}
  \underline{\hspace{12cm}}
  \end{center}

  \item Circule no texto as palavras com a grafia FÁBULA.

  \item Leia o texto de novo em voz alta, agora sem ajuda.

\end{enumerate}

\section{Lição 18.3 --- Escrita com a grafia FÁBULA}

\noindent\textbf{Regra da unidade:} A fábula é uma história curta com animais que agem como pessoas e termina com uma moral, um ensinamento.


\noindent\textbf{Ditado guiado:} o professor dita as palavras abaixo; o aluno escreve e depois corrige com o professor.

\begin{center}
  \underline{\hspace{2.3cm}} \quad \underline{\hspace{2.3cm}} \quad \underline{\hspace{2.3cm}} \quad \underline{\hspace{2.3cm}}

  \underline{\hspace{2.3cm}} \quad \underline{\hspace{2.3cm}} \quad \underline{\hspace{2.3cm}} \quad \underline{\hspace{2.3cm}}
\end{center}

\noindent\textbf{Complete as palavras com a grafia FÁBULA:}
\begin{center}
  leã\underline{~} \quad rat\underline{~} \quad formig\underline{~} \quad cigarr\underline{~} \quad mora\underline{~}
\end{center}

\noindent\textbf{Escreva a palavra que o professor mostrar (banco de palavras) e separe em sílabas:}

\begin{center}
  \underline{\hspace{2cm}} $=$ \underline{\hspace{2cm}} \quad \underline{\hspace{2cm}} $=$ \underline{\hspace{2cm}} \quad \underline{\hspace{2cm}} $=$ \underline{\hspace{2cm}}
  \end{center}

%% PLANCHETA 18.1

\plancheta{Folha de Prática 18.1 --- Produção: a carta (Parte 1)}

\begin{exercicio}[Folha de Prática 18.1 --- Produção: a carta (Parte 1)]

\noindent\textbf{Nome:} \underline{\hspace{3.2cm}} \quad \textbf{Data:} \underline{\hspace{1.6cm}} \quad \textbf{Nota:} \underline{\hspace{0.9cm}}/10


\subsection*{PARTE 1: RECONHECIMENTO}

\begin{enumerate}

  \item Forme as sílabas com os sons da unidade:

\begin{center}
\resizebox{\textwidth}{!}{%
\begin{tabular}{cccc}
  \textbf{LE} & \textbf{ÃO} & \textbf{RA} & \textbf{TO} \\
  \end{tabular}}
\end{center}

  \item Sublinhe as palavras da unidade:

\begin{center}
  \uline{leão} \quad \uline{rato} \quad \uline{formiga} \quad \uline{cigarra} \quad \uline{moral} \quad \uline{fábula} \quad corrida \quad sabedoria \quad gato \quad sapo \quad \uline{rato} \quad bola \quad menino
\end{center}

  \item Identifique o som-alvo em cada palavra:
  \begin{center}
  LEÃO $\rightarrow$ som-alvo: \underline{\hspace{1cm}} \quad RATO $\rightarrow$ som-alvo: \underline{\hspace{1cm}} \quad FORMIGA $\rightarrow$ som-alvo: \underline{\hspace{1cm}}
  \end{center}

  \item Separe em sílabas:
  \begin{center}
  LEÃO $\rightarrow$ \underline{\hspace{1cm}} \quad RATO $\rightarrow$ \underline{\hspace{1cm}} \quad FORMIGA $\rightarrow$ \underline{\hspace{1cm}}
  \end{center}

  \item Conte as sílabas:
  \begin{center}
  LEÃO $\rightarrow$ \underline{\hspace{0.9cm}} \quad RATO $\rightarrow$ \underline{\hspace{0.9cm}} \quad FORMIGA $\rightarrow$ \underline{\hspace{0.9cm}}
  \end{center}

\end{enumerate}

\subsection*{PARTE 2: PRODUÇÃO GUIADA}

\begin{enumerate}[resume]

  \item Complete com a grafia da unidade:
  \begin{center}
  FÁBULA\underline{~~~} \quad FÁBULA\underline{~~~} \quad FÁBULA\underline{~~~} \quad FÁBULA\underline{~~~} \quad FÁBULA\underline{~~~}
  \end{center}

  \item Escreva 3 palavras com o som-alvo:
  \begin{center}
  1. \underline{\hspace{3cm}} \quad 2. \underline{\hspace{3cm}} \quad 3. \underline{\hspace{3cm}}
  \end{center}

  \item Copie as palavras:
  \begin{center}
  leão \quad rato \quad formiga \quad cigarra \quad moral
  \end{center}

  \item Separe em sílabas:
  \begin{center}
  LEÃO = \underline{\hspace{1cm}} \underline{\hspace{1cm}} \quad RATO = \underline{\hspace{1cm}} \underline{\hspace{1cm}} \quad FORMIGA = \underline{\hspace{1cm}} \underline{\hspace{1cm}}
  \end{center}

  \item Escreva 2 frases usando palavras da unidade:
  \begin{enumerate}[label=\alph*)]
    \item \underline{\hspace{8cm}}
    \item \underline{\hspace{8cm}}
  \end{enumerate}

\end{enumerate}

\end{exercicio}

%% PLANCHETA 18.2

\plancheta{Folha de Prática 18.2 --- Produção: a carta (Parte 2)}

\begin{exercicio}[Folha de Prática 18.2 --- Produção: a carta (Parte 2)]

\noindent\textbf{Nome:} \underline{\hspace{3.2cm}} \quad \textbf{Data:} \underline{\hspace{1.6cm}} \quad \textbf{Nota:} \underline{\hspace{0.9cm}}/10


\subsection*{PARTE 3: INTEGRAÇÃO}

\begin{enumerate}[resume]

  \item Leia as palavras em voz alta:
  \begin{center}
  $\rightarrow$ leão ~ rato ~ formiga ~ cigarra ~ moral ~ fábula
  \end{center}

  \item Leia as frases em voz alta:
  \begin{center}
  $\rightarrow$ O leão poupou o rato.\\
  $\rightarrow$ A formiga trabalhou no verão.\\
  $\rightarrow$ A cigarra cantou no frio.\\
  $\rightarrow$ A moral ensina a lição.\\
  \end{center}

  \item Circule a opção correta:
  \begin{enumerate}[label=(\alph*)]

    \item O som da unidade é: ( ) oral ( ) nasal 

    \item A primeira sílaba de LEÃO é: ( ) LE ( ) ÃO

    \item As palavras da unidade aparecem: ( ) no início ( ) no meio ( ) no fim das palavras

  \end{enumerate}

  \item Leia e complete:
  \begin{center}
  leã\underline{~~~} \quad rat\underline{~~~} \quad for\underline{~~~} \quad cig\underline{~~~}
  \end{center}

  \item Leitura em voz alta com o professor:
  \begin{center}
  leão \quad rato \quad formiga \quad cigarra \quad moral \quad fábula \quad corrida \quad sabedoria
  \end{center}

\end{enumerate}

\end{exercicio}

%% PLANCHETA 18.3

\plancheta{Folha de Prática 18.3 --- Produção com FÁBULA}

\begin{exercicio}[Folha de Prática 18.3 --- Produção com FÁBULA]

\noindent\textbf{Nome:} \underline{\hspace{3.2cm}} \quad \textbf{Data:} \underline{\hspace{1.6cm}} \quad \textbf{Nota:} \underline{\hspace{0.9cm}}/15


\noindent\textbf{Comando:} Explicar com suas palavras o que a moral da fábula significa na vida real.


\begin{enumerate}

  \item Planeje: escolha 3 palavras da unidade para usar:
  \begin{center}
  1. \underline{\hspace{3cm}} \quad 2. \underline{\hspace{3cm}} \quad 3. \underline{\hspace{3cm}}
  \end{center}

  \item Escreva o texto (mínimo 3 frases):
  \begin{center}
  \underline{\hspace{12cm}}\\\underline{\hspace{12cm}}\\\underline{\hspace{12cm}}\\\underline{\hspace{12cm}}\\\underline{\hspace{12cm}}
  \end{center}

  \item Releia e corrija: marque palavras com a grafia FÁBULA que você usou.

  \item Ilustre o texto:
  \begin{center}
  \begin{tikzpicture}
    \draw[thick, dashed] (0,0) rectangle (12,6);
  \end{tikzpicture}
  \end{center}

\end{enumerate}

\end{exercicio}

%% PLANCHETA 18.4

\plancheta{Folha de Prática 18.4 --- Ditado e Avaliação da Unidade 18}

\begin{exercicio}[Folha de Prática 18.4 --- Ditado e Avaliação da Unidade 18]

\noindent\textbf{Nome:} \underline{\hspace{3.2cm}} \quad \textbf{Data:} \underline{\hspace{1.6cm}} \quad \textbf{Nota:} \underline{\hspace{0.9cm}}/20


\subsection*{PARTE 1: DITADO (10 itens)}

O professor dita as palavras e frases abaixo; o aluno escreve com letra cursiva.

\begin{enumerate}

  \item \underline{\hspace{10.5cm}}

  \item \underline{\hspace{10.5cm}}

  \item \underline{\hspace{10.5cm}}

  \item \underline{\hspace{10.5cm}}

  \item \underline{\hspace{10.5cm}}

  \item \underline{\hspace{10.5cm}}

\end{enumerate}

\subsection*{PARTE 2: AUTOAVALIAÇÃO DO ALUNO}

\begin{enumerate}[label=\alph*)]

  \item Eu li as palavras da unidade sem ajuda.  \quad ( ) Sim   ( ) Não

  \item Eu escrevi o ditado com as grafias certas.  \quad ( ) Sim   ( ) Não

  \item Eu separei as palavras em sílabas.  \quad ( ) Sim   ( ) Não

  \item Eu sei explicar a regra com as minhas palavras.  \quad ( ) Sim   ( ) Não

\end{enumerate}

\subsection*{PARTE 3: REGISTRO DO PROFESSOR (S/F/R/N)}

\begin{center}
\renewcommand{\arraystretch}{1.4}
\begin{tabular}{@{}p{9cm}cccc@{}}
\toprule
Aspecto observado & S & F & R & N \\
\midrule
Lê com precisão as palavras-alvo & & & & \\
Escreve com ortografia estável & & & & \\
Generaliza a regra em escrita espontânea & & & & \\
Mantém atenção durante o ditado & & & & \\
\bottomrule\end{tabular}\end{center}

\end{exercicio}

%% PLANCHETA 18.5

\plancheta{Folha de Prática 18.5 --- Leitura e Fluência da Unidade 18}

\begin{exercicio}[Folha de Prática 18.5 --- Leitura e Fluência da Unidade 18]

\noindent\textbf{Nome:} \underline{\hspace{3.2cm}} \quad \textbf{Data:} \underline{\hspace{1.6cm}} \quad \textbf{Nota:} \underline{\hspace{0.9cm}}/10


\subsection*{PARTE 1: LEITURA EM VOZ ALTA}

Leia o texto abaixo três vezes; a cada leitura, marque um círculo no cronômetro do professor.

\begin{quote}
O leão poupou o rato. A formiga trabalhou no verão.

A cigarra cantou no frio. A moral ensina a lição.
\end{quote}

\begin{center}
1ª leitura: \underline{\hspace{1.5cm}} \quad 2ª leitura: \underline{\hspace{1.5cm}} \quad 3ª leitura: \underline{\hspace{1.5cm}}
\end{center}

\subsection*{PARTE 2: COMPREENSÃO}

\begin{enumerate}

  \item Quem são os personagens do texto?
  \begin{center}
  \underline{\hspace{12cm}}
  \end{center}

  \item Qual é a grafia-alvo que mais aparece no texto? Escreva três exemplos:
  \begin{center}
  \underline{\hspace{3.5cm}} \quad \underline{\hspace{3.5cm}} \quad \underline{\hspace{3.5cm}}
  \end{center}

  \item Você gostou do texto? Por quê?
  \begin{center}
  \underline{\hspace{12cm}}
  \end{center}

\end{enumerate}

\subsection*{PARTE 3: DESENHO}

Desenhe uma cena do texto:

\begin{center}
\begin{tikzpicture}
  \draw[thick, dashed] (0,0) rectangle (12,5);
\end{tikzpicture}
\end{center}

\end{exercicio}

\begin{dica}
**Dica para o Professor:** Após a leitura, converse com a turma: 'Já aconteceu algo parecido com você?' A relação com a vida garante a compreensão profunda.
\end{dica}

%% ============================================
%% UNIDADE 19 -- Concordância básica
%% ============================================

\chapter{Unidade 19 --- Concordância básica}

\metadata{I-19}{Silábico-Alfabético Avançado}{EF04LP01, EF04LP02, EF04LP03, EF04LP05, EF04LP06, EF04LP07, EF04LP09}{D1}{EF04LP01, EF04LP02, EF04LP03, EF04LP05, EF04LP06, EF04LP07, EF04LP09}{Resolucao CNE/CEB 7/2010, art. 12}

\textbf{Regra:} O trava-língua repete sons parecidos para desafiar a boca: 'O rato roeu a roupa do rei de Roma'. A articulação clara é o objetivo.

\textbf{Palavra geradora do bloco:} \textbf{RATO} --- o rato aparece no trava-língua clássico; a repetição de R treina a vibração.

\section{Lição 19.1 --- Concordar sujeito e verbo em frases curtas}

\textbf{Foco da lição:} Concordar sujeito e verbo em frases curtas


\begin{boquinha}[Boquinha da Técnica]
No trava-língua, a boca precisa ser rápida e precisa; comece devagar e aumente a velocidade. O espelho ajuda a ver os lábios trabalhando.
\end{boquinha}

\noindent\textbf{Formação guiada:}

\begin{center}
\resizebox{\textwidth}{!}{%
\begin{tabular}{ccccc}
  \textbf{RA} & \textbf{RO} & \textbf{RUI} & \textbf{ROM} & \textbf{RIM} \\
  \end{tabular}}
\end{center}

\noindent\textbf{Palavras da unidade:} rato, roupa, rei, Roma, roeu, rápido, roda, risada.

\noindent\textbf{Frases para leitura oral:}
\begin{itemize}[leftmargin=1.6em]
  \item O rato roeu a roupa do rei.
  \item A roda gira rápido na rua.
  \item O rei de Roma riu da risada.
  \item Três pratos de trigo para três tigres.
\end{itemize}

\section{Lição 19.2 --- Leitura do texto curto}

\noindent\textbf{Leia o texto em voz alta, com atenção às palavras da unidade:}


\begin{quote}
O rato roeu a roupa do rei. A roda gira rápido na rua.

O rei de Roma riu da risada. Três pratos de trigo para três tigres.
\end{quote}

\textbf{Depois de ler, responda por escrito:}
\begin{enumerate}

  \item O que acontece no texto que você leu?
  \begin{center}
  \underline{\hspace{12cm}}
  \end{center}

  \item Quantas frases o texto tem?
  \begin{center}
  \underline{\hspace{12cm}}
  \end{center}

  \item Circule no texto as palavras com a grafia TRAVA-LÍNGUA.

  \item Leia o texto de novo em voz alta, agora sem ajuda.

\end{enumerate}

\section{Lição 19.3 --- Escrita com a grafia TRAVA-LÍNGUA}

\noindent\textbf{Regra da unidade:} O trava-língua repete sons parecidos para desafiar a boca: 'O rato roeu a roupa do rei de Roma'. A articulação clara é o objetivo.


\noindent\textbf{Ditado guiado:} o professor dita as palavras abaixo; o aluno escreve e depois corrige com o professor.

\begin{center}
  \underline{\hspace{2.3cm}} \quad \underline{\hspace{2.3cm}} \quad \underline{\hspace{2.3cm}} \quad \underline{\hspace{2.3cm}}

  \underline{\hspace{2.3cm}} \quad \underline{\hspace{2.3cm}} \quad \underline{\hspace{2.3cm}} \quad \underline{\hspace{2.3cm}}
\end{center}

\noindent\textbf{Complete as palavras com a grafia TRAVA-LÍNGUA:}
\begin{center}
  rat\underline{~} \quad roup\underline{~} \quad re\underline{~} \quad Rom\underline{~} \quad roe\underline{~}
\end{center}

\noindent\textbf{Escreva a palavra que o professor mostrar (banco de palavras) e separe em sílabas:}

\begin{center}
  \underline{\hspace{2cm}} $=$ \underline{\hspace{2cm}} \quad \underline{\hspace{2cm}} $=$ \underline{\hspace{2cm}} \quad \underline{\hspace{2cm}} $=$ \underline{\hspace{2cm}}
  \end{center}

%% PLANCHETA 19.1

\plancheta{Folha de Prática 19.1 --- Concordância básica (Parte 1)}

\begin{exercicio}[Folha de Prática 19.1 --- Concordância básica (Parte 1)]

\noindent\textbf{Nome:} \underline{\hspace{3.2cm}} \quad \textbf{Data:} \underline{\hspace{1.6cm}} \quad \textbf{Nota:} \underline{\hspace{0.9cm}}/10


\subsection*{PARTE 1: RECONHECIMENTO}

\begin{enumerate}

  \item Forme as sílabas com os sons da unidade:

\begin{center}
\resizebox{\textwidth}{!}{%
\begin{tabular}{ccccc}
  \textbf{RA} & \textbf{RO} & \textbf{RUI} & \textbf{ROM} & \textbf{RIM} \\
  \end{tabular}}
\end{center}

  \item Sublinhe as palavras da unidade:

\begin{center}
  \uline{rato} \quad \uline{roupa} \quad \uline{rei} \quad \uline{Roma} \quad \uline{roeu} \quad \uline{rápido} \quad roda \quad risada \quad gato \quad sapo \quad \uline{rato} \quad bola \quad menino
\end{center}

  \item Identifique o som-alvo em cada palavra:
  \begin{center}
  RATO $\rightarrow$ som-alvo: \underline{\hspace{1cm}} \quad ROUPA $\rightarrow$ som-alvo: \underline{\hspace{1cm}} \quad REI $\rightarrow$ som-alvo: \underline{\hspace{1cm}}
  \end{center}

  \item Separe em sílabas:
  \begin{center}
  RATO $\rightarrow$ \underline{\hspace{1cm}} \quad ROUPA $\rightarrow$ \underline{\hspace{1cm}} \quad REI $\rightarrow$ \underline{\hspace{1cm}}
  \end{center}

  \item Conte as sílabas:
  \begin{center}
  RATO $\rightarrow$ \underline{\hspace{0.9cm}} \quad ROUPA $\rightarrow$ \underline{\hspace{0.9cm}} \quad REI $\rightarrow$ \underline{\hspace{0.9cm}}
  \end{center}

\end{enumerate}

\subsection*{PARTE 2: PRODUÇÃO GUIADA}

\begin{enumerate}[resume]

  \item Complete com a grafia da unidade:
  \begin{center}
  TRAVA-LÍNGUA\underline{~~~} \quad TRAVA-LÍNGUA\underline{~~~} \quad TRAVA-LÍNGUA\underline{~~~} \quad TRAVA-LÍNGUA\underline{~~~} \quad TRAVA-LÍNGUA\underline{~~~}
  \end{center}

  \item Escreva 3 palavras com o som-alvo:
  \begin{center}
  1. \underline{\hspace{3cm}} \quad 2. \underline{\hspace{3cm}} \quad 3. \underline{\hspace{3cm}}
  \end{center}

  \item Copie as palavras:
  \begin{center}
  rato \quad roupa \quad rei \quad Roma \quad roeu
  \end{center}

  \item Separe em sílabas:
  \begin{center}
  RATO = \underline{\hspace{1cm}} \underline{\hspace{1cm}} \quad ROUPA = \underline{\hspace{1cm}} \underline{\hspace{1cm}} \quad REI = \underline{\hspace{1cm}} \underline{\hspace{1cm}}
  \end{center}

  \item Escreva 2 frases usando palavras da unidade:
  \begin{enumerate}[label=\alph*)]
    \item \underline{\hspace{8cm}}
    \item \underline{\hspace{8cm}}
  \end{enumerate}

\end{enumerate}

\end{exercicio}

%% PLANCHETA 19.2

\plancheta{Folha de Prática 19.2 --- Concordância básica (Parte 2)}

\begin{exercicio}[Folha de Prática 19.2 --- Concordância básica (Parte 2)]

\noindent\textbf{Nome:} \underline{\hspace{3.2cm}} \quad \textbf{Data:} \underline{\hspace{1.6cm}} \quad \textbf{Nota:} \underline{\hspace{0.9cm}}/10


\subsection*{PARTE 3: INTEGRAÇÃO}

\begin{enumerate}[resume]

  \item Leia as palavras em voz alta:
  \begin{center}
  $\rightarrow$ rato ~ roupa ~ rei ~ Roma ~ roeu ~ rápido
  \end{center}

  \item Leia as frases em voz alta:
  \begin{center}
  $\rightarrow$ O rato roeu a roupa do rei.\\
  $\rightarrow$ A roda gira rápido na rua.\\
  $\rightarrow$ O rei de Roma riu da risada.\\
  $\rightarrow$ Três pratos de trigo para três tigres.\\
  \end{center}

  \item Circule a opção correta:
  \begin{enumerate}[label=(\alph*)]

    \item O som da unidade é: ( ) oral ( ) nasal ( ) vibrante

    \item A primeira sílaba de RATO é: ( ) RA ( ) RO

    \item As palavras da unidade aparecem: ( ) no início ( ) no meio ( ) no fim das palavras

  \end{enumerate}

  \item Leia e complete:
  \begin{center}
  rat\underline{~~~} \quad rou\underline{~~~} \quad rei\underline{~~~} \quad Rom\underline{~~~}
  \end{center}

  \item Leitura em voz alta com o professor:
  \begin{center}
  rato \quad roupa \quad rei \quad Roma \quad roeu \quad rápido \quad roda \quad risada
  \end{center}

\end{enumerate}

\end{exercicio}

%% PLANCHETA 19.3

\plancheta{Folha de Prática 19.3 --- Produção com TRAVA-LÍNGUA}

\begin{exercicio}[Folha de Prática 19.3 --- Produção com TRAVA-LÍNGUA]

\noindent\textbf{Nome:} \underline{\hspace{3.2cm}} \quad \textbf{Data:} \underline{\hspace{1.6cm}} \quad \textbf{Nota:} \underline{\hspace{0.9cm}}/15


\noindent\textbf{Comando:} Criar o próprio trava-língua com 3 palavras que comecem com o mesmo som.


\begin{enumerate}

  \item Planeje: escolha 3 palavras da unidade para usar:
  \begin{center}
  1. \underline{\hspace{3cm}} \quad 2. \underline{\hspace{3cm}} \quad 3. \underline{\hspace{3cm}}
  \end{center}

  \item Escreva o texto (mínimo 3 frases):
  \begin{center}
  \underline{\hspace{12cm}}\\\underline{\hspace{12cm}}\\\underline{\hspace{12cm}}\\\underline{\hspace{12cm}}\\\underline{\hspace{12cm}}
  \end{center}

  \item Releia e corrija: marque palavras com a grafia TRAVA-LÍNGUA que você usou.

  \item Ilustre o texto:
  \begin{center}
  \begin{tikzpicture}
    \draw[thick, dashed] (0,0) rectangle (12,6);
  \end{tikzpicture}
  \end{center}

\end{enumerate}

\end{exercicio}

%% PLANCHETA 19.4

\plancheta{Folha de Prática 19.4 --- Ditado e Avaliação da Unidade 19}

\begin{exercicio}[Folha de Prática 19.4 --- Ditado e Avaliação da Unidade 19]

\noindent\textbf{Nome:} \underline{\hspace{3.2cm}} \quad \textbf{Data:} \underline{\hspace{1.6cm}} \quad \textbf{Nota:} \underline{\hspace{0.9cm}}/20


\subsection*{PARTE 1: DITADO (10 itens)}

O professor dita as palavras e frases abaixo; o aluno escreve com letra cursiva.

\begin{enumerate}

  \item \underline{\hspace{10.5cm}}

  \item \underline{\hspace{10.5cm}}

  \item \underline{\hspace{10.5cm}}

  \item \underline{\hspace{10.5cm}}

  \item \underline{\hspace{10.5cm}}

  \item \underline{\hspace{10.5cm}}

\end{enumerate}

\subsection*{PARTE 2: AUTOAVALIAÇÃO DO ALUNO}

\begin{enumerate}[label=\alph*)]

  \item Eu li as palavras da unidade sem ajuda.  \quad ( ) Sim   ( ) Não

  \item Eu escrevi o ditado com as grafias certas.  \quad ( ) Sim   ( ) Não

  \item Eu separei as palavras em sílabas.  \quad ( ) Sim   ( ) Não

  \item Eu sei explicar a regra com as minhas palavras.  \quad ( ) Sim   ( ) Não

\end{enumerate}

\subsection*{PARTE 3: REGISTRO DO PROFESSOR (S/F/R/N)}

\begin{center}
\renewcommand{\arraystretch}{1.4}
\begin{tabular}{@{}p{9cm}cccc@{}}
\toprule
Aspecto observado & S & F & R & N \\
\midrule
Lê com precisão as palavras-alvo & & & & \\
Escreve com ortografia estável & & & & \\
Generaliza a regra em escrita espontânea & & & & \\
Mantém atenção durante o ditado & & & & \\
\bottomrule\end{tabular}\end{center}

\end{exercicio}

\begin{dica}
**Dica para o Professor:** Grave a criança no trava-língua e reproduza: a autocorreção auditiva é mais eficaz que a correção do professor.
\end{dica}

%% ============================================
%% UNIDADE 20 -- Produção: o diário
%% ============================================

\chapter{Unidade 20 --- Produção: o diário}

\metadata{I-20}{Silábico-Alfabético Avançado}{EF04LP01, EF04LP02, EF04LP03, EF04LP05, EF04LP06, EF04LP07, EF04LP09}{D1}{EF04LP01, EF04LP02, EF04LP03, EF04LP05, EF04LP06, EF04LP07, EF04LP09}{Resolucao CNE/CEB 7/2010, art. 12}

\textbf{Regra:} A adivinha é uma charada em versos: a resposta está escondida nas pistas do texto. Ler as pistas é parte do jogo.

\textbf{Palavra geradora do bloco:} \textbf{OVO} --- o ovo está em adivinhas clássicas ('branco por fora, amarelo por dentro') e a palavra é curta e familiar.

\section{Lição 20.1 --- Produzir diários com registro pessoal}

\textbf{Foco da lição:} Produzir diários com registro pessoal


\begin{boquinha}[Boquinha da Técnica]
Na adivinha, leia as pistas devagar; a boca destaca as palavras-chave. Depois, responda e justifique: 'por que você acha isso?'
\end{boquinha}

\noindent\textbf{Formação guiada:}

\begin{center}
\resizebox{\textwidth}{!}{%
\begin{tabular}{cccc}
  \textbf{OVO} & \textbf{PISTA} & \textbf{CHAR} & \textbf{ADA} \\
  \end{tabular}}
\end{center}

\noindent\textbf{Palavras da unidade:} ovo, pista, charada, resposta, segredo, palma, escuro, claro.

\noindent\textbf{Frases para leitura oral:}
\begin{itemize}[leftmargin=1.6em]
  \item O ovo tem duas cores.
  \item A pista esconde a resposta.
  \item A charada fala do escuro.
  \item O segredo ficou claro no fim.
\end{itemize}

\section{Lição 20.2 --- Leitura do texto curto}

\noindent\textbf{Leia o texto em voz alta, com atenção às palavras da unidade:}


\begin{quote}
O ovo tem duas cores. A pista esconde a resposta.

A charada fala do escuro. O segredo ficou claro no fim.
\end{quote}

\textbf{Depois de ler, responda por escrito:}
\begin{enumerate}

  \item O que acontece no texto que você leu?
  \begin{center}
  \underline{\hspace{12cm}}
  \end{center}

  \item Quantas frases o texto tem?
  \begin{center}
  \underline{\hspace{12cm}}
  \end{center}

  \item Circule no texto as palavras com a grafia ADIVINHA.

  \item Leia o texto de novo em voz alta, agora sem ajuda.

\end{enumerate}

\section{Lição 20.3 --- Escrita com a grafia ADIVINHA}

\noindent\textbf{Regra da unidade:} A adivinha é uma charada em versos: a resposta está escondida nas pistas do texto. Ler as pistas é parte do jogo.


\noindent\textbf{Ditado guiado:} o professor dita as palavras abaixo; o aluno escreve e depois corrige com o professor.

\begin{center}
  \underline{\hspace{2.3cm}} \quad \underline{\hspace{2.3cm}} \quad \underline{\hspace{2.3cm}} \quad \underline{\hspace{2.3cm}}

  \underline{\hspace{2.3cm}} \quad \underline{\hspace{2.3cm}} \quad \underline{\hspace{2.3cm}} \quad \underline{\hspace{2.3cm}}
\end{center}

\noindent\textbf{Complete as palavras com a grafia ADIVINHA:}
\begin{center}
  ov\underline{~} \quad pist\underline{~} \quad charad\underline{~} \quad respost\underline{~} \quad segred\underline{~}
\end{center}

\noindent\textbf{Escreva a palavra que o professor mostrar (banco de palavras) e separe em sílabas:}

\begin{center}
  \underline{\hspace{2cm}} $=$ \underline{\hspace{2cm}} \quad \underline{\hspace{2cm}} $=$ \underline{\hspace{2cm}} \quad \underline{\hspace{2cm}} $=$ \underline{\hspace{2cm}}
  \end{center}

%% PLANCHETA 20.1

\plancheta{Folha de Prática 20.1 --- Produção: o diário (Parte 1)}

\begin{exercicio}[Folha de Prática 20.1 --- Produção: o diário (Parte 1)]

\noindent\textbf{Nome:} \underline{\hspace{3.2cm}} \quad \textbf{Data:} \underline{\hspace{1.6cm}} \quad \textbf{Nota:} \underline{\hspace{0.9cm}}/10


\subsection*{PARTE 1: RECONHECIMENTO}

\begin{enumerate}

  \item Forme as sílabas com os sons da unidade:

\begin{center}
\resizebox{\textwidth}{!}{%
\begin{tabular}{cccc}
  \textbf{OVO} & \textbf{PISTA} & \textbf{CHAR} & \textbf{ADA} \\
  \end{tabular}}
\end{center}

  \item Sublinhe as palavras da unidade:

\begin{center}
  \uline{ovo} \quad \uline{pista} \quad \uline{charada} \quad \uline{resposta} \quad \uline{segredo} \quad \uline{palma} \quad escuro \quad claro \quad gato \quad sapo \quad rato \quad bola \quad menino
\end{center}

  \item Identifique o som-alvo em cada palavra:
  \begin{center}
  OVO $\rightarrow$ som-alvo: \underline{\hspace{1cm}} \quad PISTA $\rightarrow$ som-alvo: \underline{\hspace{1cm}} \quad CHARADA $\rightarrow$ som-alvo: \underline{\hspace{1cm}}
  \end{center}

  \item Separe em sílabas:
  \begin{center}
  OVO $\rightarrow$ \underline{\hspace{1cm}} \quad PISTA $\rightarrow$ \underline{\hspace{1cm}} \quad CHARADA $\rightarrow$ \underline{\hspace{1cm}}
  \end{center}

  \item Conte as sílabas:
  \begin{center}
  OVO $\rightarrow$ \underline{\hspace{0.9cm}} \quad PISTA $\rightarrow$ \underline{\hspace{0.9cm}} \quad CHARADA $\rightarrow$ \underline{\hspace{0.9cm}}
  \end{center}

\end{enumerate}

\subsection*{PARTE 2: PRODUÇÃO GUIADA}

\begin{enumerate}[resume]

  \item Complete com a grafia da unidade:
  \begin{center}
  ADIVINHA\underline{~~~} \quad ADIVINHA\underline{~~~} \quad ADIVINHA\underline{~~~} \quad ADIVINHA\underline{~~~} \quad ADIVINHA\underline{~~~}
  \end{center}

  \item Escreva 3 palavras com o som-alvo:
  \begin{center}
  1. \underline{\hspace{3cm}} \quad 2. \underline{\hspace{3cm}} \quad 3. \underline{\hspace{3cm}}
  \end{center}

  \item Copie as palavras:
  \begin{center}
  ovo \quad pista \quad charada \quad resposta \quad segredo
  \end{center}

  \item Separe em sílabas:
  \begin{center}
  OVO = \underline{\hspace{1cm}} \underline{\hspace{1cm}} \quad PISTA = \underline{\hspace{1cm}} \underline{\hspace{1cm}} \quad CHARADA = \underline{\hspace{1cm}} \underline{\hspace{1cm}}
  \end{center}

  \item Escreva 2 frases usando palavras da unidade:
  \begin{enumerate}[label=\alph*)]
    \item \underline{\hspace{8cm}}
    \item \underline{\hspace{8cm}}
  \end{enumerate}

\end{enumerate}

\end{exercicio}

%% PLANCHETA 20.2

\plancheta{Folha de Prática 20.2 --- Produção: o diário (Parte 2)}

\begin{exercicio}[Folha de Prática 20.2 --- Produção: o diário (Parte 2)]

\noindent\textbf{Nome:} \underline{\hspace{3.2cm}} \quad \textbf{Data:} \underline{\hspace{1.6cm}} \quad \textbf{Nota:} \underline{\hspace{0.9cm}}/10


\subsection*{PARTE 3: INTEGRAÇÃO}

\begin{enumerate}[resume]

  \item Leia as palavras em voz alta:
  \begin{center}
  $\rightarrow$ ovo ~ pista ~ charada ~ resposta ~ segredo ~ palma
  \end{center}

  \item Leia as frases em voz alta:
  \begin{center}
  $\rightarrow$ O ovo tem duas cores.\\
  $\rightarrow$ A pista esconde a resposta.\\
  $\rightarrow$ A charada fala do escuro.\\
  $\rightarrow$ O segredo ficou claro no fim.\\
  \end{center}

  \item Circule a opção correta:
  \begin{enumerate}[label=(\alph*)]

    \item O som da unidade é: ( ) oral ( ) nasal 

    \item A primeira sílaba de OVO é: ( ) OVO ( ) PIS

    \item As palavras da unidade aparecem: ( ) no início ( ) no meio ( ) no fim das palavras

  \end{enumerate}

  \item Leia e complete:
  \begin{center}
  ovo\underline{~~~} \quad pis\underline{~~~} \quad cha\underline{~~~} \quad res\underline{~~~}
  \end{center}

  \item Leitura em voz alta com o professor:
  \begin{center}
  ovo \quad pista \quad charada \quad resposta \quad segredo \quad palma \quad escuro \quad claro
  \end{center}

\end{enumerate}

\end{exercicio}

%% PLANCHETA 20.3

\plancheta{Folha de Prática 20.3 --- Produção com ADIVINHA}

\begin{exercicio}[Folha de Prática 20.3 --- Produção com ADIVINHA]

\noindent\textbf{Nome:} \underline{\hspace{3.2cm}} \quad \textbf{Data:} \underline{\hspace{1.6cm}} \quad \textbf{Nota:} \underline{\hspace{0.9cm}}/15


\noindent\textbf{Comando:} Escrever uma adivinha para a turma, usando pelo menos duas pistas.


\begin{enumerate}

  \item Planeje: escolha 3 palavras da unidade para usar:
  \begin{center}
  1. \underline{\hspace{3cm}} \quad 2. \underline{\hspace{3cm}} \quad 3. \underline{\hspace{3cm}}
  \end{center}

  \item Escreva o texto (mínimo 3 frases):
  \begin{center}
  \underline{\hspace{12cm}}\\\underline{\hspace{12cm}}\\\underline{\hspace{12cm}}\\\underline{\hspace{12cm}}\\\underline{\hspace{12cm}}
  \end{center}

  \item Releia e corrija: marque palavras com a grafia ADIVINHA que você usou.

  \item Ilustre o texto:
  \begin{center}
  \begin{tikzpicture}
    \draw[thick, dashed] (0,0) rectangle (12,6);
  \end{tikzpicture}
  \end{center}

\end{enumerate}

\end{exercicio}

%% PLANCHETA 20.4

\plancheta{Folha de Prática 20.4 --- Ditado e Avaliação da Unidade 20}

\begin{exercicio}[Folha de Prática 20.4 --- Ditado e Avaliação da Unidade 20]

\noindent\textbf{Nome:} \underline{\hspace{3.2cm}} \quad \textbf{Data:} \underline{\hspace{1.6cm}} \quad \textbf{Nota:} \underline{\hspace{0.9cm}}/20


\subsection*{PARTE 1: DITADO (10 itens)}

O professor dita as palavras e frases abaixo; o aluno escreve com letra cursiva.

\begin{enumerate}

  \item \underline{\hspace{10.5cm}}

  \item \underline{\hspace{10.5cm}}

  \item \underline{\hspace{10.5cm}}

  \item \underline{\hspace{10.5cm}}

  \item \underline{\hspace{10.5cm}}

  \item \underline{\hspace{10.5cm}}

\end{enumerate}

\subsection*{PARTE 2: AUTOAVALIAÇÃO DO ALUNO}

\begin{enumerate}[label=\alph*)]

  \item Eu li as palavras da unidade sem ajuda.  \quad ( ) Sim   ( ) Não

  \item Eu escrevi o ditado com as grafias certas.  \quad ( ) Sim   ( ) Não

  \item Eu separei as palavras em sílabas.  \quad ( ) Sim   ( ) Não

  \item Eu sei explicar a regra com as minhas palavras.  \quad ( ) Sim   ( ) Não

\end{enumerate}

\subsection*{PARTE 3: REGISTRO DO PROFESSOR (S/F/R/N)}

\begin{center}
\renewcommand{\arraystretch}{1.4}
\begin{tabular}{@{}p{9cm}cccc@{}}
\toprule
Aspecto observado & S & F & R & N \\
\midrule
Lê com precisão as palavras-alvo & & & & \\
Escreve com ortografia estável & & & & \\
Generaliza a regra em escrita espontânea & & & & \\
Mantém atenção durante o ditado & & & & \\
\bottomrule\end{tabular}\end{center}

\end{exercicio}

\begin{dica}
**Dica para o Professor:** A adivinha exercita leitura inferencial: o aluno aprende que nem toda resposta está escrita; às vezes está escondida nas pistas.
\end{dica}

%% ============================================
%% UNIDADE 22 -- Sarau de leitura do 4º Ano
%% ============================================

\chapter{Unidade 22 --- Sarau de leitura do 4º Ano}

\metadata{I-22}{Silábico-Alfabético Avançado}{EF04LP01, EF04LP02, EF04LP03, EF04LP05, EF04LP06, EF04LP07, EF04LP09}{D1}{EF04LP01, EF04LP02, EF04LP03, EF04LP05, EF04LP06, EF04LP07, EF04LP09}{Resolucao CNE/CEB 7/2010, art. 12}

\textbf{Regra:} A revisão final organiza as grafias do ano em um mapa; o sarau celebra a fluência conquistada com leitura expressiva para a turma.

\textbf{Palavra geradora do bloco:} \textbf{SARAU} --- o sarau transforma a leitura em celebração: a meta do ano inteiro é ler com prazer em público.

\section{Lição 22.1 --- Celebrar a fluência e a entonação na leitura}

\textbf{Foco da lição:} Celebrar a fluência e a entonação na leitura


\begin{boquinha}[Boquinha da Técnica]
No sarau, a boca é o instrumento da festa: leitura clara, ritmo e emoção. Cada aluno apresenta um texto que escolheu.
\end{boquinha}

\noindent\textbf{Formação guiada:}

\begin{center}
\resizebox{\textwidth}{!}{%
\begin{tabular}{cccc}
  \textbf{SA} & \textbf{RAU} & \textbf{FESTA} & \textbf{VOZ} \\
  \end{tabular}}
\end{center}

\noindent\textbf{Palavras da unidade:} sarau, festa, poema, leitura, palco, plateia, aplauso, memória.

\noindent\textbf{Frases para leitura oral:}
\begin{itemize}[leftmargin=1.6em]
  \item O sarau tem poema e festa.
  \item A plateia escuta em silêncio.
  \item O palco recebe o leitor.
  \item O aplauso premia a leitura.
\end{itemize}

\section{Lição 22.2 --- Leitura do texto curto}

\noindent\textbf{Leia o texto em voz alta, com atenção às palavras da unidade:}


\begin{quote}
O sarau tem poema e festa. A plateia escuta em silêncio.

O palco recebe o leitor. O aplauso premia a leitura.
\end{quote}

\textbf{Depois de ler, responda por escrito:}
\begin{enumerate}

  \item O que acontece no texto que você leu?
  \begin{center}
  \underline{\hspace{12cm}}
  \end{center}

  \item Quantas frases o texto tem?
  \begin{center}
  \underline{\hspace{12cm}}
  \end{center}

  \item Circule no texto as palavras com a grafia SARAU.

  \item Leia o texto de novo em voz alta, agora sem ajuda.

\end{enumerate}

\section{Lição 22.3 --- Escrita com a grafia SARAU}

\noindent\textbf{Regra da unidade:} A revisão final organiza as grafias do ano em um mapa; o sarau celebra a fluência conquistada com leitura expressiva para a turma.


\noindent\textbf{Ditado guiado:} o professor dita as palavras abaixo; o aluno escreve e depois corrige com o professor.

\begin{center}
  \underline{\hspace{2.3cm}} \quad \underline{\hspace{2.3cm}} \quad \underline{\hspace{2.3cm}} \quad \underline{\hspace{2.3cm}}

  \underline{\hspace{2.3cm}} \quad \underline{\hspace{2.3cm}} \quad \underline{\hspace{2.3cm}} \quad \underline{\hspace{2.3cm}}
\end{center}

\noindent\textbf{Complete as palavras com a grafia SARAU:}
\begin{center}
  sara\underline{~} \quad fest\underline{~} \quad poem\underline{~} \quad leitur\underline{~} \quad palc\underline{~}
\end{center}

\noindent\textbf{Escreva a palavra que o professor mostrar (banco de palavras) e separe em sílabas:}

\begin{center}
  \underline{\hspace{2cm}} $=$ \underline{\hspace{2cm}} \quad \underline{\hspace{2cm}} $=$ \underline{\hspace{2cm}} \quad \underline{\hspace{2cm}} $=$ \underline{\hspace{2cm}}
  \end{center}

%% PLANCHETA 22.1

\plancheta{Folha de Prática 22.1 --- Sarau de leitura do 4º Ano (Parte 1)}

\begin{exercicio}[Folha de Prática 22.1 --- Sarau de leitura do 4º Ano (Parte 1)]

\noindent\textbf{Nome:} \underline{\hspace{3.2cm}} \quad \textbf{Data:} \underline{\hspace{1.6cm}} \quad \textbf{Nota:} \underline{\hspace{0.9cm}}/10


\subsection*{PARTE 1: RECONHECIMENTO}

\begin{enumerate}

  \item Forme as sílabas com os sons da unidade:

\begin{center}
\resizebox{\textwidth}{!}{%
\begin{tabular}{cccc}
  \textbf{SA} & \textbf{RAU} & \textbf{FESTA} & \textbf{VOZ} \\
  \end{tabular}}
\end{center}

  \item Sublinhe as palavras da unidade:

\begin{center}
  \uline{sarau} \quad \uline{festa} \quad \uline{poema} \quad \uline{leitura} \quad \uline{palco} \quad \uline{plateia} \quad aplauso \quad memória \quad gato \quad sapo \quad rato \quad bola \quad menino
\end{center}

  \item Identifique o som-alvo em cada palavra:
  \begin{center}
  SARAU $\rightarrow$ som-alvo: \underline{\hspace{1cm}} \quad FESTA $\rightarrow$ som-alvo: \underline{\hspace{1cm}} \quad POEMA $\rightarrow$ som-alvo: \underline{\hspace{1cm}}
  \end{center}

  \item Separe em sílabas:
  \begin{center}
  SARAU $\rightarrow$ \underline{\hspace{1cm}} \quad FESTA $\rightarrow$ \underline{\hspace{1cm}} \quad POEMA $\rightarrow$ \underline{\hspace{1cm}}
  \end{center}

  \item Conte as sílabas:
  \begin{center}
  SARAU $\rightarrow$ \underline{\hspace{0.9cm}} \quad FESTA $\rightarrow$ \underline{\hspace{0.9cm}} \quad POEMA $\rightarrow$ \underline{\hspace{0.9cm}}
  \end{center}

\end{enumerate}

\subsection*{PARTE 2: PRODUÇÃO GUIADA}

\begin{enumerate}[resume]

  \item Complete com a grafia da unidade:
  \begin{center}
  SARAU\underline{~~~} \quad SARAU\underline{~~~} \quad SARAU\underline{~~~} \quad SARAU\underline{~~~} \quad SARAU\underline{~~~}
  \end{center}

  \item Escreva 3 palavras com o som-alvo:
  \begin{center}
  1. \underline{\hspace{3cm}} \quad 2. \underline{\hspace{3cm}} \quad 3. \underline{\hspace{3cm}}
  \end{center}

  \item Copie as palavras:
  \begin{center}
  sarau \quad festa \quad poema \quad leitura \quad palco
  \end{center}

  \item Separe em sílabas:
  \begin{center}
  SARAU = \underline{\hspace{1cm}} \underline{\hspace{1cm}} \quad FESTA = \underline{\hspace{1cm}} \underline{\hspace{1cm}} \quad POEMA = \underline{\hspace{1cm}} \underline{\hspace{1cm}}
  \end{center}

  \item Escreva 2 frases usando palavras da unidade:
  \begin{enumerate}[label=\alph*)]
    \item \underline{\hspace{8cm}}
    \item \underline{\hspace{8cm}}
  \end{enumerate}

\end{enumerate}

\end{exercicio}

%% PLANCHETA 22.2

\plancheta{Folha de Prática 22.2 --- Sarau de leitura do 4º Ano (Parte 2)}

\begin{exercicio}[Folha de Prática 22.2 --- Sarau de leitura do 4º Ano (Parte 2)]

\noindent\textbf{Nome:} \underline{\hspace{3.2cm}} \quad \textbf{Data:} \underline{\hspace{1.6cm}} \quad \textbf{Nota:} \underline{\hspace{0.9cm}}/10


\subsection*{PARTE 3: INTEGRAÇÃO}

\begin{enumerate}[resume]

  \item Leia as palavras em voz alta:
  \begin{center}
  $\rightarrow$ sarau ~ festa ~ poema ~ leitura ~ palco ~ plateia
  \end{center}

  \item Leia as frases em voz alta:
  \begin{center}
  $\rightarrow$ O sarau tem poema e festa.\\
  $\rightarrow$ A plateia escuta em silêncio.\\
  $\rightarrow$ O palco recebe o leitor.\\
  $\rightarrow$ O aplauso premia a leitura.\\
  \end{center}

  \item Circule a opção correta:
  \begin{enumerate}[label=(\alph*)]

    \item O som da unidade é: ( ) oral ( ) nasal ( ) vibrante

    \item A primeira sílaba de SARAU é: ( ) SA ( ) RAU

    \item As palavras da unidade aparecem: ( ) no início ( ) no meio ( ) no fim das palavras

  \end{enumerate}

  \item Leia e complete:
  \begin{center}
  sar\underline{~~~} \quad fes\underline{~~~} \quad poe\underline{~~~} \quad lei\underline{~~~}
  \end{center}

  \item Leitura em voz alta com o professor:
  \begin{center}
  sarau \quad festa \quad poema \quad leitura \quad palco \quad plateia \quad aplauso \quad memória
  \end{center}

\end{enumerate}

\end{exercicio}

%% PLANCHETA 22.3

\plancheta{Folha de Prática 22.3 --- Produção com SARAU}

\begin{exercicio}[Folha de Prática 22.3 --- Produção com SARAU]

\noindent\textbf{Nome:} \underline{\hspace{3.2cm}} \quad \textbf{Data:} \underline{\hspace{1.6cm}} \quad \textbf{Nota:} \underline{\hspace{0.9cm}}/15


\noindent\textbf{Comando:} Apresentar no sarau um texto escolhido, lido com fluência e expressividade.


\begin{enumerate}

  \item Planeje: escolha 3 palavras da unidade para usar:
  \begin{center}
  1. \underline{\hspace{3cm}} \quad 2. \underline{\hspace{3cm}} \quad 3. \underline{\hspace{3cm}}
  \end{center}

  \item Escreva o texto (mínimo 3 frases):
  \begin{center}
  \underline{\hspace{12cm}}\\\underline{\hspace{12cm}}\\\underline{\hspace{12cm}}\\\underline{\hspace{12cm}}\\\underline{\hspace{12cm}}
  \end{center}

  \item Releia e corrija: marque palavras com a grafia SARAU que você usou.

  \item Ilustre o texto:
  \begin{center}
  \begin{tikzpicture}
    \draw[thick, dashed] (0,0) rectangle (12,6);
  \end{tikzpicture}
  \end{center}

\end{enumerate}

\end{exercicio}

%% PLANCHETA 22.4

\plancheta{Folha de Prática 22.4 --- Ditado e Avaliação da Unidade 22}

\begin{exercicio}[Folha de Prática 22.4 --- Ditado e Avaliação da Unidade 22]

\noindent\textbf{Nome:} \underline{\hspace{3.2cm}} \quad \textbf{Data:} \underline{\hspace{1.6cm}} \quad \textbf{Nota:} \underline{\hspace{0.9cm}}/20


\subsection*{PARTE 1: DITADO (10 itens)}

O professor dita as palavras e frases abaixo; o aluno escreve com letra cursiva.

\begin{enumerate}

  \item \underline{\hspace{10.5cm}}

  \item \underline{\hspace{10.5cm}}

  \item \underline{\hspace{10.5cm}}

  \item \underline{\hspace{10.5cm}}

  \item \underline{\hspace{10.5cm}}

  \item \underline{\hspace{10.5cm}}

\end{enumerate}

\subsection*{PARTE 2: AUTOAVALIAÇÃO DO ALUNO}

\begin{enumerate}[label=\alph*)]

  \item Eu li as palavras da unidade sem ajuda.  \quad ( ) Sim   ( ) Não

  \item Eu escrevi o ditado com as grafias certas.  \quad ( ) Sim   ( ) Não

  \item Eu separei as palavras em sílabas.  \quad ( ) Sim   ( ) Não

  \item Eu sei explicar a regra com as minhas palavras.  \quad ( ) Sim   ( ) Não

\end{enumerate}

\subsection*{PARTE 3: REGISTRO DO PROFESSOR (S/F/R/N)}

\begin{center}
\renewcommand{\arraystretch}{1.4}
\begin{tabular}{@{}p{9cm}cccc@{}}
\toprule
Aspecto observado & S & F & R & N \\
\midrule
Lê com precisão as palavras-alvo & & & & \\
Escreve com ortografia estável & & & & \\
Generaliza a regra em escrita espontânea & & & & \\
Mantém atenção durante o ditado & & & & \\
\bottomrule\end{tabular}\end{center}

\end{exercicio}

\begin{dica}
**Dica para o Professor:** O sarau é a culminância: celebre cada criança publicamente; a autoestima leitora construída hoje sustenta os anos seguintes.
\end{dica}

\input{checkpoints/rastreio-u6}
%% FIM DA PARTE II
